% ===================================================================
%  CV-Dyn2DGS: Chan-Vese 표면에 정렬된 2D Gaussian Surfel을 이용한
%              실시간 심장 표면 시각화
%
%  STATUS: 초고. 제 4-7장(방법, 오차 분석, 구현, 실험 설계)은 작성 완료.
%          제 8-9장(결과, 논의)은 측정값이 존재하지 않으므로 자리표시자만 있음.
%          자세한 내용은 paper/README.md 참조.
%
%  BUILD:  make            (xelatex + bibtex)
%          make pdflatex   (ko.TeX 폰트가 설치된 경우)
% ===================================================================
\documentclass[11pt,oneside]{report}

% ===================================================================
%  CV-Dyn2DGS thesis preamble
% ===================================================================
%  Korean body text via ko.TeX. Compiles with either
%    xelatex  (recommended; uses system Korean fonts)
%    pdflatex (needs a ko.TeX font set installed)
%  See paper/README.md for the exact commands.
% ===================================================================

\usepackage{kotex}
\usepackage[a4paper,margin=2.5cm]{geometry}

\usepackage{amsmath,amssymb,amsthm,mathtools}
\usepackage{graphicx}
\usepackage{booktabs}
\usepackage{array}
\usepackage{longtable}
\usepackage{multirow}
\usepackage{xcolor}
\usepackage{caption}
% natbib must be loaded before hyperref, and is REQUIRED by the
% \bibliographystyle{plainnat} used in main.tex: plainnat.bst emits \natexlab
% and the \bibitem[...]{...} form that only natbib defines. The body uses \citep.
\usepackage[numbers,square,sort&compress]{natbib}
\usepackage[hidelinks]{hyperref}
\usepackage{enumitem}
\usepackage{tikz}
\usetikzlibrary{arrows.meta,positioning,fit,backgrounds}

% -------------------------------------------------------------------
%  DRAFT MACHINERY
% -------------------------------------------------------------------
%  The empirical claims of this thesis (RQ1-RQ6) are all unmeasured at the time
%  of writing: the implementation exists but has not been executed. Rather than
%  leave blanks that could be mistaken for results, every unmeasured quantity is
%  marked, and a banner states the status on every page.
%
%  Clear \draftmodetrue ONLY after results/ has been produced by a real run of
%  `cvdyn2dgs experiments` and paper/tables/ regenerated by
%  scripts/make_tables.py.
% -------------------------------------------------------------------
\newif\ifdraftmode
\draftmodetrue

% "not measured" cell - deliberately conspicuous, never an empty cell
\newcommand{\NM}{%
  \ifdraftmode\textcolor{red}{\texttt{n/a}}\else\texttt{n/a}\fi}

% inline placeholder for a value that a real run must supply
\newcommand{\PH}[1]{\textcolor{red}{\texttt{<#1>}}}

% a whole paragraph that cannot be written before results exist
\newcommand{\pending}[1]{%
  \par\noindent\fcolorbox{red}{red!4}{%
    \begin{minipage}{0.97\linewidth}
    \small\textcolor{red}{\textbf{[결과 대기]}}\enspace #1
    \end{minipage}}\par\medskip}

% a note about a deliberate deviation or an honesty caveat
\newcommand{\caveat}[1]{%
  \par\noindent\fcolorbox{black!35}{black!3}{%
    \begin{minipage}{0.97\linewidth}
    \small\textbf{유의.}\enspace #1
    \end{minipage}}\par\medskip}

\ifdraftmode
  \usepackage{fancyhdr}
  \pagestyle{fancy}
  \fancyhf{}
  \fancyhead[C]{\small\textcolor{red}{초고 --- 실험 결과 미측정 (구현 미실행)}}
  \fancyfoot[C]{\thepage}
  \renewcommand{\headrulewidth}{0.4pt}
\fi

% -------------------------------------------------------------------
%  THEOREM ENVIRONMENTS
% -------------------------------------------------------------------
\theoremstyle{plain}
\newtheorem{theorem}{정리}[section]
\newtheorem{proposition}[theorem]{명제}
\newtheorem{lemma}[theorem]{보조정리}
\newtheorem{corollary}[theorem]{계}

\theoremstyle{definition}
\newtheorem{definition}[theorem]{정의}
\newtheorem{assumption}[theorem]{가정}
\newtheorem{example}[theorem]{예}

\theoremstyle{remark}
\newtheorem{remark}[theorem]{비고}

\renewcommand{\proofname}{증명}

% -------------------------------------------------------------------
%  NOTATION MACROS
% -------------------------------------------------------------------
% Level set and segmentation
\newcommand{\dom}{\Omega}                        % image domain
\newcommand{\lev}{\phi}                          % level-set function
\newcommand{\levt}{\lev_t}
\newcommand{\surf}{\Gamma}                       % surface
\newcommand{\surft}{\surf_t}
\newcommand{\img}{I}                             % image
\newcommand{\imgt}{\img_t}
\newcommand{\cin}{c^{\mathrm{in}}_t}
\newcommand{\cout}{c^{\mathrm{out}}_t}
\newcommand{\lin}{\lambda_{\mathrm{in}}}
\newcommand{\lout}{\lambda_{\mathrm{out}}}
\newcommand{\Heps}{H_\varepsilon}
\newcommand{\deps}{\delta_\varepsilon}
\newcommand{\ECV}{E_{\mathrm{CV}}}

% Spacing-aware operators (contribution 1)
\newcommand{\gradh}{\nabla_h}
\newcommand{\divh}{\operatorname{div}_h}
\newcommand{\spacing}{h}
\newcommand{\hmax}{h_{\max}}

% Surfels
\newcommand{\Gcan}{\mathcal{G}^{2D}_0}           % canonical surfel set
\newcommand{\anchor}{p}
\newcommand{\nrm}{n}
\newcommand{\framemat}{E}
\newcommand{\scalemat}{S}
\newcommand{\amp}{a}
\newcommand{\opa}{o}
\newcommand{\store}{\mathcal{D}}                 % stored representation

% Rendering
\newcommand{\camc}{o_c}
\newcommand{\raydir}{d_q}
\newcommand{\depthq}{\tau}
\newcommand{\alphaq}{\alpha}
\newcommand{\Chat}{\hat{C}}

% Errors / metrics
\newcommand{\Esurf}{E_{\mathrm{surf}}}
\newcommand{\Eflicker}{E_{\mathrm{flicker}}}
\newcommand{\Sfull}{S_{\mathrm{full}}}
\newcommand{\Sours}{S_{\mathrm{ours}}}
\newcommand{\CR}{\mathrm{CR}}

% Misc
\newcommand{\Real}{\mathbb{R}}
\newcommand{\norm}[1]{\left\lVert #1 \right\rVert}
\newcommand{\abs}[1]{\left\lvert #1 \right\rvert}
\newcommand{\inner}[2]{\left\langle #1,\, #2 \right\rangle}
\newcommand{\transp}{^{\top}}
\DeclareMathOperator*{\argmin}{arg\,min}
\DeclareMathOperator{\dist}{dist}
\DeclareMathOperator{\sgn}{sgn}
\DeclareMathOperator{\Length}{Length}

% Code cross-reference: ties an equation to the module that implements it
\newcommand{\impl}[1]{\textcolor{black!55}{\small\texttt{#1}}}
\newcommand{\implnote}[1]{\par\smallskip\noindent\impl{구현: #1}\par\smallskip}


\title{%
  \textbf{CV-Dyn2DGS}\\[0.4em]
  {\large Chan--Vese 표면에 정렬된 2D Gaussian Surfel을 이용한\\
   실시간 심장 표면 시각화}\\[0.8em]
  {\normalsize\itshape Chan--Vese-Guided Dynamic 2D Gaussian Surfels\\
   for Real-Time Cardiac Surface Visualization}}
\author{양수헌 (Sukhun Yang)\\[0.3em]
        {\small 서울대학교}\\[0.2em]
        {\small 지도교수: 강명주}}
\date{2026}

\begin{document}

\maketitle

\ifdraftmode
\begin{center}
\fcolorbox{red}{red!5}{%
\begin{minipage}{0.88\linewidth}
\small
\textcolor{red}{\textbf{초고 상태 고지}}\par\smallskip
본 원고의 제 4장(방법), 제 5장(오차 분석), 제 6장(구현), 제 7장(실험 설계)은
작성이 완료되었다. 그러나 제 8장(결과)의 모든 표와 제 9장(논의)은
\textbf{측정값이 존재하지 않는다}. 구현은 완성되었으나 실행된 적이 없으며
(제 \ref{sec:impl-status}절), 따라서 본 연구의 실험적 주장
RQ1--RQ6은 어느 것도 확인되지도 반박되지도 않은 상태이다.
미측정 항목은 모두 \NM 로 표시하였고, 결과에 의존하는 문단은
\textcolor{red}{[결과 대기]} 상자로 대체하였다.
어떤 수치도 추정하거나 인용 없이 기재하지 않았다.
\end{minipage}}
\end{center}
\vfill
\fi

\begin{abstract}
Cine 심장 자기공명영상(cine CMR)은 한 심장주기를 여러 시간 프레임으로 나누어
촬영한 3차원 영상의 시계열이다. 이를 분리된 3차원 표면으로 변환하면 의료진이
원 영상과 움직이는 심장 모델을 나란히 보고, 회전시키거나 시간에 따른 형태 변화를
관찰할 수 있다. 그러나 각 프레임마다 2D Gaussian Splatting(2DGS)을 처음부터
최적화하면 전처리시간과 저장량이 프레임 수에 비례해 증가하고, primitive identity가
프레임마다 달라져 재생 시 흔들림이 발생한다.

본 논문은 CV-Dyn2DGS를 제안한다. 첫 프레임의 Chan--Vese 표면 위에 법선방향 두께가
없는 2차원 Gaussian 원반(surfel)을 배치하여 정준(canonical) 2DGS를 구축하고, 이후
프레임에서는 이전 level-set을 seed로 현재 표면을 구한 뒤 기존 surfel의 중심을
새 표면의 법선 방향으로 정사영하고 접평면 프레임을 최소회전으로 전송한다. 표면
이동만으로 설명되지 않는 밝기 변화는 작은 프레임별 외형 잔차로 보정한다. 저장되는
것은 정준 2DGS 하나, 프레임별 표면, 그리고 작은 잔차뿐이며, dense deformation
field나 학습된 motion field는 저장하지 않는다.

이론적 기여는 네 가지이다. 첫째, 비등방 복셀 간격 $\spacing=(h_x,h_y,h_z)$을 명시적으로
반영한 간격 인지 이산화를 정립한다. 단기축 cine CMR은 $h_z \gg h_x \approx h_y$이므로
간격을 무시한 스텐실은 곡률 기반 정규화를 슬라이스 방향으로 편향시킨다. 둘째, 인접
프레임 해의 근접성을 이용한 좁은 띠 웜스타트가 수렴을 보존하며 반복 횟수를 줄임을
국소 Polyak--{\L}ojasiewicz 조건 아래 보인다. 셋째, 일차 Taylor 제약에 대한 최소
노름 논법으로 최소법선 보정을 유도하고, 그 반복이 잔차를 이차 속도로 감소시킴을
보인다. 넷째, 접평면 프레임, 원근 정합 래스터화, 외형 잔차 최소제곱, 중간 시점
보간의 네 축에 대해 명시적 오차 한계를 유도하고 각 한계를 검증 지표에 대응시킨다.

본 논문은 위 이론을 PyTorch로 구현한 시스템을 함께 제시한다. 구현은 간격 인지
연산자, 좁은 띠 Chan--Vese, 정준 surfel 초기화, 법선 정사영, 접선 전송, 원근 정합
타일 래스터화, 희소 렌더링 연산자를 이용한 정칙화 최소제곱, 저계수 잔차 압축,
그리고 mesh 및 thin-3DGS 비교 기준선을 포함한다.

\medskip
\noindent\textbf{실험적 주장의 현재 상태.}
본 논문의 효율 및 품질 주장 RQ1--RQ6은 \emph{가설이며 아직 확인되지 않았다}.
구현은 완성되었으나 실행된 적이 없다(제 \ref{sec:impl-status}절). 따라서 제 8장의
모든 표는 자리표시자이며, 어떤 성능 수치도 제시하지 않는다. 다만 PyTorch를
필요로 하지 않는 이산 논리 커널 41개 항목은 독립 참조와 대조하여 검증되었다
(제 \ref{sec:kernel-verification}절).

\end{abstract}

\tableofcontents

\chapter{서론}
\label{ch:intro}

\section{문제와 동기}

Cine 심장 자기공명영상은 한 심장주기를 $T$개의 시간 프레임으로 나누어 촬영한
3차원 영상의 시계열 $\imgt:\dom\subset\Real^3\to\Real$, $t=0,\dots,T-1$이다. 각
프레임에서 좌심실 등 관심 구조의 표면을 분할하고 이를 4차원(공간 3차원 + 시간)으로
시각화하면, 의료진은 원 영상과 움직이는 3차원 심장을 나란히 보며 자유롭게 회전하고
특정 시점을 정지할 수 있다.

이 목표를 달성하는 데는 두 개의 서로 다른 시간 척도가 관여한다. 하나는
\emph{전처리시간}, 즉 영상을 입력한 뒤 모든 프레임의 표면 표현을 만드는 데 걸리는
시간이며, 다른 하나는 \emph{재생시간}, 즉 전처리가 끝난 표현을 화면에 한 프레임
그리는 데 걸리는 시간이다. 본 논문에서 ``실시간''은 후자만을 뜻한다. 즉 환자별
전처리를 한 번 완료한 후의 상호작용적 재생이며, 촬영과 동시에 분할한다는 뜻이
아니다.

표면을 화면에 그리는 표현으로는 삼각형 mesh가 표준이지만, 표면 해상도와 shading에
따라 각진 인상이 생길 수 있다. 최근의 2D Gaussian Splatting(2DGS)~\citep{huang2024_2dgs}은
3차원 공간에 놓인 납작한 2차원 Gaussian 원반을 primitive로 사용하여, 부피를 가진
3차원 Gaussian~\citep{kerbl2023_3dgs}보다 view-consistent한 depth와 normal을 얻는다.
그러나 2DGS를 매 프레임 독립적으로 최적화하면 세 가지 문제가 생긴다. 전처리시간이
$T$에 비례하고, 저장량도 $T$에 비례하며, primitive identity가 프레임마다 달라져
재생 시 시간적 흔들림이 발생한다.

\section{핵심 착안}

본 논문의 착안은 다음 관찰에 있다. 인접한 cine 프레임 사이의 해부학적 형상 변화는
작다. 따라서 프레임마다 표현을 처음부터 만들 필요가 없고, 첫 프레임의 표현을
\emph{현재 프레임의 표면으로 밀어 옮기면} 된다. 이때 밀어 옮기는 근거를 학습된
변형장이 아니라 \emph{분할이 이미 제공하는 미분 정보}에서 가져온다.

구체적으로, Chan--Vese 활성 윤곽~\citep{chan2001_active}을 level-set
방법~\citep{osher1988_fronts}과 결합하면 각 프레임의 표면은 음함수 $\levt$의
영도집합 $\surft=\{\levt=0\}$으로 주어지고, 그 법선은 $\gradh\levt/\norm{\gradh\levt}$로
즉시 얻어진다. 이전 프레임의 surfel 중심을 이 법선 방향으로 정사영하면 새 표면
위로 보낼 수 있으며, 이는 3차원 변형장을 저장하거나 신경망을 학습하지 않고 국소
계산만으로 수행된다.

primitive로 2차원 원반을 택하는 이유는 차원 정합이다. Chan--Vese 표면 $\surft$는
$\Real^3$ 속의 2차원 곡면이므로, primitive도 2차원 원반이면 차원이 일치한다.
3차원 타원체는 법선방향 두께라는 자유도를 추가로 가지며, 이는 얇은 표면을 나타내는
데 불필요하고 렌더링 시 표면 흐림의 원인이 된다. 또한 원반은 3차원 공간의 평면
도형이므로 카메라 광선과의 교차를 직접 계산하여 원근 정합(perspective-correct)하게
투영할 수 있다.

\section{연구 질문}

본 논문의 핵심 연구 질문은 다음과 같다.

\begin{quote}
프레임별 2DGS를 독립적으로 최적화하는 대신, Chan--Vese 표면에 정렬된 정준 2D
Gaussian surfel을 법선 정사영하고 작은 잔차만 저장하면, mesh보다 매끄러운 표면
품질과 시간적 안정성을 확보하면서 전처리시간과 저장량을 줄이고 30 FPS 이상의
4D 심장 시각화를 달성할 수 있는가?
\end{quote}

이 질문은 제 \ref{ch:protocol}장에서 여섯 개의 검증 가능한 하위 질문
RQ1--RQ6으로 분해된다. 중요한 점은 이들이 모두 \emph{가설}이라는 것이다. 특히
제안 표현이 반드시 더 작다는 보장은 없다. 프레임별 표면 파일과 잔차가 크면
전 프레임 2DGS를 저장하는 것보다 오히려 클 수 있으며, 제
\ref{sec:storage-breakeven}절에서 그 손익분기 조건을 명시한다. 마찬가지로 단색
표면만 표시한다면 mesh가 더 단순하고 적절할 수 있으므로, 동일 표면의
mesh 렌더러를 필수 기준선으로 둔다. 2DGS가 실질적 이득을 주지 못한다는 결론도
유효한 연구 결과로 받아들인다.

\section{기여}

본 논문의 기여는 이론, 시스템, 검증의 세 층으로 나뉜다.

\paragraph{이론적 기여 (제 \ref{ch:method}, \ref{ch:error}장).}
\begin{enumerate}[leftmargin=2em]
  \item \textbf{간격 인지 이산화.} 비등방 복셀 간격을 명시적으로 반영한 경사 및
  발산 연산자를 정의하고, 이로부터 얻어지는 이산 경사 흐름 갱신식을 제시한다
  (제 \ref{sec:spacing-aware}절).
  \item \textbf{웜스타트 수렴.} 프레임 $t-1$의 해를 좁은 띠로 제한하여 프레임 $t$의
  초기값으로 쓰는 웜스타트가 수렴을 보존하고 반복 횟수를 줄임을, 국소
  Polyak--{\L}ojasiewicz 조건 아래 보인다 (정리 \ref{thm:warmstart}).
  \item \textbf{최소법선 정사영.} 일차 Taylor 제약에 대한 최소 노름 논법으로 보정
  식을 유도하고, 그 반복이 잔차를 이차 속도로 감소시킴을 보인다
  (정리 \ref{thm:projection}, 보조정리 \ref{lem:quadratic}).
  \item \textbf{네 축의 명시적 오차 분석.} 접평면 프레임, 원근 정합 래스터화,
  외형 잔차 최소제곱, 중간 시점 보간 각각에 대해 오차 한계를 유도하고 검증
  지표에 대응시킨다 (제 \ref{ch:error}장).
\end{enumerate}

\paragraph{시스템 기여 (제 \ref{ch:impl}장).}
위 이론을 PyTorch로 구현하였다. 이 과정에서 원 제안서에 없던 설계 결정들이
필요했으며, 이들은 그 자체로 보고할 가치가 있다. 특히 (i) 기하를 최적화 불가능한
버퍼로 두어 어떤 손실도 표면을 움직일 수 없게 구조적으로 보장한 점,
(ii) cine CMR에서 감독 신호가 무엇인지를 광선 행진으로 정의한 점,
(iii) 좁은 띠를 희소 목록이 아닌 조밀 crop으로 실현한 점,
(iv) 빠른 래스터라이저가 도입하는 두 근사를 브루트포스 참조 구현으로 측정 가능하게
한 점이다 (제 \ref{sec:impl-decisions}절).

\paragraph{검증 기여 (제 \ref{sec:kernel-verification}절).}
PyTorch를 필요로 하지 않는 이산 논리 커널을 표준 라이브러리만으로 재구현하여
독립 참조와 대조하였다. 41개 항목이 통과하였으며, 이는 본 논문에서 실제로
실행된 유일한 검증이다.

\section{본 논문이 주장하지 않는 것}
\label{sec:not-claimed}

명확히 해 둘 필요가 있는 범위 제한이 있다.

\begin{itemize}[leftmargin=2em]
  \item \textbf{최초성 주장을 하지 않는다.} 3D/4D Chan--Vese, 이전 윤곽을 이용한
  심장 분할, 정적 및 동적 2DGS, 정준 primitive와 변형을 이용한 4D-GS, Gaussian
  기반 심장 표현, 시간가변 심장 mesh는 모두 선행연구가 존재한다
  (제 \ref{ch:related}장). 본 논문이 다루는 것은 이들의 특정 조합과 그 효율--품질
  절충의 정량적 분석이다.
  \item \textbf{조직 추적이 아니다.} 법선 정사영은 가능한 변위 중 최소 노름을
  택하므로 접선 방향 성분을 결정하지 않는다. 결과는 기하학적 표면 운동이며,
  추가 검증 없이 심근 물질점 대응이나 strain으로 해석할 수 없다
  (제 \ref{sec:tangential}절).
  \item \textbf{volumetric renderer가 아니다.} 2DGS는 표면 렌더러이므로 심장 내부의
  임의 단면을 복원하지 않는다.
  \item \textbf{digital twin이 아니다.} 조직 물성, 전기전도, 혈류를 모델링하지
  않으며 치료 반응을 예측하지 않는다. 결과물은 시간가변 해부학 시각화 모델이다.
\end{itemize}

\section{논문 구성}

제 \ref{ch:related}장에서 선행연구를 정리하고 본 연구의 위치를 밝힌다.
제 \ref{ch:notation}장에서 기호와 문제를 정의한다. 제 \ref{ch:method}장이 방법의
본체로, 간격 인지 Chan--Vese부터 저장 표현까지를 전개한다. 제 \ref{ch:error}장에서
네 축의 오차 분석을 제시한다. 제 \ref{ch:impl}장에서 구현을, 제 \ref{ch:protocol}장에서
실험 설계를 기술한다. 제 \ref{ch:results}장은 결과 자리표시자이며,
제 \ref{ch:discussion}, \ref{ch:limitations}, \ref{ch:conclusion}장에서 논의,
한계, 결론을 다룬다.

\chapter{선행연구}
\label{ch:related}

\section{세 가지 문제의 구분}

비슷해 보이는 연구라도 실제로 푸는 문제는 다르다. \emph{분할}은 각 프레임에서
심장 영역이 어디인지 찾는 문제이고, \emph{운동 추적}은 이전 프레임의 동일한
조직점이 다음 프레임에서 어디로 갔는지 찾는 문제이며, \emph{렌더링}은 얻어진
구조와 밝기를 사람이 빠르게 관찰할 수 있는 영상으로 만드는 문제이다. 본 연구는
첫째와 셋째를 연결한다. 법선 정사영만으로는 접선 방향의 실제 조직 이동을 알 수
없으므로 둘째를 해결한다고 주장하지 않는다.

\section{활성 윤곽과 4D level-set}

Chan과 Vese는 영상 gradient 대신 경계 안팎의 밝기 평균 차이와 경계 길이를
최소화하는 활성 윤곽을 제안하였다~\citep{chan2001_active}. Level-set
방법~\citep{osher1988_fronts}은 곡면을 함수의 영도집합으로 나타내므로 2차원
곡선에서 3차원 곡면으로 자연스럽게 확장되고 위상 변화를 자동으로 다룬다.
Chan--Vese의 이론적 토대는 Mumford--Shah 범함수~\citep{mumford1989_optimal}의
조각마다 상수 축약이다.

Drapaca 등은 multi-phase Chan--Vese를 공간과 시간에 걸친 4D level-set으로 확장하여
연속 MR 영상의 변화하는 조직 경계를 분할하였다~\citep{drapaca2005_segmentation}.
따라서 ``시간축을 포함한 Chan--Vese'' 자체는 새로운 주장이 될 수 없다.

심장 영상에서도 이 방향은 충분히 연구되었다. Duan 등은 실시간 분할을 목표로 Active
Geometric Functions를 제안하였고~\citep{duan2010_realtime}, Queir\'os 등은 cine CMR의
ED에서 3차원 표면을 만든 뒤 해부학적 제약을 둔 optical flow로 전체 주기를
추적하였다~\citep{queiros2014_fast}. Avendi 등은 학습된 국소화 뒤에 변형 모형
(deformable model) 정련을 붙여 좌심실과 우심실을 분할하였는데, 이는 학습 요소와
변분적 표면 진화를 결합한 초기 사례이다~\citep{avendi2016_combined,avendi2017_rv}.
최근의 CSTM은 학습 기반 시공간 메모리로 전체 cine 시퀀스의 4D 분할을 수행하며
basal/apical slice의 through-plane motion 문제를 다룬다~\citep{ye2025_cstm}. 한편
현재 심장 분할의 사실상 기준선은 nnU-Net 계열이며~\citep{isensee2021_nnunet},
Chan--Vese의 Dice가 이들보다 낮다는 사실은 본 연구에 대한 반박이 아니다. 본 연구가
표면 공급원에서 요구하는 것은 mask가 아니라 부호 거리장 $\levt$와 그 기울기이며,
제 \ref{sec:external-baselines}절의 층위 1 실험은 공급원을 교체하여 이 구분을
측정 가능한 형태로 만든다.

\caveat{%
정정. 본 절의 Avendi 인용은 초기 원고에서 ``Chan--Vese 계열의 영역 항''을 사용한
연구로 잘못 기술되어 있었고, 출처 또한 존재하지 않는 서지사항이었다. 확인 결과 해당
연구는 CNN/auto-encoder 국소화와 변형 모형 정련의 결합이므로 위와 같이 고쳤다.}

\caveat{%
본 연구가 이 계열과 겹치는 부분은 (i) 이전 시간의 결과를 다음 시간에 재사용하고,
(ii) 심장 표면 시계열을 만들며, (iii) 시간적 일관성과 계산 속도를 평가한다는 점이다.
차이는 Chan--Vese 결과를 최종 mask로 끝내지 않고, $\levt$와 $\gradh\levt$를 surfel
중심과 접평면의 명시적 전송 규칙에 사용한다는 점이다. 따라서 본 연구의 독창성은
새로운 분할 이론이 아니라, 기존 변분 표면을 surface-native 렌더링 표현과 연결하는
수학적 갱신 규칙과 그 효율 분석에 있다.}

또한 본 연구가 ``4D Chan--Vese''라는 말을 쓰지 않는 이유를 밝혀 둔다. 그 용어에는
(i) $x,y,z,t$ 전체를 하나의 energy에서 동시에 최적화하는 결합 4D level-set과
(ii) 각 시간에서 3D Chan--Vese를 풀되 이전 결과를 시작점으로 쓰는 연속적 3D
Chan--Vese라는 두 뜻이 섞인다. 본 논문의 방법은 두 번째이므로
\emph{frame-to-frame warm-started 3D Chan--Vese}라고 정확히 표현한다.

\section{3DGS에서 2DGS로}

3D Gaussian Splatting은 장면을 부피를 가진 비등방 Gaussian 타원체로 나타내고 빠른
미분가능 래스터화로 실시간 novel-view 렌더링을 가능하게 하였다~\citep{kerbl2023_3dgs}.
그러나 얇은 표면을 3차원 blob으로 근사하면 두께와 시점에 따라 표면이 퍼지거나
geometry가 일관되지 않을 수 있다.

Huang 등의 2DGS는 각 primitive를 3차원 공간에 놓인 2차원 Gaussian 원반으로 바꾸고
카메라 광선과 원반의 교점을 원근 정합하게 계산한다~\citep{huang2024_2dgs}. Depth
distortion과 normal consistency 정규화를 사용하여 3DGS보다 표면 복원에 적합한
geometry를 얻는 것이 핵심이다. 본 연구에서 2DGS를 선택하는 이유는 점을 예쁘게
그리기 위해서가 아니라, Chan--Vese가 산출한 2차원 곡면과 rendering primitive의
차원을 맞추기 위해서이다.

정적 2DGS 자체는 본 연구의 독창성이 아니다. 본 연구는 알려진 2DGS 래스터화 원리를
사용하되, surfel의 중심과 접평면을 cine CMR의 Chan--Vese level-set으로부터 직접
갱신한다. 즉 새로 제안하는 대상은 렌더러가 아니라 시간 갱신과 저장 표현이다.

\section{동적 2DGS와 정준 표현 재사용}

Dynamic 2D Gaussians는 sparse control point로 2D Gaussian의 움직임을 제어하여 동적
객체의 렌더링과 mesh 추출을 수행한다~\citep{zhang2025_d2dgs}. Space-time 2DGS도
정준 2D Gaussian과 학습된 변형을 결합하고 depth 및 normal 정규화로 동적 표면을
복원한다~\citep{wang2024_st2dgs}. 따라서 ``2D Gaussian을 시간에 따라 이동한다''는
생각 역시 이미 존재한다. 4D-GS는 하나의 정준 3D Gaussian 집합에 신경 변형을 적용하여
동적 장면을 나타내므로, ``첫 표현을 만들고 시간에 따라 primitive를 재사용한다''는
큰 구조도 이미 존재한다~\citep{wu2024_4dgs}.

이들과 본 연구의 차이는 네 가지이다. 본 연구는 (i) cine CMR의 알려진 물리 좌표를
사용하고, (ii) 각 프레임의 Chan--Vese 영도집합을 목표 표면으로 삼으며, (iii) 신경
변형이나 control-point 망을 학습하지 않고, (iv) 최소 크기의 법선 이동으로 surfel을
갱신한다. 이러한 제약은 자유로운 동작 표현력을 낮추지만, 추가 학습자료 없이 환자별로
계산할 수 있고 각 이동의 근거가 $\levt$와 $\gradh\levt$로 명시된다.

\section{Dyna3DGR: 가장 가까운 선행연구}
\label{sec:dyna3dgr}

Dyna3DGR은 explicit 3D Gaussian과 implicit neural motion field를 함께 최적화하여
cine CMR에서 구조와 운동을 self-supervised하게 복원한다~\citep{fu2025_dyna3dgr}.
ACDC에서 학습용 point correspondence 없이 운동 추적을 평가했다는 점, 환자 하나에
대해 instance optimization을 수행한다는 점, Gaussian을 시간에 따라 재사용한다는
점에서 본 연구와 가장 많이 겹친다.

\begin{table}[htbp]
\centering
\small
\caption{Dyna3DGR과 CV-Dyn2DGS의 구조적 차이. 본 연구에서는 Dyna3DGR을 재현하지
않았으므로 수치 비교를 제시하지 않는다(제 \ref{sec:no-dyna-repro}절).}
\label{tab:dyna3dgr}
\begin{tabular}{@{}lll@{}}
\toprule
 & Dyna3DGR & CV-Dyn2DGS \\
\midrule
Primitive & 부피를 채우는 3D 타원체 & 표면 위 2D 원반 \\
시간 변화 & 연속 신경 motion field & 프레임별 $\surft$ + 국소 법선 정사영 \\
목표 & 위상/시간 일관 운동 추적 & 표면 외형의 저비용 저장과 실시간 재생 \\
표현 범위 & volumetric intensity와 내부 운동 & 명시적 장기 표면 \\
학습 & 환자별 self-supervised & 환자별, 신경망 없음 \\
취약점 & 최적화 비용, 모델 복잡성 & 접선 운동, 큰 프레임 간 이동, 위상 변화 \\
\bottomrule
\end{tabular}
\end{table}

\caveat{%
\label{sec:no-dyna-repro}%
Dyna3DGR은 본 연구에서 \textbf{재현하지 않았다}. ACDC 데이터와 GPU 학습이 필요하며
본 연구의 실행 환경에 GPU가 없었다. 따라서 어떤 수치도 Dyna3DGR에 귀속시키지 않으며,
표 \ref{tab:dyna3dgr}의 구조적 비교만 제시한다. 저자 답변에 보고된 환자별 최적화 약
11분이라는 값은 \emph{출처의 주장}으로만 언급하며 본 연구의 측정값과 나란히 놓지
않는다.

\textbf{단, 한 가지를 정정한다.} 원 제안서는 저자 공개 구현이 없어 재현이 어렵다고
가정하였으나, 확인 결과 Dyna3DGR의 구현은 \texttt{github.com/windrise/Dyna3DGR}에
공개되어 있다. 즉 재현 불가의 근거는 ``코드가 없다''가 아니라 ``본 원고 작성 환경에
GPU가 없다''는 것뿐이며, GPU가 있는 환경에서는 재현이 정당하게 요구될 수 있다.
제안서 §5.7의 면제 조항(``재현이 어려우면 구조적 차이만 보고'')을 이 논문이 계속
원용하는 것은 따라서 잠정적이다. 제 \ref{sec:external-baselines}절은 이를 미실시
비교로 명시적으로 기록한다.}

\section{시간가변 심장 mesh: surface-only 목적의 강한 대안}

Meng 등은 ED myocardial mesh의 vertex displacement를 SAX와 LAX cine MRI로부터
추정하는 딥러닝 방법을 제안하였다~\citep{meng2022_mesh}. 고정 mesh connectivity와
vertex correspondence를 유지하므로 기능적 측정에 유리하다. Liu 등의 개인화 4D
whole-heart mesh 연구는 multi-view cine MRI, mesh VAE, 다중 스케일 시간 모델링과
미분가능 contour rendering을 결합하여 네 개 심방/심실의 부드러운 mesh 궤적을
복원한다~\citep{liu2026_personalized}. biv-me는 이러한 아이디어를 실사용 가능한
공개 파이프라인 수준으로 확장한 매우 강한 선행연구로, cine CMR로부터 시간가변
biventricular mesh를 생성하고 volume, strain, wall thickness를 계산하며 여러
cohort에서 대규모 검증을 수행하였다~\citep{dillon2026_bivme}.

\caveat{%
즉 ``심장 표면을 자동으로 만들고 돌려 본다''는 목적만으로는 본 연구가 이들보다
새롭거나 유리하다고 말할 수 없다. CV-Dyn2DGS가 mesh 연구와 구별되려면 두 가지가
실험으로 확인되어야 한다. 첫째, 같은 Chan--Vese 표면에서 2D surfel이 mesh보다
silhouette, normal/depth 일관성 또는 MRI-derived appearance를 더 잘 전달해야 한다.
둘째, 프레임별 2DGS를 독립 생성하는 것보다 정준 2DGS와 작은 잔차의 저장 및 갱신
비용이 작아야 한다. 이 장점이 나타나지 않으면 ``단순 표면 viewer에는 mesh가 더
적절하다''가 올바른 결론이다. 또한 vertex correspondence가 중요한 strain 계산에서는
mesh 운동 추적이 본 연구보다 적합하다.}

\section{의료영상 Gaussian Splatting}

ClipGS는 정적 CT 및 해부학 volume을 대상으로 clipping plane을 지원하는 Gaussian
시각화를 제안하였다~\citep{li2025_clipgs}. 이는 의료용 GS viewer가 빠르게 동작할 수
있음을 보여주지만 장기의 시간 변형을 풀지는 않는다. GaussianPile은 slice thickness의
point-spread function을 고려한 비등방 Gaussian으로 slice 기반 의료 volume을
복원한다~\citep{kong2026_gaussianpile}. Park 등은 sparse non-aligned spine MRI에서
해부학적으로 유용한 새로운 단면을 Gaussian으로 렌더링한다~\citep{park2026_rendering}.

이들은 MRI/의료 volume의 내부 밝기를 GS로 표현한다는 점에서 더 강하지만, 반복되는
심장 표면을 이전 프레임에서 저비용으로 갱신하는 문제는 다루지 않는다. 역으로 본
연구의 2D 원반은 표면을 표현하므로 ``임의 단면을 정확히 복원하는 full volume
renderer''라고 주장할 수 없다.

\section{연구 공백의 정확한 진술}
\label{sec:gap}

위 검토에서 이미 다루어진 것은 3D/4D Chan--Vese, 이전 윤곽을 이용한 심장 분할,
정적 및 동적 2DGS, 정준 primitive와 변형을 이용한 4D-GS, 동적 심장 Gaussian,
시간가변 심장 mesh, 그리고 GS 기반 의료 volume 시각화이다. 따라서 본 논문은 최초의
4D Chan--Vese, 최초의 2DGS, 최초의 움직이는 2D Gaussian, 최초의 실시간 의료 GS
viewer 중 어느 것도 주장하지 않는다.

현재 조사 범위에서 직접 확인되지 않은 조합은 다음이다.

\begin{quote}
cine CMR의 프레임별 Chan--Vese implicit surface 값과 법선으로 정준 2D Gaussian
surfel의 중심과 접평면을 갱신하고, 별도의 신경/dense deformation field 없이
표면 시계열과 작은 외형 잔차만 저장하는 표현.
\end{quote}

이 문장은 체계적 문헌고찰에 의한 ``세계 최초'' 주장이 아니라, 본 연구 단계에서
확인된 구체적 연구 공백이다. 그러므로 기여의 강도는 조합 자체보다 실험 결과로
결정된다. 즉 (i) 동적 2DGS 및 Dyna3DGR보다 단순한 갱신이 어느 정도의 품질을
유지하는지, (ii) 독립 2DGS보다 시간과 저장량을 실제로 줄이는지, (iii) mesh 및 얇은
3DGS보다 2D 표면 primitive가 의미 있는지를 동시에 보여야 한다.

\chapter{기호와 문제 정의}
\label{ch:notation}

\section{입력}

한 환자의 cine CMR을
\begin{equation}
  \imgt : \dom \subset \Real^3 \to \Real, \qquad t = 0, 1, \dots, T-1
  \label{eq:input}
\end{equation}
로 나타낸다. $\imgt(x)$는 시간 $t$와 물리 좌표 $x$에서의 MRI 밝기이다. 복셀 간격은
\begin{equation}
  \spacing = (h_x, h_y, h_z)
  \label{eq:spacing}
\end{equation}
이며, 단기축 cine CMR에서는 일반적으로 $h_z \ge h_x \approx h_y$이다. 모든 거리,
곡률, 표면 오차는 복셀 지표가 아니라 \textbf{mm 단위}로 계산한다. 이 규약은 선택이
아니다. $h_z/h_x \approx 6$인 격자에서 복셀 지표로 계산한 표면 거리는 물리적 의미가
없다.

$\dom$은 유계 열린 영역이고 경계 $\partial\dom$은 조각마다 매끄럽다고 가정한다.
각 $\imgt$는 유계이며 제곱적분가능하다고 가정한다.

\section{표면의 level-set 표현}

각 프레임의 심장 표면은 level-set 함수 $\levt$의 영도집합으로 표현한다.
\begin{equation}
  \surft = \{\, x \in \dom : \levt(x) = 0 \,\}.
  \label{eq:levelset}
\end{equation}

부호 규약은 다음과 같다.
\begin{equation}
  \levt(x) > 0 : \text{심장 내부}, \qquad
  \levt(x) < 0 : \text{심장 외부}, \qquad
  \levt(x) = 0 : \text{심장 표면}.
  \label{eq:sign}
\end{equation}

\caveat{%
식 \eqref{eq:sign}은 신호 처리 문헌에서 흔한 ``내부가 음수'' 규약과 \textbf{반대}이다.
본 논문과 구현 전체에서 이 규약을 일관되게 사용하며, 구현의 모든 비교 연산은 이
규약에 대해 명시적으로 작성되었다. 부호를 뒤집으면 경사 흐름의 구동력 방향이
반대가 되어 표면이 발산한다.}

$\levt$가 $\surft$ 근방에서 매끄럽고 $\gradh\levt \ne 0$이면 단위 법선과 평균 곡률은
\begin{equation}
  \nrm = \frac{\gradh\levt}{\norm{\gradh\levt}}, \qquad
  \kappa = \divh\!\left(\frac{\gradh\levt}{\norm{\gradh\levt}}\right)
  \label{eq:normal-curvature}
\end{equation}
로 주어진다. 식 \eqref{eq:sign}의 규약에서 $\gradh\levt$는 \emph{내부를 향한다}.

\section{목표 구조의 범위}

기본 연구 범위에서는 좌심실 blood pool 또는 좌심실 내막 표면처럼 비교적 명확한
하나의 구조를 사용한다. 전체 좌/우심실 및 심근은 multiphase level-set을 요구하는
확장 범위로 둔다. 구현의 실제 데이터 로더는 좌심실 이외의 구조를 요청하면
조용히 다른 것을 반환하지 않고 예외를 발생시킨다.

\section{출력과 저장 표현}
\label{sec:storage-def}

첫 프레임에서 생성한 정준 2D Gaussian surfel 집합을
\begin{equation}
  \Gcan = \bigl\{\, g^0_i = (\anchor^0_i,\ e^0_{i,1},\ e^0_{i,2},\
  s_{i,1},\ s_{i,2},\ \amp^0_i,\ \opa^0_i) \,\bigr\}_{i=1}^{N}
  \label{eq:canonical-set}
\end{equation}
로 둔다. 각 성분의 뜻은 다음과 같다.

\begin{itemize}[leftmargin=2em]
  \item $\anchor^0_i \in \dom$: surfel의 중심(anchor).
  \item $e^0_{i,1}, e^0_{i,2} \in \Real^3$: 접평면을 이루는 서로 직교하는 두 단위
  접선축.
  \item $s_{i,1}, s_{i,2} > 0$: 두 접선축 방향의 scale. \textbf{법선방향 scale은 없다.}
  \item $\amp^0_i \in \Real$: MRI-derived 외형 진폭.
  \item $\opa^0_i \in [0,1]$: 불투명도.
\end{itemize}

법선방향 scale의 부재가 이 표현의 핵심 구조적 성질이다. primitive는 두께가 없는
2차원 원반이며 3차원 타원체가 아니다. 법선 방향 정보는 scale이 아니라 단위 법선
$\nrm^t_i$와 접평면 프레임 $\framemat^t_i$로만 표현된다.

최종 저장 표현은
\begin{equation}
  \boxed{\;
  \store = \Bigl\{\, \Gcan,\ \{\surft\}_{t=0}^{T-1},\ \{\Delta \amp_t\}_{t=0}^{T-1} \,\Bigr\}
  \;}
  \label{eq:store}
\end{equation}
이다. 즉 dense deformation field를 저장하지 않는다. 프레임별 위치와 회전 전체를
저장하는 대신, 재생 시 현재 표면으로 anchor를 법선 정사영하고 표면 법선으로 접평면을
계산하여 surfel의 위치와 방향을 갱신한다.

\caveat{%
식 \eqref{eq:store}에 프레임별 anchor가 없다는 사실은 임의 시점 탐색(seek)에 대한
비자명한 함의를 가진다. 전처리는 anchor를 $\anchor^{t-1}\to\anchor^t$로 \emph{순차}
정사영하지만, 저장된 것만으로 프레임 $t$로 바로 이동하려면 \emph{정준} anchor를
$\surft$로 직접 정사영해야 한다. 두 경로는 동일하지 않다. 제
\ref{sec:seek-problem}절에서 이 문제를 명시하고 구현에서 두 방식을 모두 제공하여
그 차이를 측정 가능하게 하였다.}

\section{Real-time의 정의}

본 논문에서 real-time은 MRI scanner가 신호를 획득하는 동시에 처리한다는 뜻이
아니다. 환자별 cine CMR에 대해 한 번의 전처리를 완료한 후, 저장된 표현을 사용하여
다음 상호작용을 30 FPS 이상으로 수행하는 것을 뜻한다.

\begin{itemize}[leftmargin=2em]
  \item 심장주기 재생 및 정지,
  \item 자유로운 카메라 회전과 확대,
  \item MRI 영상과 3차원 심장 표현의 동시 표시,
  \item 선택적 clipping과 표면 depth/normal map 관찰.
\end{itemize}

따라서 전처리시간과 재생시간은 별도의 지표로 측정한다. 전처리가 느려도 viewer는
실시간일 수 있으며, 반대로 전처리가 빨라도 렌더링이 느릴 수 있다.

\section{매끄러운 Heaviside와 Dirac 근사}

영역의 지시 함수와 경계 측도를 변분 계산에 쓸 수 있게 매끄럽게 근사한다. 정규화
폭 $\varepsilon > 0$에 대하여
\begin{align}
  \Heps(z) &= \frac{1}{2}\left(1 + \frac{2}{\pi}\arctan\frac{z}{\varepsilon}\right),
  \label{eq:heaviside}\\
  \deps(z) &= \Heps'(z)
  = \frac{1}{\pi}\cdot\frac{\varepsilon}{\varepsilon^2 + z^2}
  \label{eq:dirac}
\end{align}
로 정의한다.

\begin{proposition}[근사의 기본 성질]
\label{prop:heaviside}
식 \eqref{eq:heaviside}, \eqref{eq:dirac}에 대하여 다음이 성립한다.
\begin{enumerate}[label=(\roman*)]
  \item 모든 $\varepsilon>0$에서 $\Heps \in C^\infty(\Real)$이고
  $0 < \Heps < 1$이며 $\Heps$는 순증가한다.
  \item $\int_\Real \deps(z)\,dz = 1$.
  \item $\varepsilon \to 0^+$에서 $z \ne 0$인 각 점에서 $\Heps(z) \to H(z)$이고,
  분포의 의미에서 $\deps \to \delta$이다.
\end{enumerate}
\end{proposition}

\begin{proof}
(i) $\arctan$이 매끄럽고 유계이며 순증가하므로 $\Heps$도 그러하다. 치역은
$\arctan$의 치역 $(-\pi/2,\pi/2)$로부터 $(0,1)$이다.

(ii) 치환 $w = z/\varepsilon$을 쓰면
\[
  \int_\Real \frac{1}{\pi}\frac{\varepsilon}{\varepsilon^2+z^2}\,dz
  = \frac{1}{\pi}\int_\Real \frac{dw}{1+w^2}
  = \frac{1}{\pi}\bigl[\arctan w\bigr]^{\infty}_{-\infty} = 1.
\]

(iii) $z>0$이면 $\arctan(z/\varepsilon)\to\pi/2$이므로 $\Heps(z)\to1$이고, $z<0$이면
$\to 0$이다. (ii)에 의해 $\deps$는 총질량 1을 유지하면서 원점 근방으로 질량이
집중되므로, 임의의 유계 연속 시험함수 $\psi$에 대해 $\int \deps \psi \to \psi(0)$이다.
\end{proof}

(i)의 \emph{순 양성} $\Heps > 0$이 중요하다. 이것이 식 \eqref{eq:cin}, \eqref{eq:cout}의
분모가 영이 되지 않음을 보장한다.

\implnote{%
식 \eqref{eq:heaviside}, \eqref{eq:dirac}는 \texttt{levelset/operators.py}의
\texttt{heaviside\_eps}, \texttt{dirac\_eps}. 명제 \ref{prop:heaviside}의 (ii)는
\texttt{experiments/theory\_checks.py::check\_smooth\_heaviside}에서 수치 확인.}

\section{인위적 시간과 물리적 시간}

본 논문에는 두 개의 시간 변수가 있으며 혼동하지 않도록 유의한다.

\begin{itemize}[leftmargin=2em]
  \item $t$: cine 영상의 \textbf{프레임 지표}, 즉 실제 심장 시간.
  \item $s \ge 0$: \textbf{인위적 최적화 시간}. 각 고정된 프레임 $t$에 대하여 energy를
  최소화하는 경사 하강의 진행을 매개하는 가상의 변수이며 물리적 심장 시간이 아니다.
\end{itemize}

웜스타트는 프레임 $t$에서 $s\to\infty$로 얻은 정상 상태를 프레임 $t+1$의 $s=0$
초기값으로 넘겨주는 절차이다. surfel 기하의 위 첨자 $t$는 언제나 물리적 프레임을
가리키며 $s$와 무관하다.

\chapter{방법}
\label{ch:method}

\section{전체 파이프라인}
\label{sec:pipeline}

한 환자에 대한 처리는 전처리와 재생으로 나뉜다. 전처리는 다음 순서로 진행된다.

\begin{enumerate}[leftmargin=2.2em]
  \item 첫 프레임에서 제공된 ED label 또는 내부 seed로 3D Chan--Vese를 풀어
  $\lev_0$을 얻는다 (제 \ref{sec:spacing-aware}절).
  \item $\surf_0$ 위에 2D Gaussian 원반을 배치하고 MRI 밝기, 불투명도, scale을
  맞추어 정준 2DGS $\Gcan$을 만든다 (제 \ref{sec:canonical}절).
  \item 프레임 $t \ge 1$에서 $\lev_{t-1}$을 seed로 Chan--Vese를 짧게 수행하여
  $\levt$를 구한다 (제 \ref{sec:warmstart}절).
  \item 각 surfel 중심을 $\surft$의 법선 방향으로 정사영하고 (제
  \ref{sec:projection}절), 두 접선축을 새 접평면으로 전송한다 (제
  \ref{sec:transport}절).
  \item 표면 이동만으로 설명되지 않는 밝기 차이를 작은 잔차로 fitting한다
  (제 \ref{sec:residual}절).
  \item 마지막 프레임까지 반복하고 $\store$만 저장한다 (제
  \ref{sec:storage-breakeven}절).
\end{enumerate}

재생에서는 Chan--Vese나 비선형 최적화를 다시 수행하지 않고, 저장된 표면을 이용한
중심/접평면 갱신과 원근 정합 래스터화만 수행한다 (제 \ref{sec:raster}절).

\paragraph{계산 비용의 구조.}
각 프레임에서 전체 2DGS를 처음부터 최적화하는 비용을 $C_{\mathrm{full}}$이라 하면
독립 최적화의 총비용은
\begin{equation}
  C_{\mathrm{independent}} = T\,C_{\mathrm{full}}
  \label{eq:cost-independent}
\end{equation}
이다. 제안 방법은 첫 프레임에서만 전체 fitting을 수행하므로
\begin{equation}
  C_{\mathrm{ours}} = C_{\mathrm{full}}
  + \sum_{t=1}^{T-1}\Bigl(C^t_{\mathrm{CV}} + C^t_{\mathrm{project}}
  + C^t_{\mathrm{orient}} + C^t_{\mathrm{res}}\Bigr)
  \label{eq:cost-ours}
\end{equation}
이다. 본 방법이 실제로 빠르려면
\begin{equation}
  C^t_{\mathrm{CV}} + C^t_{\mathrm{project}} + C^t_{\mathrm{orient}}
  + C^t_{\mathrm{res}} \;<\; C_{\mathrm{full}}
  \label{eq:cost-hypothesis}
\end{equation}
이어야 한다.

\caveat{%
식 \eqref{eq:cost-hypothesis}은 \textbf{검증할 가설이며 미리 참이라고 가정하는
결론이 아니다}. 네 항 각각을 별도로 측정하여 항별로 확인한다. 구현은 이 네 항을
독립된 타이머 스테이지로 계측한다.}

\section{간격 인지 3D Chan--Vese}
\label{sec:spacing-aware}

\subsection{에너지 범함수}

영상 분할의 변분적 토대는 Mumford--Shah 범함수~\citep{mumford1989_optimal}이다.
근사 함수 $u$가 경계 $C=\surft$로 나뉜 두 영역에서 각각 상수라고 제한하는
조각마다 상수 축약을 적용하면 매끄러움 항이 소멸하고 자료 충실도 항과 경계 측도
항만 남는다. 지시 함수를 $\Heps(\levt)$와 $1-\Heps(\levt)$로, 경계 측도를
$\deps(\levt)\norm{\gradh\levt}$의 적분으로 대체하면 다음을 얻는다.

\begin{definition}[Chan--Vese 에너지, 간격 인지 형태]
\label{def:cv-energy}
가중치 $\mu, \lin, \lout > 0$와 프레임 $t$의 level-set 함수 $\levt$, 영역 평균
$\cin, \cout$에 대하여
\begin{align}
  \ECV(\levt, \cin, \cout)
  &= \mu \int_\dom \deps(\levt)\,\norm{\gradh\levt}_2 \,dx \notag\\
  &\quad + \lin \int_\dom \abs{\imgt - \cin}^2\,\Heps(\levt)\,dx \notag\\
  &\quad + \lout \int_\dom \abs{\imgt - \cout}^2\,\bigl(1 - \Heps(\levt)\bigr)\,dx.
  \label{eq:cv-energy}
\end{align}
\end{definition}

첫째 항은 길이/면적 정규화이다. $\deps(\levt)\norm{\gradh\levt}$는 영도집합
$\surft$ 위의 근사 면적 측도이므로, 이 항은 표면 면적을 벌점으로 부과하여 매끄러운
경계를 유도한다. 둘째와 셋째 항은 각각 내부와 외부의 자료 충실도이다.

분할 문제는 $\min_{\levt,\cin,\cout}\ECV$로 정식화되며, 세 변수에 대한 결합 최소화를
교대 최소화로 접근한다.

\subsection{최적 영역 평균}

\begin{proposition}[최적 영역 평균]
\label{prop:region-means}
$\levt$를 고정하면 $\ECV$를 최소화하는 영역 평균은
\begin{align}
  \cin  &= \frac{\displaystyle\int_\dom \imgt\,\Heps(\levt)\,dx}
                {\displaystyle\int_\dom \Heps(\levt)\,dx},
  \label{eq:cin}\\[0.4em]
  \cout &= \frac{\displaystyle\int_\dom \imgt\,\bigl(1-\Heps(\levt)\bigr)\,dx}
                {\displaystyle\int_\dom \bigl(1-\Heps(\levt)\bigr)\,dx}
  \label{eq:cout}
\end{align}
이다.
\end{proposition}

\begin{proof}
$\cin$은 식 \eqref{eq:cv-energy}의 둘째 항에만 나타나므로
\[
  \frac{\partial \ECV}{\partial \cin}
  = \lin \int_\dom -2\,(\imgt - \cin)\,\Heps(\levt)\,dx = 0.
\]
$\lin>0$이므로 $\int \imgt \Heps(\levt)\,dx = \cin \int \Heps(\levt)\,dx$이고,
명제 \ref{prop:heaviside}(i)에 의해 $\Heps>0$이므로 $\int \Heps(\levt)\,dx>0$이다.
따라서 나누어 식 \eqref{eq:cin}을 얻는다. 둘째 미분은
$\partial^2 \ECV/\partial(\cin)^2 = 2\lin\int \Heps(\levt)\,dx > 0$이므로 최소점이다.
$\cout$도 같은 논법으로 얻는다.
\end{proof}

식 \eqref{eq:cin}, \eqref{eq:cout}은 $\levt$가 고정된 상태의 닫힌 형태 해이므로 교대
최소화의 각 단계에서 즉시 계산된다. $\Heps$가 급격한 이진 지시 함수가 아니므로
경계 근방의 복셀도 부분적으로 두 평균에 기여한다.

\subsection{Euler--Lagrange 방정식과 경사 흐름}

\begin{theorem}[Chan--Vese Euler--Lagrange 방정식과 경사 흐름]
\label{thm:gradient-flow}
$\cin, \cout$을 고정하고 $\levt \in C^2(\dom)$에서 $\ECV$의 최소화를 고려하자.
$\ECV$의 정류점은 다음 Euler--Lagrange 조건을 만족한다.
\begin{equation}
  \deps(\levt)\left[
    \mu\,\divh\!\left(\frac{\gradh\levt}{\norm{\gradh\levt}}\right)
    - \lin (\imgt - \cin)^2 + \lout (\imgt - \cout)^2
  \right] = 0
  \quad \text{in } \dom,
  \label{eq:euler-lagrange}
\end{equation}
자연/Neumann 경계 조건
\begin{equation}
  \frac{\deps(\levt)}{\norm{\gradh\levt}}\,\frac{\partial \levt}{\partial \nrm} = 0
  \quad \text{on } \partial\dom
  \label{eq:neumann}
\end{equation}
와 함께 성립한다. 이에 대응하며, 경사가 작은 곳의 수치 안정성을 위해 분모를
정규화한 경사 하강 흐름은 인위적 시간 $s$에 대하여
\begin{equation}
  \frac{\partial \levt}{\partial s}
  = \deps(\levt)\left[
    \mu\,\divh\!\left(\frac{\gradh\levt}{\norm{\gradh\levt}+\varepsilon}\right)
    - \lin (\imgt - \cin)^2 + \lout (\imgt - \cout)^2
  \right]
  \label{eq:gradient-flow}
\end{equation}
이다.
\end{theorem}

\begin{proof}
$\levt$의 임의의 매끄러운 섭동 $\psi \in C^\infty_c(\dom)$ 방향으로의 일차 변분을
계산한다. $E(\tau) := \ECV(\levt + \tau\psi, \cin, \cout)$로 두고
$\frac{d}{d\tau}E(\tau)\big|_{\tau=0}$을 세 항별로 계산한다. 이하 $\lev := \levt$로
줄여 쓰고, 연속 수준의 연산자를 $\nabla,\operatorname{div}$로 표기한다.

\textbf{(1) 길이 항.} $L(\lev) = \mu\int_\dom \deps(\lev)\norm{\nabla\lev}\,dx$에 대하여
$\frac{d}{d\tau}\deps(\lev+\tau\psi)\big|_0 = \deps'(\lev)\psi$이고
$\frac{d}{d\tau}\norm{\nabla(\lev+\tau\psi)}\big|_0 = \frac{\nabla\lev\cdot\nabla\psi}{\norm{\nabla\lev}}$
이므로
\begin{equation}
  \frac{dL}{d\tau}\bigg|_0
  = \mu \int_\dom \left[
    \deps'(\lev)\,\psi\,\norm{\nabla\lev}
    + \deps(\lev)\frac{\nabla\lev\cdot\nabla\psi}{\norm{\nabla\lev}}
  \right] dx.
  \label{eq:proof-length}
\end{equation}
둘째 피적분항에 벡터장 $V = \deps(\lev)\frac{\nabla\lev}{\norm{\nabla\lev}}$에 대한
발산 정리를 적용한다. $\operatorname{div}(\psi V) = \nabla\psi\cdot V + \psi\operatorname{div}V$이므로
\[
  \int_\dom \nabla\psi \cdot V\,dx
  = \int_{\partial\dom} \psi\,V\cdot \nrm\,dS - \int_\dom \psi\,\operatorname{div}V\,dx.
\]
곱셈 규칙과 $\nabla\lev\cdot\nabla\lev = \norm{\nabla\lev}^2$을 사용하면
\[
  \operatorname{div}V
  = \deps'(\lev)\,\norm{\nabla\lev}
  + \deps(\lev)\operatorname{div}\!\left(\frac{\nabla\lev}{\norm{\nabla\lev}}\right).
\]
이를 식 \eqref{eq:proof-length}에 대입하면 첫째 항의
$\deps'(\lev)\psi\norm{\nabla\lev}$와 $\operatorname{div}V$의 동일한 성분이 \emph{정확히
상쇄}되어
\begin{equation}
  \frac{dL}{d\tau}\bigg|_0
  = -\mu\int_\dom \psi\,\deps(\lev)\operatorname{div}\!\left(\frac{\nabla\lev}{\norm{\nabla\lev}}\right) dx
  + \mu\int_{\partial\dom} \psi\,\frac{\deps(\lev)}{\norm{\nabla\lev}}
    \frac{\partial\lev}{\partial\nrm}\,dS
  \label{eq:proof-length-final}
\end{equation}
를 얻는다.

\textbf{(2) 내부 충실도 항.}
$F_{\mathrm{in}}(\lev) = \lin\int_\dom \abs{\imgt-\cin}^2 \Heps(\lev)\,dx$에서 $\cin$은
고정이므로 $\Heps$만 미분되고, $\Heps' = \deps$이므로
\begin{equation}
  \frac{dF_{\mathrm{in}}}{d\tau}\bigg|_0
  = \lin \int_\dom \abs{\imgt-\cin}^2 \deps(\lev)\,\psi\,dx.
\end{equation}

\textbf{(3) 외부 충실도 항.}
$\frac{d}{d\tau}\bigl(1-\Heps(\lev+\tau\psi)\bigr)\big|_0 = -\deps(\lev)\psi$이므로
\begin{equation}
  \frac{dF_{\mathrm{out}}}{d\tau}\bigg|_0
  = -\lout \int_\dom \abs{\imgt-\cout}^2 \deps(\lev)\,\psi\,dx.
\end{equation}

\textbf{(4) 결합.} 세 항을 더하면 전체 변분은
\[
  \frac{dE}{d\tau}\bigg|_0
  = \int_\dom \psi\,\deps(\lev)\left[
    -\mu\operatorname{div}\!\left(\frac{\nabla\lev}{\norm{\nabla\lev}}\right)
    + \lin(\imgt-\cin)^2 - \lout(\imgt-\cout)^2
  \right] dx + (\text{경계항})
\]
이다. $\lev$가 정류점이면 모든 시험함수 $\psi$에 대해 이 값이 영이어야 하므로,
변분법의 기본 보조정리에 의해 대괄호에 $\deps(\lev)$를 곱한 것이 $\dom$ 내부에서
영이어야 한다. 부호를 정리하면 식 \eqref{eq:euler-lagrange}이다. 경계항
\eqref{eq:proof-length-final}이 모든 $\psi$에 대해 소멸하려면 식 \eqref{eq:neumann}이
성립해야 한다.

\textbf{(5) 경사 흐름.} $L^2(\dom)$ 내적에 대한 $\ECV$의 경사는 위 변분의 피적분
계수이다. 경사 하강은 $\partial_s\lev = -\delta \ECV/\delta\lev$로 정의되므로 부호를
정리하여 식 \eqref{eq:gradient-flow}을 얻는다. 연속 연산자를 간격 인지 연산자로
대체하고, 경사가 작은 곳에서 정규화된 곡률 항의 분모가 영이 되어 폭주하는 것을 막기
위해 $\norm{\gradh\lev}$를 $\norm{\gradh\lev}+\varepsilon$으로 정규화하였다.
\end{proof}

\begin{remark}[부호의 해석]
식 \eqref{eq:sign}의 규약에서 어떤 복셀이 내부($\lev>0$)에 있으면서 그 밝기가
$\cin$보다 $\cout$에 더 가까우면 대괄호가 음수가 되어 $\lev$가 감소하고, 그 복셀은
내부에서 빠져나간다. 이것이 의도한 방향이다. 부호를 하나만 뒤집으면 표면이
발산한다.
\end{remark}

\begin{remark}[$\divh,\gradh$로의 대체의 지위]
식 \eqref{eq:gradient-flow}에서 연속 연산자를 간격 인지 연산자로 대체한 것은 비등방
계량에서 성립하는 정확한 연속 항등식이 아니라, 연속 $L^2$ 경사 흐름의 \emph{간격
인지 이산화} 대응물로 이해해야 한다. 이 이산화를 다음 절에서 구체적으로 구성한다.
\end{remark}

\subsection{간격 인지 이산화: 첫 번째 기여}
\label{sec:discretization}

자기공명영상은 복셀 간격이 축마다 다른 비등방 격자에서 표본화되며, 슬라이스 간격
$h_z$가 평면 내 화소 간격 $h_x \approx h_y$보다 크다. 표준 유한차분이 이 간격을
무시하면 경사와 발산이 축 방향으로 왜곡되어 곡률 기반 정규화 항이 편향된다.

격자점 $(i,j,k)$에서의 이산 level-set 값을 $\lev_{i,j,k}$라 하고, 전진 차분을
\begin{equation}
  D^+_x \lev_{i,j,k} = \frac{\lev_{i+1,j,k}-\lev_{i,j,k}}{h_x},\quad
  D^+_y \lev_{i,j,k} = \frac{\lev_{i,j+1,k}-\lev_{i,j,k}}{h_y},\quad
  D^+_z \lev_{i,j,k} = \frac{\lev_{i,j,k+1}-\lev_{i,j,k}}{h_z},
  \label{eq:forward-diff}
\end{equation}
후진 차분을
\begin{equation}
  D^-_x \lev_{i,j,k} = \frac{\lev_{i,j,k}-\lev_{i-1,j,k}}{h_x},\quad
  D^-_y \lev_{i,j,k} = \frac{\lev_{i,j,k}-\lev_{i,j-1,k}}{h_y},\quad
  D^-_z \lev_{i,j,k} = \frac{\lev_{i,j,k}-\lev_{i,j,k-1}}{h_z}
  \label{eq:backward-diff}
\end{equation}
로 정의한다. 각 차분이 $1/h_x, 1/h_y, 1/h_z$로 조정된 점이 핵심이다. 간격 인지
경사는 중심 차분을 사용하여
\begin{equation}
  \gradh\lev_{i,j,k} =
  \left(
    \frac{\lev_{i+1,j,k}-\lev_{i-1,j,k}}{2h_x},\
    \frac{\lev_{i,j+1,k}-\lev_{i,j-1,k}}{2h_y},\
    \frac{\lev_{i,j,k+1}-\lev_{i,j,k-1}}{2h_z}
  \right),
  \label{eq:central-grad}
\end{equation}
간격 인지 발산은 벡터장 $V=(V_x,V_y,V_z)$에 대해 후진 차분을 사용하여
\begin{equation}
  \divh V_{i,j,k} =
  \frac{V_{x,i,j,k}-V_{x,i-1,j,k}}{h_x}
  + \frac{V_{y,i,j,k}-V_{y,i,j-1,k}}{h_y}
  + \frac{V_{z,i,j,k}-V_{z,i,j,k-1}}{h_z}
  \label{eq:backward-div}
\end{equation}
로 정의한다. 곡률 항은 $\kappa_h = \divh\bigl(\gradh\lev/\norm{\gradh\lev}\bigr)$로
이산화된다.

\paragraph{차분 조합의 근거.}
곡률 항에서 정규화된 경사장은 \emph{전진} 차분으로 만들고 발산은 \emph{후진} 차분으로
취한다. 이 짝은 임의의 선택이 아니다. 이산 부분적분 항등식에서 경계항을 제외한
부분이 연속 경우 식 \eqref{eq:proof-length-final}의 부호 구조를 보존하도록 한다.
반면 법선(식 \eqref{eq:normal-curvature})과 정사영(제 \ref{sec:projection}절)에는
중심 차분 \eqref{eq:central-grad}을 사용한다. 그 절단 오차가 $O(\hmax^2)$이며, 이것이
명제 \ref{prop:normal-error}에서 법선 각오차를 지배하는 항이다.

\begin{remark}[비등방성의 반영]
간격을 무시한 이산화는 사실상 $h_x=h_y=h_z=1$로 두는 것에 해당한다. cine CMR에서
$h_z > h_x \approx h_y$이므로 간격을 무시하면 $z$ 방향 차분이 실제보다 과대평가되어
표면 곡률과 정규화가 슬라이스 방향으로 왜곡된다. 식 \eqref{eq:central-grad},
\eqref{eq:backward-div}의 스케일링이 이 왜곡을 제거한다.
\end{remark}

\caveat{%
이 기여가 실제로 효과가 있는지는 \emph{측정해야 한다}. 구현은 모든 연산자에
\texttt{spacing\_aware} 스위치를 두어 간격을 무시한 버전을 함께 제공하며, 제
\ref{sec:ablation-design}절의 ablation은 등방 격자와 $h_z/h_x \approx 6.4$인 격자에서
두 버전의 Dice를 비교한다. 격차가 비등방성과 함께 커지지 않는다면 이 기여는 이론이
주장하는 일을 하지 않는 것이다.}

\subsection{이산 갱신식}

간격 인지 연산자를 식 \eqref{eq:gradient-flow}에 대입하고 인위적 시간을 $\Delta s$의
명시적 Euler 전진으로 이산화하면 반복 지수 $n$에 대한 갱신식
\begin{equation}
  \lev^{n+1}_{i,j,k} = \lev^{n}_{i,j,k}
  + \Delta s\,\deps\bigl(\lev^n_{i,j,k}\bigr)\left[
    \mu\,\divh\!\left(\frac{\gradh\lev^n}{\norm{\gradh\lev^n}+\varepsilon}\right)_{i,j,k}
    - \lin(\imgt{}_{,i,j,k}-\cin)^2 + \lout(\imgt{}_{,i,j,k}-\cout)^2
  \right]
  \label{eq:euler-update}
\end{equation}
을 얻는다. 각 반복에서 식 \eqref{eq:cin}, \eqref{eq:cout}으로 영역 평균을 갱신하고
식 \eqref{eq:euler-update}로 $\levt$를 갱신하는 교대 절차를 수렴할 때까지 반복한다.

\paragraph{보폭의 선택.}
식 \eqref{eq:euler-update}은 명시적 기법이므로 보폭을 고정하지 않고
\begin{equation}
  \Delta s = \mathrm{cfl}\cdot \frac{\min(\spacing)}{\max_{i,j,k}\abs{\partial_s \lev}}
  \label{eq:cfl}
\end{equation}
로 적응적으로 선택한다. 이는 한 반복에서 계면이 $\mathrm{cfl}\cdot\min(\spacing)$보다
많이 움직이지 않게 한다. 고정 보폭은 보조정리 \ref{lem:linear-conv}의 조건
$\alpha \le 1/L$을 조용히 위반할 수 있어 측정된 반복 횟수를 무의미하게 만든다.

\section{웜스타트: 두 번째 기여}
\label{sec:warmstart}

\subsection{설정}

각 프레임 $t$에 대하여 에너지 $E_t(\lev) := \ECV(\lev,\cin,\cout)$을 최소화하는 정상
상태 해를 $\lev^*_t$라 하자. 웜스타트 초기화란 프레임 $t$의 경사 하강을 시작할 때
초기값으로 이전 프레임의 수렴 해를 사용하는 것이다. 본 논문의 웜스타트는 이전 해
전체가 아니라 \emph{좁은 띠}로 제한된 초기값
\begin{equation}
  \lev^{(0)}_t = \lev^*_{t-1}
  \quad \text{on } \bigl\{\, x \in \dom : \abs{\lev^*_{t-1}(x)} < b \,\bigr\}
  \label{eq:warmstart-init}
\end{equation}
을 사용한다. $b>0$은 띠의 폭이며, 띠 밖에서는 부호만 보존하는 상수 외삽을 둔다.
이와 대비되는 콜드스타트는 프레임과 무관한 고정 초기값 $\lev_{\mathrm{cold}}$를
사용한다.

\begin{assumption}[작은 변화 가정]
\label{ass:small-change}
인접 프레임의 수렴 해가 서로 가까워, 어떤 작은 $\eta>0$에 대하여
\begin{equation}
  \norm{\lev^*_{t-1} - \lev^*_t}_{L^2(\dom)} \le \eta
  \label{eq:small-change}
\end{equation}
가 성립한다. 또한 $\lev^*_{t-1}$이 $\lev^*_t$ 주위의 국소 수렴 영역 안에 놓이며,
좁은 띠 제한 \eqref{eq:warmstart-init}이 이 근접성을 훼손하지 않는다고 가정한다.
\end{assumption}

식 \eqref{eq:small-change}의 좌변은 웜스타트의 초기 오차
$e^{\mathrm{warm}}_0 \le \eta$이다. 콜드스타트의 초기 오차
$e^{\mathrm{cold}}_0 = \norm{\lev_{\mathrm{cold}}-\lev^*_t}$은 일반적으로 이보다
훨씬 크다. 좁은 띠 초기화는 표면 근방에만 계산을 집중시켜 프레임당 비용을 줄이는
장치이며, 표면 근방에서 $\lev^*_{t-1}$과 $\lev^*_t$가 일치하므로 초기 오차의 크기를
바꾸지 않는다.

\subsection{도입하는 알려진 결과}

정리 \ref{thm:warmstart}은 새로운 수렴 정리를 처음부터 증명하는 것이 아니라, 경사
흐름과 경사 하강에 대한 알려진 결과를 웜스타트 설정에 적용한다. 어느 부분이 인용된
기지 결과이고 어느 부분이 본 논문의 특수화인지 명시한다.

\begin{lemma}[경사 흐름의 에너지 소산; 기지 결과]
\label{lem:dissipation}
$\lev(s)$가 경사 흐름 $\partial_s\lev = -\delta E_t/\delta\lev$의 매끄러운 해이면
\begin{equation}
  \frac{d}{ds}E_t(\lev(s))
  = -\norm{\frac{\delta E_t}{\delta \lev}(\lev(s))}^2 \le 0
  \label{eq:dissipation}
\end{equation}
이 성립하고, 등호는 정류점에서만 성립한다.
\end{lemma}

\begin{proof}
연쇄 법칙과 $L^2$ 내적을 사용하면
\[
  \frac{d}{ds}E_t(\lev(s))
  = \inner{\frac{\delta E_t}{\delta\lev}}{\partial_s \lev}
  = \inner{\frac{\delta E_t}{\delta\lev}}{-\frac{\delta E_t}{\delta\lev}}
  = -\norm{\frac{\delta E_t}{\delta\lev}}^2 \le 0.
\]
값이 영이 되는 것은 $\delta E_t/\delta\lev = 0$, 즉 정류점일 때뿐이다.
\end{proof}

\begin{assumption}[국소 Polyak--{\L}ojasiewicz 및 평활성]
\label{ass:pl}
조건 (i)은 Polyak~\citep{polyak1963_gradient}이 도입하고 Karimi
등~\citep{karimi2016_linear}이 경사 및 근접 경사법의 선형 수렴 조건으로 정리한
것이다.
해 $\lev^*_t$의 반지름 $r>0$ 근방
$\mathcal{B}_r = \{\lev : \norm{\lev-\lev^*_t}\le r\}$에서 다음을 가정한다.
\begin{enumerate}[label=(\roman*)]
  \item (국소 PL 조건) 어떤 상수 $m>0$에 대하여 모든 $\lev\in\mathcal{B}_r$에서
  \begin{equation}
    \tfrac{1}{2}\norm{\frac{\delta E_t}{\delta\lev}(\lev)}^2
    \ge m\bigl(E_t(\lev)-E_t(\lev^*_t)\bigr).
    \label{eq:pl}
  \end{equation}
  \item (국소 $L$-평활성) 어떤 상수 $L>0$에 대하여 $\delta E_t/\delta\lev$가
  $\mathcal{B}_r$에서 Lipschitz 상수 $L$로 Lipschitz 연속이다.
\end{enumerate}
\end{assumption}

\caveat{%
Chan--Vese 에너지 \eqref{eq:cv-energy}는 $\lev$에 대해 \textbf{전역적으로 볼록하지
않다}. 비선형 곡률 항과 매끄러운 Heaviside의 비볼록성 때문이다. 따라서 전역 강볼록성을
주장할 수 없다. 그러나 고립된 국소 최소점의 충분히 작은 근방에서는, $E_t$가 $C^2$이고
그곳의 이차 형식이 양의 정부호라는 일반적 조건 아래 가정 \ref{ass:pl}이 성립한다.
본 논문은 이 국소 조건을 가정으로 도입하며 그것이 성립하는 근방으로 논의를 제한한다.
가정 \ref{ass:small-change}은 바로 웜스타트 초기값이 이 근방에 들어옴을 보장하는
역할을 한다($\eta \le r$).}

\begin{lemma}[국소 선형 수렴; 기지 결과의 적용]
\label{lem:linear-conv}
가정 \ref{ass:pl}이 성립하고 경사 하강 반복
$\lev^{(k+1)} = \lev^{(k)} - \alpha\,\delta E_t/\delta\lev(\lev^{(k)})$를 고정 보폭
$0<\alpha\le 1/L$로 수행하며 반복열이 $\mathcal{B}_r$을 벗어나지 않는다고 하자.
그러면 에너지 잉여 $\Delta_k := E_t(\lev^{(k)})-E_t(\lev^*_t)$에 대하여
\begin{equation}
  \Delta_k \le \rho^k \Delta_0, \qquad \rho := 1-\alpha m \in (0,1)
  \label{eq:linear-conv}
\end{equation}
가 성립한다.
\end{lemma}

\begin{proof}
$L$-평활성으로부터 표준 하강 보조정리~\citep[Lemma 1.2.3]{nesterov2004_introductory}에
의해
\[
  E_t(\lev^{(k+1)}) \le E_t(\lev^{(k)})
  - \alpha\left(1-\tfrac{\alpha L}{2}\right)\norm{\frac{\delta E_t}{\delta\lev}(\lev^{(k)})}^2
  \le E_t(\lev^{(k)}) - \tfrac{\alpha}{2}\norm{\frac{\delta E_t}{\delta\lev}(\lev^{(k)})}^2,
\]
마지막 부등식은 $\alpha\le 1/L$이면 $1-\alpha L/2 \ge 1/2$임을 사용한다. 여기에 PL
조건 \eqref{eq:pl}을 대입하면
$\Delta_{k+1} \le \Delta_k - \alpha m \Delta_k = (1-\alpha m)\Delta_k$이고,
$0<\alpha m<1$이므로 귀납으로 식 \eqref{eq:linear-conv}을 얻는다.
\end{proof}

\subsection{웜스타트 수렴 정리}

\begin{theorem}[웜스타트 수렴과 반복 횟수 절감]
\label{thm:warmstart}
가정 \ref{ass:small-change}와 \ref{ass:pl}이 $\lev^*_t$의 근방 $\mathcal{B}_r$에서
성립한다고 하자. 경사 하강을 보폭 $0<\alpha\le 1/L$로 수행하고, 목표 허용 오차
$\mathrm{tol}>0$에 대해 에너지 잉여가 $\Delta_K \le \mathrm{tol}$이 되는 최소 반복
횟수를 $K(e_0)$라 하자. 여기서 $e_0 = \norm{\lev^{(0)}-\lev^*_t}$이다. 그러면 다음이
성립한다.
\begin{enumerate}[label=(\roman*)]
  \item (수렴 보존) 초기값이 $\mathcal{B}_r$ 안에 있으면 경사 하강은 $\lev^*_t$로
  선형 속도로 수렴하며, 에너지는 식 \eqref{eq:dissipation}에 의해 단조 감소한다.
  \item (정량적 반복 횟수 한계) 초기 에너지 잉여를 $L$-평활성으로 위로 어림하면
  $\Delta_0 \le \tfrac{L}{2}e_0^2$이고, 식 \eqref{eq:linear-conv}로부터
  \begin{equation}
    K(e_0) = \left\lceil \frac{\log(\Delta_0/\mathrm{tol})}{\log(1/\rho)} \right\rceil
    \le \left\lceil \frac{\log\bigl(L e_0^2/(2\,\mathrm{tol})\bigr)}{\log(1/\rho)} \right\rceil,
    \qquad \rho = 1-\alpha m.
    \label{eq:iteration-bound}
  \end{equation}
  \item (웜스타트의 이득) $e^{\mathrm{warm}}_0 \le \eta < e^{\mathrm{cold}}_0$이면
  식 \eqref{eq:iteration-bound}의 상계가 $e_0$의 단조증가 함수이므로
  \begin{equation}
    K(e^{\mathrm{warm}}_0) \le K(e^{\mathrm{cold}}_0).
    \label{eq:warmstart-gain}
  \end{equation}
\end{enumerate}
\end{theorem}

\begin{proof}
\textbf{(i)} 가정 \ref{ass:small-change}에서 $\eta \le r$이면 웜스타트 초기값
$\lev^{(0)}=\lev^*_{t-1}$은 $\mathcal{B}_r$ 안에 있다. 보조정리
\ref{lem:linear-conv}에 의해 에너지 잉여가 $\rho^k$로 감소하므로 $\Delta_k\to0$이고,
PL 조건에서 $E_t(\lev)-E_t(\lev^*_t)\ge0$이 최소점에서만 영이 되므로
$\lev^{(k)}\to\lev^*_t$이다.

\textbf{(ii)} $\delta E_t/\delta\lev(\lev^*_t)=0$이고 $\delta E_t/\delta\lev$가
$\mathcal{B}_r$에서 $L$-Lipschitz이므로 표준 이차 상계
\[
  E_t(\lev)-E_t(\lev^*_t)
  \le \inner{\frac{\delta E_t}{\delta\lev}(\lev^*_t)}{\lev-\lev^*_t}
  + \frac{L}{2}\norm{\lev-\lev^*_t}^2
  = \frac{L}{2}\norm{\lev-\lev^*_t}^2
\]
에서 $\lev=\lev^{(0)}$로 두면 $\Delta_0 \le \tfrac{L}{2}e_0^2$을 얻는다. 다음으로
$\Delta_K \le \rho^K\Delta_0 \le \mathrm{tol}$이 되기 위한 충분조건에 로그를 취하고
$\log\rho<0$임에 유의하여 정리하면
$K \ge \log(\Delta_0/\mathrm{tol})/\log(1/\rho)$이고, 이를 넘는 최소 정수가
식 \eqref{eq:iteration-bound}이다.

\textbf{(iii)} 식 \eqref{eq:iteration-bound}의 우변은 로그 안이 $e_0^2$에 비례하므로
$e_0$의 단조증가 함수이다. 따라서 결론이 따라온다.

(i)과 (ii)는 최적화 이론의 기지 결과(보조정리 \ref{lem:dissipation},
\ref{lem:linear-conv})를 인용하여 얻었고, 초기 오차를 인접 프레임 해의 근접성
\eqref{eq:small-change}과 연결하여 반복 횟수 절감 \eqref{eq:warmstart-gain}으로
특수화한 부분이 본 논문의 기여이다.
\end{proof}

\begin{remark}[선형 수렴 대 수축]
식 \eqref{eq:iteration-bound}은 에너지 잉여가 비율 $\rho$로 감소하는 선형 수렴에 대한
한계이다. 반복열 자체가 상수 $\rho'\in(0,1)$로 수축하는 경우에도 같은 로그 논법으로
$K(e_0)\le\lceil \log(e_0/\mathrm{tol})/\log(1/\rho')\rceil$을 얻으며 결론
\eqref{eq:warmstart-gain}은 동일하게 성립한다. 어느 관점에서도 더 작은 초기 오차가
더 적은 반복을 낳는다는 정성적 결론은 변하지 않는다.
\end{remark}

\caveat{%
정리 \ref{thm:warmstart}은 \emph{순서}를 예측하며 절대 반복 횟수를 예측하지 않는다.
또한 구현에서 웜스타트를 끄면 좁은 띠 crop도 함께 꺼지므로(띠가 웜스타트로부터
정의되므로), 벽시계 시간의 격차는 두 효과가 섞인다. 반복 횟수의 격차는 초기화만을
분리한다. RQ1은 두 값을 별도로 보고한다.}

\section{정준 2D Gaussian surfel}
\label{sec:canonical}

\subsection{기하}

프레임 $t$에서 surfel anchor $\anchor^t_i$가 정사영으로 $\surft$ 위에 놓인 상태라
하자. 그 점에서의 단위 법선은 간격 인지 경사로부터
\begin{equation}
  \nrm^t_i = \frac{\gradh\levt(\anchor^t_i)}{\norm{\gradh\levt(\anchor^t_i)}+\varepsilon}
  \label{eq:unit-normal}
\end{equation}
로 정의한다. 접평면 프레임 행렬은 두 접선축을 열로 하여
\begin{equation}
  \framemat^t_i = \bigl[\,e^t_{i,1}\;\; e^t_{i,2}\,\bigr] \in \Real^{3\times2},
  \qquad (\framemat^t_i)\transp \framemat^t_i = I_2,
  \qquad (\framemat^t_i)\transp \nrm^t_i = 0
  \label{eq:frame}
\end{equation}
을 만족하도록 구성한다. scale 행렬은
$\scalemat_i = \operatorname{diag}(s_{i,1},s_{i,2}) \in \Real^{2\times2}$이고, 국소
좌표 $u\in\Real^2$에 대응하는 표면 위의 점과 Gaussian 가중치는
\begin{equation}
  X^t_i(u) = \anchor^t_i + \framemat^t_i \scalemat_i u,
  \qquad
  G_i(u) = \exp\!\left(-\tfrac{1}{2}u\transp u\right)
  \label{eq:surfel-point}
\end{equation}
로 주어진다. 정의에 의해 $\{\nrm^t_i, e^t_{i,1}, e^t_{i,2}\}$은 $\anchor^t_i$에서의
오른손 정규직교 프레임을 이룬다.

식 \eqref{eq:surfel-point}에는 법선방향 두께나 3차원 공분산이 없다. 따라서
Chan--Vese의 2차원 곡면과 primitive의 차원이 직접 대응한다.

\subsection{초기 배치와 scale}

anchor는 $\surf_0$의 marching tetrahedra mesh에서 \emph{면적 균일}하게 추출한다.
면적에 비례하는 확률로 삼각형을 고르고 그 안에서 균일하게 점을 뽑으므로, 표면 위의
밀도가 고른 표본이 된다. 이후 식 \eqref{eq:projection-iter}의 정사영으로 anchor를
$\{\lev_0=0\}$ 위로 조인다. 즉 mesh는 \emph{표본 추출기}일 뿐이며 저장 표현의
일부가 아니다.

면적 $A$를 $N$개 원반이 덮을 때 평균 anchor 간격은 $\approx\sqrt{A/N}$이므로 기본
반경을
\begin{equation}
  s = \gamma\sqrt{A/N}
  \label{eq:scale-init}
\end{equation}
로 둔다. $\gamma$는 hole과 overlap 사이의 절충을 정하는 계수이며, 두 값을 모두
측정한다.

곡률 적응 확장에서는 명제 \ref{prop:planar-disk}의 평면 원반 오차
$\tfrac12\kappa s^2$이 예산 이하가 되도록 반경을 추가로 제한한다.
\begin{equation}
  s_i \le \sqrt{2\,\mathrm{budget}/\abs{\kappa_i}}.
  \label{eq:curvature-scale}
\end{equation}
이는 오차 분석을 직접적, 정량적으로 사용하는 것이다. 원반은 표면이 휘는 곳에서만
작아지고 다른 곳에서는 작아지지 않는다.

\section{법선 정사영: 세 번째 기여}
\label{sec:projection}

\subsection{최소 법선 보정의 유도}

이전 프레임의 anchor $\anchor^{t-1}_i$를 현재 표면 $\surft$로 옮기는 문제를
생각한다. 형상이 프레임 간에 조금 변하므로 이전 anchor는 일반적으로 $\surft$ 위에
없고 $\levt(\anchor^{t-1}_i)\ne0$이다.

\begin{theorem}[최소 법선 보정과 반복 투영]
\label{thm:projection}
$\levt$가 $\anchor := \anchor^{t-1}_i$ 근방에서 $C^2$이고 $\gradh\levt(\anchor)\ne0$
이라 하자. $\anchor$에서 표면 $\surft$로 향하는, 경사 방향에 평행한 일차 보정은
정규화 상수 $\varepsilon\ge0$에 대하여
\begin{equation}
  d^t_i = -\frac{\levt(\anchor^{t-1}_i)}
                {\norm{\gradh\levt(\anchor^{t-1}_i)}^2 + \varepsilon}
          \,\gradh\levt(\anchor^{t-1}_i)
  \label{eq:minimal-step}
\end{equation}
로 주어진다. 이를 반복하면 Newton 유사 투영 반복
\begin{equation}
  \anchor^{t,(k+1)}_i = \anchor^{t,(k)}_i
  - \frac{\levt\bigl(\anchor^{t,(k)}_i\bigr)}
         {\norm{\gradh\levt\bigl(\anchor^{t,(k)}_i\bigr)}^2+\varepsilon}
    \,\gradh\levt\bigl(\anchor^{t,(k)}_i\bigr),
  \qquad \anchor^{t,(0)}_i = \anchor^{t-1}_i
  \label{eq:projection-iter}
\end{equation}
을 얻으며, 이 반복은 $\surft$ 근방에서 잔차 $\levt(\anchor^{(k)})$를 영으로 국소
수렴시킨다.
\end{theorem}

\begin{proof}
$\levt$를 $\anchor$ 주위에서 일차 Taylor 전개하면 변위 $d$에 대하여
\begin{equation}
  \levt(\anchor+d) = \levt(\anchor) + \gradh\levt(\anchor)\transp d + o(\norm{d}).
  \label{eq:taylor}
\end{equation}
표면 위에 놓이는 조건 $\levt(\anchor+d)=0$을 일차항까지 부과하면
\begin{equation}
  \levt(\anchor) + \gradh\levt(\anchor)\transp d = 0.
  \label{eq:linear-constraint}
\end{equation}
식 \eqref{eq:linear-constraint}은 $d\in\Real^3$에 대한 \emph{하나의} 스칼라 방정식이므로
과소 결정된다. 표면의 형상만으로는 접선 방향 이동이 결정되지 않기 때문이다. 따라서
유일성을 위해 최소 노름 보정을 택한다.
\begin{equation}
  d = \argmin\bigl\{\,\norm{d'} \;:\; \gradh\levt(\anchor)\transp d' = -\levt(\anchor) \,\bigr\}.
  \label{eq:least-norm}
\end{equation}
임의의 허용해 $d'$을 경사 성분과 접선 성분으로
$d' = \beta\,\gradh\levt(\anchor) + w$, $w\transp\gradh\levt(\anchor)=0$로 분해하면
제약은 $\beta$만으로 결정되고 노름은
$\norm{d'}^2 = \beta^2\norm{\gradh\levt(\anchor)}^2 + \norm{w}^2$이므로 $\norm{w}$는
목적함수만 키운다. 따라서 최소해는 $w=0$, 즉 경사 방향에 평행하다.
$d=\beta\,\gradh\levt(\anchor)$를 식 \eqref{eq:linear-constraint}에 대입하면
\[
  \levt(\anchor) + \beta\norm{\gradh\levt(\anchor)}^2 = 0
  \;\Longrightarrow\;
  \beta = -\frac{\levt(\anchor)}{\norm{\gradh\levt(\anchor)}^2}.
\]
이로써 $\varepsilon=0$인 최소 법선 보정을 얻고, 경사가 작은 곳에서 분모가 영에
가까워져 보정이 폭주하는 것을 막기 위해 $\varepsilon\ge0$을 더하면 식
\eqref{eq:minimal-step}이 된다. 이 보정은 식 \eqref{eq:unit-normal}의 법선 방향에
평행하므로 법선 정사영이다.

반복 투영: 일차 보정은 일차 Taylor 근사에 기반하므로 곡률이 있는 곳에서는 한 번의
보정으로 정확히 $\surft$에 닿지 않는다. 보정을 반복하면 식
\eqref{eq:projection-iter}을 얻는다. 잔차의 이차 감소는 보조정리
\ref{lem:quadratic}에서 정량적으로 보이며, 그로부터 국소 수렴이 따라온다.
\end{proof}

\begin{proposition}[부호 거리 근방에서의 1--2회 수렴]
\label{prop:newton-steps}
$\levt$가 $\surft$ 근방에서 부호 거리 함수에 가까워 $\norm{\gradh\levt}\approx1$이
성립하면, $\varepsilon$이 무시할 만큼 작을 때 반복 \eqref{eq:projection-iter}의 한
단계는
\begin{equation}
  \anchor^{(k+1)} = \anchor^{(k)} - \levt\bigl(\anchor^{(k)}\bigr)\gradh\levt\bigl(\anchor^{(k)}\bigr)
  \label{eq:newton-step}
\end{equation}
로 축약되며, 이는 스칼라 함수 $\levt$에 대한 표준 Newton 단계와 일치한다. $\levt$가
국소적으로 거의 아핀이면 잔차는 한 번의 반복으로 거의 영이 되고, 곡률에 의한 이차
잔차는 다음 한 번의 반복에서 제거된다. 따라서 반복 횟수 $K_p$는 1 또는 2로 충분하다.
\end{proposition}

\begin{proof}
부호 거리 근방에서 $\norm{\gradh\levt}^2\approx1$이므로 $\varepsilon\to0$이면 식
\eqref{eq:projection-iter}의 분모가 1이 되어 식 \eqref{eq:newton-step}이 된다.
$\levt$가 국소적으로 아핀이면 식 \eqref{eq:taylor}의 $o(\norm{d})$ 항이 소멸하여 식
\eqref{eq:linear-constraint}이 정확한 방정식이 되고 한 번의 단계로
$\levt(\anchor^{(1)})=0$이 된다. 아핀이 아닌 경우 남는 잔차는 보조정리
\ref{lem:quadratic}에 의해 이차이므로 부호 거리 근방에서 매우 작아 추가 한 번의
반복으로 제거된다.
\end{proof}

이 성질이 구현에서 $\levt$를 주기적으로 부호 거리 함수로 재초기화하는 이유이다.
재초기화는 영도집합을 움직이지 않으면서 $\norm{\gradh\levt}\approx1$을 회복시킨다.

\subsection{접선 방향의 미결정성}
\label{sec:tangential}

식 \eqref{eq:least-norm}의 최소 노름 선택은 \emph{가장 가까운 표면으로 가는 기하학적
선택}이며 실제 심근세포 하나를 추적하는 방법이 아니다. 표면을 따라 미끄러지는 접선
운동은 두 표면의 모양만 보고는 결정할 수 없기 때문이다.

\caveat{%
그러므로 본 연구의 결과는 ``심장 모양의 시간 변화''에는 사용할 수 있으나, 추가 검증
없이 ``실제 조직의 strain''이라고 해석할 수 없다. vertex correspondence가 필요한
기능적 계측에는 mesh 운동 추적이 더 적합하다.}

또한 식 \eqref{eq:projection-iter}은 저장된 변형장이 아니라 각 프레임에서 표면이
제공하는 국소 계산이다. 따라서 위상이 바뀌거나 서로 가까운 두 표면 분기가 존재하면
잘못된 곳으로 투영될 수 있다. 구현은 최대 이동거리 제한(trust region)과 잔차가
증가한 anchor의 되돌림을 사용하여 이 실패를 줄이고, 두 사건의 발생 빈도를 보고한다.

\section{접평면 프레임 전송}
\label{sec:transport}

단위 법선 $\nrm^t_i$는 접평면을 결정하지만 그 평면 안에서의 축의 방향, 즉 면내
회전까지 고정하지는 못한다. 이 gauge 자유도를 프레임 간 연속적으로 결정하기 위해,
이전 접선축을 새 접평면으로 정사영한 뒤 정규화하는 최소 회전 전송을 사용한다.
\begin{align}
  \bar e^t_{i,1} &= e^{t-1}_{i,1} - \bigl(e^{t-1}_{i,1}\cdot \nrm^t_i\bigr)\nrm^t_i,
  \label{eq:transport-project}\\
  e^t_{i,1} &= \frac{\bar e^t_{i,1}}{\norm{\bar e^t_{i,1}}+\varepsilon},
  \qquad
  e^t_{i,2} = \nrm^t_i \times e^t_{i,1}.
  \label{eq:transport-normalize}
\end{align}
식 \eqref{eq:transport-project}은 이전 축에서 법선 성분을 제거하여 새 접평면 안에
놓이도록 하고, 식 \eqref{eq:transport-normalize}은 이를 단위 벡터로 만든 뒤 외적으로
둘째 축을 정한다. 이렇게 하면 프레임이 필요 이상으로 회전하지 않아 접선축의 시간적
연속성이 유지된다.

이는 실제 조직의 접선 이동을 추정하는 식이 아니라 rendering orientation을
연속적으로 정하는 규칙임을 명시한다. 전송의 정확성과 퇴화는 제 \ref{sec:frame-error}절에서
분석한다.

\paragraph{밀도 변화의 보정.}
법선 투영이 반복되면 수축 영역에 surfel이 몰리고 팽창 영역에 빈틈이 생길 수 있다.
이는 표면 잔차 $\Esurf$로는 드러나지 않는다. anchor가 정확히 $\surft$ 위에 있으면서도
분포가 나쁠 수 있기 때문이다. 따라서 최근접 중심거리 분포를 별도로 측정하고, 필요한
경우 접평면 안에서만 작용하는 약한 repulsion과 제한된 densification/pruning을 적용한다.
법선 성분을 제거한 뒤 이동하고 다시 정사영하므로, 이 보정이 표면 정확도를 분포
품질과 교환하지 않는다.

\section{외형 잔차}
\label{sec:residual}

표면 이동만으로 MRI 밝기 변화를 모두 설명할 수 없으므로 surfel $i$의 프레임 $t$에서의
진폭을 기준값과 잔차의 합
\begin{equation}
  \amp^t_i = \amp^0_i + \Delta \amp^t_i
  \label{eq:amplitude-split}
\end{equation}
으로 둔다. 중심, 접평면, scale을 고정하면 alpha 가중치가 주어진 2DGS 렌더링은 진폭에
대해 선형이므로(명제 \ref{prop:linearity}), 관측 pixel vector를 $y_t$, 렌더링 행렬을
$A_t$라 할 때
\begin{equation}
  \Delta \amp^*_t = \argmin_{\Delta \amp_t}
  \norm{W\bigl(A_t(\amp^0+\Delta \amp_t) - y_t\bigr)}^2_2
  + \lambda_a \norm{\Delta \amp_t}^2_2
  + \lambda_T \norm{\Delta \amp_t - \Delta \amp_{t-1}}^2_2
  \label{eq:residual-ls}
\end{equation}
를 푼다. 첫 항은 MRI appearance를 맞추고, 둘째 항은 잔차를 작게 유지하며, 셋째 항은
프레임 간 밝기 깜빡임을 억제한다.

정칙화는 선택 사항이 아니다. 알려진 slice plane에서 보이지 않는 surfel이 많으므로
$A_t$는 큰 영공간을 가지며, 자료 항만으로는 그들의 잔차가 결정되지 않는다. Tikhonov
항이 그것을 0으로 고정한다.

저장량을 더 줄이는 확장에서는
\begin{equation}
  \Delta \amp^t_i = \sum_{\ell=1}^{r} U_{i\ell} z_{\ell t},
  \qquad r \ll \min(N,T)
  \label{eq:lowrank}
\end{equation}
와 같은 저계수 시간 잔차를 사용한다. 그 절단 오차는 명제 \ref{prop:lowrank}에서
정확히 주어진다.

\section{원근 정합 래스터화}
\label{sec:raster}

단순히 3차원 Gaussian을 화면에 아핀 투영하지 않고, 각 pixel 광선과 surfel 평면의
교점을 직접 구한다. 카메라 중심을 $\camc$, pixel $q$ 방향의 \emph{단위} 벡터를
$\raydir$라 하면 광선은
\begin{equation}
  r_q(\depthq) = \camc + \depthq\,\raydir, \qquad \depthq \ge 0
  \label{eq:ray}
\end{equation}
이다. surfel $i$는 $\anchor^t_i$를 지나고 법선 $\nrm^t_i$를 갖는 평면 위에 놓이므로,
광선이 이 평면과 만나는 깊이와 교점은
\begin{equation}
  \depthq^t_i(q) = \frac{(\nrm^t_i)\transp\bigl(\anchor^t_i - \camc\bigr)}
                        {(\nrm^t_i)\transp \raydir},
  \qquad
  x^t_i(q) = \camc + \depthq^t_i(q)\,\raydir
  \label{eq:ray-plane}
\end{equation}
로 계산한다. $\raydir$가 단위이므로 $\depthq$는 mm 단위 거리이며 복셀 간격과 직접
비교 가능하다. 교점의 2차원 국소 좌표는
\begin{equation}
  u^t_i(q) = \scalemat_i^{-1}(\framemat^t_i)\transp\bigl(x^t_i(q) - \anchor^t_i\bigr),
  \label{eq:local-coord}
\end{equation}
화면 pixel에서의 불투명도는
\begin{equation}
  \alphaq^t_i(q) = \opa^t_i \exp\!\left(-\tfrac{1}{2}\norm{u^t_i(q)}^2_2\right)
  \label{eq:alpha}
\end{equation}
이다. $(\nrm^t_i)\transp\raydir$가 0에 가까운 스침각 surfel과 카메라 뒤의 교점
($\depthq \le 0$)은 제거한다. 남은 surfel을 $\depthq^t_i(q)$ 순서로 alpha
composition하면
\begin{equation}
  \Chat_t(q) = \sum_i \left(\prod_{j<i}\bigl(1-\alphaq^t_j(q)\bigr)\right)
               \alphaq^t_i(q)\,c^t_i
  \label{eq:compositing}
\end{equation}
를 얻는다. $c^t_i$는 진폭 $\amp^t_i$를 grayscale 또는 의료영상 window/level colormap으로
변환한 값이다.

\subsection{정준 프레임 fitting의 손실}

첫 프레임의 2DGS fitting에는 세 종류의 제약을 추가한다. 투영된 Chan--Vese mask와
rendered alpha mask 사이의 $\mathcal{L}_{\mathrm{mask}}$, rendered depth에서 계산한
normal과 식 \eqref{eq:unit-normal} 사이의 $\mathcal{L}_{\mathrm{normal}}$, 같은 pixel에
기여하는 surfel들의 depth가 불필요하게 벌어지는 것을 억제하는
$\mathcal{L}_{\mathrm{dist}}$이다. 전체 손실을
\begin{equation}
  \mathcal{L} = \mathcal{L}_{\mathrm{app}}
  + \lambda_m \mathcal{L}_{\mathrm{mask}}
  + \lambda_n \mathcal{L}_{\mathrm{normal}}
  + \lambda_d \mathcal{L}_{\mathrm{dist}}
  \label{eq:total-loss}
\end{equation}
로 둔다.

\caveat{%
기하가 Chan--Vese 표면에 고정되므로 이 손실은 경계를 바꿀 수 없다. 구현은 anchor,
접선축, 법선을 최적화 불가능한 버퍼로 등록하여 이를 학습률이 아니라 \emph{구조적으로}
보장한다(제 \ref{sec:impl-decisions}절). 따라서 식 \eqref{eq:total-loss}의 기하 항들은
통상의 2DGS에서와 다른 역할을 한다. 표면을 개선할 수는 없고, 고정된 표면을 원반이
덮는 방식만 개선한다.}

\subsection{재생 시 프레임 예산}

재생 중 프레임당 계산은 표면 loading, 식 \eqref{eq:projection-iter}의 소수 갱신,
식 \eqref{eq:unit-normal}--\eqref{eq:transport-normalize}의 orientation 갱신, 식
\eqref{eq:compositing}의 래스터화이다. 목표 조건은
\begin{equation}
  T_{\mathrm{frame}} = T_{\mathrm{surface}} + T_{\mathrm{project}}
  + T_{\mathrm{orient}} + T_{\mathrm{raster}} < 33.3\,\mathrm{ms}
  \label{eq:frame-budget}
\end{equation}
이다. 네 항을 각각 계측한다.

\section{저장 표현과 손익분기}
\label{sec:storage-breakeven}

모든 프레임의 2DGS를 저장하는 경우
\begin{equation}
  \Sfull = T\,N\,P_{2D}
  \label{eq:s-full}
\end{equation}
이고, 제안 표현 \eqref{eq:store}은
\begin{equation}
  \Sours = N P_{2D} + \sum_t S(\surft) + T\,N\,P_r
  \label{eq:s-ours}
\end{equation}
이다. $P_{2D}$는 한 surfel의 parameter 수, $P_r$은 잔차 parameter 수이다. 압축비는
\begin{equation}
  \CR = \frac{\Sfull}{\Sours}
  \label{eq:cr}
\end{equation}
로 계산한다.

\paragraph{손익분기 조건.}
$\Sours = \Sfull$을 프레임당 표면 예산에 대해 풀면
\begin{equation}
  S^{\ast}(\surf) = \frac{T N P_{2D} - N P_{2D} - T N P_r}{T}
  = N P_{2D}\left(1 - \frac{1}{T}\right) - N P_r
  \label{eq:breakeven}
\end{equation}
을 얻는다. 프레임당 표면 저장량이 $S^\ast$를 넘으면 $\CR<1$이 되어 전 프레임 2DGS를
저장하는 것이 오히려 싸다. 식 \eqref{eq:breakeven}이 음수이면 잔차만으로도 이미
$\Sfull$을 초과하므로 이 구성에서는 제안 표현이 이길 수 없다.

\caveat{%
따라서 RQ3은 ``압축비가 얼마인가''가 아니라 ``애초에 1을 넘는가''를 먼저 묻는다.
구현은 표면을 mesh, binary mask, 좁은 띠 부호 거리의 세 방식으로 저장할 때의 크기를
\emph{모두} 보고하고, 이론적 바이트 수와 실제 파일 크기를 나란히 제시한다.}

\section{중간 시점 보간}
\label{sec:interpolation}

원 cine CMR이 20--40 phase라 하더라도 재생을 부드럽게 하기 위해 인접 표면의 부호
거리와 잔차를 보간할 수 있다. $\tau = t+\beta$, $0\le\beta\le1$에 대하여
\begin{equation}
  \lev_\tau = (1-\beta)\levt + \beta \lev_{t+1},
  \qquad
  \Delta \amp_\tau = (1-\beta)\Delta \amp_t + \beta \Delta \amp_{t+1}
  \label{eq:interpolation}
\end{equation}
로 두고 같은 법선 정사영을 적용한다.

\caveat{%
이는 \textbf{표시를 부드럽게 하는 보간일 뿐}이며, 관측되지 않은 임상 정보를 새로
복원하는 것으로 해석하지 않는다. 명제 \ref{prop:interp-eikonal}은 보간된 level-set이
참 중간 표면의 부호 거리 함수가 아님을 정량적으로 보이고, 명제
\ref{prop:interp-projection}은 그것이 정사영에 주입하는 오차를 명시한다. 임상적으로
유효한 중간 표면이 필요하면 그 시점의 영상을 직접 분할해야 한다.}

\chapter{오차 분석}
\label{ch:error}

본 장은 제 \ref{ch:method}장의 각 구성요소가 남기는 오차를 명시적으로 한계짓고, 각
한계를 제 \ref{ch:protocol}장의 검증 지표에 대응시킨다. 네 축으로 나눈다. 접평면
프레임(제 \ref{sec:frame-error}절), 원근 정합 래스터화(제 \ref{sec:raster-error}절),
외형 잔차(제 \ref{sec:residual-error}절), 중간 시점 보간(제 \ref{sec:interp-error}절).
마지막으로 프레임 간 전파를 다룬다(제 \ref{sec:propagation}절).

\section{정사영 잔차}
\label{sec:projection-error}

\begin{lemma}[이차 잔차 한계]
\label{lem:quadratic}
$\levt$가 $\anchor^{(k)}$를 포함하는 볼록 근방에서 $C^2$이고 그 근방에서 Hessian의
노름이 $\norm{D^2\levt}\le M$으로 유계라 하자. $\varepsilon=0$인 반복
\eqref{eq:projection-iter}에 대하여
\begin{equation}
  \Bigl\lvert \levt\bigl(\anchor^{(k+1)}\bigr)\Bigr\rvert
  \le \frac{M}{2\norm{\gradh\levt(\anchor^{(k)})}^2}
      \Bigl\lvert \levt\bigl(\anchor^{(k)}\bigr)\Bigr\rvert^2
  \label{eq:quadratic-bound}
\end{equation}
가 성립한다. 따라서 $\gradh\levt$가 영에서 떨어진 표면 근방에서 반복은 잔차를 이차
속도로 영에 수렴시킨다.
\end{lemma}

\begin{proof}
표기를 줄여 $g := \gradh\levt(\anchor^{(k)})$, $r_k := \levt(\anchor^{(k)})$라 하자.
$\varepsilon=0$이면 보정은
$d^{(k)} = \anchor^{(k+1)}-\anchor^{(k)} = -r_k\,g/\norm{g}^2$이다. $\levt$의 이차
Taylor 전개를 Lagrange형 나머지항과 함께 쓰면, $\anchor^{(k)}$와 $\anchor^{(k+1)}$을
잇는 선분 위의 어떤 점 $\xi$에 대하여
\begin{equation}
  \levt\bigl(\anchor^{(k+1)}\bigr)
  = \levt\bigl(\anchor^{(k)}\bigr) + g\transp d^{(k)}
  + \tfrac{1}{2}\bigl(d^{(k)}\bigr)\transp D^2\levt(\xi)\,d^{(k)}.
  \label{eq:taylor2}
\end{equation}
일차항을 계산하면
$g\transp d^{(k)} = g\transp\bigl(-r_k g/\norm{g}^2\bigr) = -r_k\norm{g}^2/\norm{g}^2 = -r_k$
이므로 식 \eqref{eq:taylor2}의 처음 두 항 $r_k + (-r_k) = 0$이 \emph{정확히 상쇄}된다.
따라서 나머지항만 남아
\[
  \levt\bigl(\anchor^{(k+1)}\bigr)
  = \tfrac{1}{2}\bigl(d^{(k)}\bigr)\transp D^2\levt(\xi)\,d^{(k)}.
\]
Hessian 노름 유계와 Cauchy--Schwarz에 의해, 그리고
$\norm{d^{(k)}} = \abs{r_k}\norm{g}/\norm{g}^2 = \abs{r_k}/\norm{g}$를 사용하면
\[
  \Bigl\lvert\levt\bigl(\anchor^{(k+1)}\bigr)\Bigr\rvert
  \le \tfrac{1}{2}M\norm{d^{(k)}}^2
  = \tfrac{1}{2}M\frac{\abs{r_k}^2}{\norm{g}^2}
\]
이고 이것이 식 \eqref{eq:quadratic-bound}이다. 표면 근방에서 $\norm{g}\ge g_0>0$이면
상수 $C := M/(2g_0^2)$에 대해 $\abs{r_{k+1}}\le C\abs{r_k}^2$이므로, 초기 잔차가
$C\abs{r_0}<1$을 만족하면 $\abs{r_k}\to0$이 이차 속도로 성립한다.
\end{proof}

\begin{remark}[정규화 $\varepsilon$의 대가]
$\varepsilon>0$은 분모를 $\norm{g}^2+\varepsilon$으로 바꾸어 일차항 상쇄를
$g\transp d^{(k)} = -r_k\norm{g}^2/(\norm{g}^2+\varepsilon)$으로 만든다. 이 경우 처음
두 항의 합은 $r_k\varepsilon/(\norm{g}^2+\varepsilon)$로,
$\varepsilon/(\norm{g}^2+\varepsilon)$에 비례하는 \emph{일차} 잔차가 남는다.
$\varepsilon \ll \norm{g}^2$이면 이 항은 무시할 만하고 식
\eqref{eq:quadratic-bound}의 이차 거동이 지배한다. 즉 $\varepsilon$은 경사가 작은 곳의
수치 안정성을 위한 대가로 아주 작은 일차 잔차를 남긴다.
\end{remark}

반복 투영이 끝난 최종 anchor에서의 잔차는 표면 정렬 오차 지표
\begin{equation}
  \Esurf = \frac{1}{N}\sum_{i=1}^{N}\abs{\levt(\anchor^t_i)}
  \label{eq:e-surf}
\end{equation}
로 곧바로 측정된다.

\section{접평면 프레임 오차}
\label{sec:frame-error}

\subsection{단위 법선의 각오차}

\begin{proposition}[단위 법선의 각오차 한계]
\label{prop:normal-error}
점 $\anchor^t_i$에서 참 경사를 $g := \nabla\levt(\anchor^t_i)\ne0$, 그 이산 근사를
$g_h := \gradh\levt(\anchor^t_i)$라 하고 절단 오차를 $e_h := g_h - g$로 두자. 중심
차분 \eqref{eq:central-grad}의 경우
$\norm{e_h} \le C_h\,\hmax^2\,\norm{D^3\levt}_\infty$로 이차 절단 오차를 가진다고
가정한다. 참 단위 법선 $\nrm = g/\norm{g}$와 $\varepsilon$ 정규화 법선
$\nrm_\varepsilon = g_h/(\norm{g_h}+\varepsilon)$ 사이의 각 $\theta$는
\begin{equation}
  \sin\theta = \frac{\norm{g_h - (g_h\cdot \nrm)\nrm}}{\norm{g_h}}
  \le \frac{\norm{e_h}}{\norm{g_h}}
  \label{eq:sin-theta}
\end{equation}
을 만족한다. 특히 $\varepsilon$ 정규화는 벡터의 크기만 바꾸고 방향은 바꾸지 않으므로
각오차에 기여하지 않으며, 표면 근방에서 $\norm{g}\ge g_0>0$이면
\begin{equation}
  \sin\theta \le \frac{C_h \hmax^2 \norm{D^3\levt}_\infty}{g_0} = O\bigl(\hmax^2\bigr)
  \label{eq:normal-rate}
\end{equation}
로 $\hmax\to0$일 때 소멸한다.
\end{proposition}

\begin{proof}
$\nrm_\varepsilon$은 $g_h$의 양의 스칼라배이므로 그 방향은 $g_h$와 같다. 따라서
$\nrm_\varepsilon$과 $\nrm$ 사이의 각은 $g_h$와 $\nrm$ 사이의 각과 같다. 단위 벡터
$\nrm$에 대한 $g_h$의 직교 성분은 $g_h^\perp = g_h - (g_h\cdot\nrm)\nrm$이고, 정의에
의해 $\sin\theta = \norm{g_h^\perp}/\norm{g_h}$이다. 이제 $g = \norm{g}\nrm$이므로
$\nrm$에 수직인 성분만 보면 $g^\perp=0$이고 따라서
\[
  g_h^\perp = (g+e_h)^\perp = g^\perp + e_h^\perp = e_h^\perp,
\]
즉 $g_h$의 수직 성분은 오차 $e_h$의 수직 성분과 같다. 그러므로
$\norm{g_h^\perp} = \norm{e_h^\perp} \le \norm{e_h}$이고 식 \eqref{eq:sin-theta}이
성립한다. 역삼각부등식에 의해 $\norm{g_h}\ge\norm{g}-\norm{e_h}$이므로 표면 근방에서
$\norm{e_h}\ll\norm{g}$일 때 최고차 항에서 $\sin\theta \le \norm{e_h}/\norm{g}$의
형태가 되고, 가정한 이차 절단 오차와 $\norm{g}\ge g_0$을 대입하면 식
\eqref{eq:normal-rate}을 얻는다. $\levt$가 표면 근방에서 부호 거리 함수에 가까우면
$\norm{g}\approx1$이므로 $g_0\approx1$로 두어도 좋다.
\end{proof}

\caveat{%
식 \eqref{eq:normal-rate}은 \emph{수렴률}에 대한 주장이므로 단일 해상도에서는 검증할
수 없다. 제 \ref{sec:kernel-verification}절의 검증은 이 항목을 포함하지 않으며,
$\spacing$을 변화시키며 $\log\sin\theta$ 대 $\log\hmax$의 기울기를 적합하는 실험이
필요하다. 합성 팬텀이 존재하는 이유가 바로 이것이다. 실제 데이터셋은 정확한 법선을
제공하지 않고 $\spacing$을 임의로 바꿀 수도 없다.}

\subsection{전송의 정규직교성}

\begin{lemma}[접평면 프레임의 정규직교성]
\label{lem:transport-orthonormal}
단위 법선 $\nrm\in\Real^3$($\norm{\nrm}=1$)과 임의의 이전 축
$e^{t-1}_1\in\Real^3$이 주어졌다고 하자. $\varepsilon=0$인 전송
\eqref{eq:transport-project}--\eqref{eq:transport-normalize}에서
$\bar e^t_1 = e^{t-1}_1 - (e^{t-1}_1\cdot\nrm)\nrm \ne 0$이면, 그 결과
$e^t_1 = \bar e^t_1/\norm{\bar e^t_1}$과 $e^t_2 = \nrm\times e^t_1$으로 이루어진
$\framemat^t = [\,e^t_1\;e^t_2\,]$은
\begin{equation}
  (\framemat^t)\transp\framemat^t = I_2, \qquad (\framemat^t)\transp\nrm = 0
  \label{eq:orthonormal}
\end{equation}
을 만족하며, $\{\nrm, e^t_1, e^t_2\}$은 오른손 정규직교 기저이다.
\end{lemma}

\begin{proof}
Gram--Schmidt 방식으로 단계별로 확인한다.

\emph{(i) $e^t_1\cdot\nrm=0$.} $\norm{\nrm}^2=1$이므로
\[
  \bar e^t_1\cdot\nrm
  = e^{t-1}_1\cdot\nrm - (e^{t-1}_1\cdot\nrm)(\nrm\cdot\nrm)
  = e^{t-1}_1\cdot\nrm - e^{t-1}_1\cdot\nrm = 0.
\]
$e^t_1$은 $\bar e^t_1$의 양의 스칼라배이므로 $e^t_1\cdot\nrm=0$이다.

\emph{(ii) $\norm{e^t_1}=1$.} $\bar e^t_1\ne0$이므로 정규화가 잘 정의된다.

\emph{(iii) $e^t_2 = \nrm\times e^t_1$의 성질.} $\nrm$과 $e^t_1$은 서로 직교하는 단위
벡터이므로 외적의 노름 공식에서 $\sin\angle=1$이고 $\norm{e^t_2}=1$이다. 또한 외적은
두 피연산자 모두에 직교하므로 $e^t_2\cdot\nrm=0$, $e^t_2\cdot e^t_1=0$이다.

\emph{(iv) 결론.} (i)--(iii)에서 식 \eqref{eq:orthonormal}이 성립하고,
$\{\nrm,e^t_1,e^t_2=\nrm\times e^t_1\}$은 순서상 오른손 규칙을 따른다.
\end{proof}

\begin{remark}[$\varepsilon>0$의 편차]
실제 사용하는 정규화는 $e^t_1 = \bar e^t_1/(\norm{\bar e^t_1}+\varepsilon)$이다. 이때
열의 크기는
$\norm{e^t_1} = \norm{\bar e^t_1}/(\norm{\bar e^t_1}+\varepsilon) = 1 - \varepsilon/(\norm{\bar e^t_1}+\varepsilon)$
로 단위 크기에서 $O(\varepsilon/\norm{\bar e^t_1})$만큼 벗어난다. 반면 직교성
$e^t_1\cdot\nrm=0$은 스칼라 정규화가 방향을 바꾸지 않으므로 \emph{정확히} 유지되며,
$e^t_2=\nrm\times e^t_1$ 역시 방향이 보존되어 직교성이 정확하다. 즉 $\varepsilon>0$은
열의 길이에만 일차 오차를 주고 직교성에는 영향을 주지 않는다. 필요하면 정규화 후 한
번 더 나누어 단위 크기를 회복할 수 있으며, 구현은 이를 기본으로 수행한다.
\end{remark}

\subsection{등방 원반에서의 gauge 불변성}

\begin{proposition}[등방 원반에서 면내 회전 gauge의 무해성]
\label{prop:gauge}
surfel이 등방, 즉 $s_{i,1}=s_{i,2}=s$이어서 $\scalemat_i = s I_2$라 하자. 임의의
$R\in SO(2)$에 대하여 접평면 프레임을 $\framemat^t_i \mapsto \framemat^t_i R$로 바꾸면,
(a) 표면 점 집합 $\{X(u) = \anchor+\framemat\scalemat u : u\in\Real^2\}$은 집합으로서
불변이고, (b) Gaussian 가중치가 재매개변수화 아래 보존되며, (c) 불투명도에 들어가는
국소 좌표의 노름 $\norm{u}$이 보존된다.
\end{proposition}

\begin{proof}
등방이므로 $\scalemat = sI_2$이고 $R\transp R = I_2$, $\norm{R\transp u}=\norm{u}$이다.

(a) 바뀐 원반의 점 집합은 $\{\anchor + s\framemat R u : u\in\Real^2\}$이다. $R$이
$\Real^2$의 전단사이므로 $\{Ru : u\in\Real^2\} = \Real^2$이고 따라서
$\{\anchor+s\framemat Ru\} = \{\anchor+s\framemat v\}$로 원래 집합과 같다.

(b) 바뀐 프레임에서 점 $X'(u)$에 대응하는 원래 좌표는 $v = Ru$이고
\[
  G(v) = \exp\bigl(-\tfrac12 v\transp v\bigr)
  = \exp\bigl(-\tfrac12 u\transp R\transp R u\bigr)
  = \exp\bigl(-\tfrac12 u\transp u\bigr) = G(u).
\]

(c) 식 \eqref{eq:local-coord}에서 등방이면 $\scalemat^{-1}=s^{-1}I_2$가 $R$과 가환이므로
\[
  u' = \scalemat^{-1}(\framemat R)\transp(x-\anchor)
     = s^{-1}R\transp \framemat\transp(x-\anchor) = R\transp u,
\]
이고 $\norm{u'}=\norm{R\transp u}=\norm{u}$이다. 따라서 식 \eqref{eq:alpha}의
불투명도는 불변이다.
\end{proof}

명제 \ref{prop:gauge}이 실용적으로 중요한 이유는 두 가지이다. 첫째, 등방 원반을 쓰면
접선 전송의 퇴화 실패 모드가 \emph{원리적으로 사라진다}. 이것이 최소 구현에서 등방
원반을 택하는 근거이다. 둘째, 퇴화 시 접평면 안의 임의의 정규직교 축을 골라도
무해하다는 결론의 근거가 된다.

\subsection{비등방 원반에서의 잔차 한계}

\begin{proposition}[비등방 원반에서의 정렬 오차 한계]
\label{prop:aniso}
surfel이 비등방이어서 $s_1\ne s_2$라 하고, 실제 사용된 프레임이 이상적 프레임을 각
$\delta$만큼 면내 회전 $R_\delta\in SO(2)$한 것이라 하자. 한 표면 점의 이상적 국소
좌표를 $u$, 어긋난 프레임에서의 좌표를 $u_\delta = \scalemat^{-1}R_\delta\scalemat u$라
할 때
\begin{equation}
  \Bigl\lvert \norm{u_\delta}^2 - \norm{u}^2 \Bigr\rvert
  \le \left\lvert \frac{s_1}{s_2}-\frac{s_2}{s_1}\right\rvert \abs{\sin\delta}\,\norm{u}^2
  \label{eq:aniso-bound}
\end{equation}
로 유계이며, 이에 따라 불투명도의 상대 오차도 같은 크기로 유계이다. 등방 극한
$s_1\to s_2$ 또는 정렬 극한 $\delta\to0$에서 우변은 영으로 수렴한다.
\end{proposition}

\begin{proof}
$M := \scalemat^{-1}R_\delta\scalemat$로 두면 $u_\delta = Mu$이고
$\norm{u_\delta}^2-\norm{u}^2 = u\transp(M\transp M - I_2)u$이다. 직접 계산하면
\[
  M = \begin{pmatrix}
    \cos\delta & -\frac{s_2}{s_1}\sin\delta\\[0.3em]
    \frac{s_1}{s_2}\sin\delta & \cos\delta
  \end{pmatrix},
\]
따라서 대칭행렬 $A := M\transp M - I_2$의 성분은
\[
  A = \begin{pmatrix}
    \bigl((s_1/s_2)^2-1\bigr)\sin^2\delta
      & \bigl(\frac{s_1}{s_2}-\frac{s_2}{s_1}\bigr)\sin\delta\cos\delta\\[0.3em]
    \bigl(\frac{s_1}{s_2}-\frac{s_2}{s_1}\bigr)\sin\delta\cos\delta
      & \bigl((s_2/s_1)^2-1\bigr)\sin^2\delta
  \end{pmatrix}.
\]
비대각 성분의 크기는
$\bigl\lvert\frac{s_1}{s_2}-\frac{s_2}{s_1}\bigr\rvert\abs{\sin\delta\cos\delta}
\le \bigl\lvert\frac{s_1}{s_2}-\frac{s_2}{s_1}\bigr\rvert\abs{\sin\delta}$이고 대각
성분은 $O(\sin^2\delta)$의 고차항이므로, 작은 $\delta$에 대하여 $A$의 스펙트럼 노름은
$\norm{A} \le \bigl\lvert\frac{s_1}{s_2}-\frac{s_2}{s_1}\bigr\rvert\abs{\sin\delta}
+ O(\sin^2\delta)$로 지배된다. Rayleigh 몫 부등식
$\abs{u\transp A u}\le\norm{A}\norm{u}^2$~\citep[제 4.2절]{horn2013_matrix}에서 일차
항까지 식 \eqref{eq:aniso-bound}을 얻는다.
\end{proof}

명제 \ref{prop:aniso}은 비등방 원반일수록 접선 전송의 정확성이 중요함을 정량적으로
말해 준다. 이방성 비 $s_1/s_2$가 1에 가까우면 전송 오차는 무해하다.

\subsection{전송의 퇴화}

\begin{proposition}[전송 퇴화와 조건수]
\label{prop:degeneracy}
이전 접선 축 $e^{t-1}_1$(단위 벡터)과 새 법선 $\nrm^t_i$ 사이의 각을 $\psi$라 하자.
그러면 정사영된 벡터의 노름은
\begin{equation}
  \norm{\bar e^t_{i,1}} = \sqrt{1-\bigl(e^{t-1}_1\cdot\nrm^t_i\bigr)^2} = \abs{\sin\psi}
  \label{eq:sin-psi}
\end{equation}
이다. 따라서 $e^{t-1}_1$이 $\nrm^t_i$과 거의 평행($\psi\to0$)해지면
$\norm{\bar e^t_{i,1}}\to0$이 되어 정규화된 축의 방향이 병적으로 민감해진다.
정량적으로, $\bar e^t_{i,1}$에 크기 $\gamma$의 섭동이 가해지면 $\varepsilon=0$일 때
정규화된 축의 방향 변화는
\begin{equation}
  \norm{\Delta e^t_1} \lesssim \frac{\gamma}{\norm{\bar e^t_{i,1}}} = \frac{\gamma}{\abs{\sin\psi}}
  \label{eq:amplification}
\end{equation}
로 증폭된다.
\end{proposition}

\begin{proof}
$e^{t-1}_1$은 단위 벡터이고 정사영 $\bar e^t_{i,1}$은 $\nrm$에 직교하는 성분이므로
피타고라스 정리에 의해
$\norm{\bar e^t_{i,1}}^2 = 1-(e^{t-1}_1\cdot\nrm)^2$이다. $e^{t-1}_1\cdot\nrm=\cos\psi$
이므로 $\norm{\bar e^t_{i,1}}^2 = \sin^2\psi$, 즉 식 \eqref{eq:sin-psi}이다. 정규화 사상
$v\mapsto v/\norm{v}$의 Jacobian은 $\frac{1}{\norm{v}}\bigl(I - \frac{vv\transp}{\norm{v}^2}\bigr)$
이고 그 접선 성분에 대한 연산자 노름은 $1/\norm{v}$이다. 따라서 식
\eqref{eq:amplification}이 성립한다.
\end{proof}

\begin{remark}[$\varepsilon$ 분모의 정당화와 실용적 처리]
식 \eqref{eq:transport-normalize}의 분모 $\norm{\bar e^t_{i,1}}+\varepsilon$은 위
증폭을 $1/\varepsilon$으로 위에서 막아 조건수를 유계로 만든다. 실용적 처리는 두
가지이다. 첫째, $\norm{\bar e^t_{i,1}}$이 임계값보다 작으면 대신 다른 이전 축
$e^{t-1}_2$를 정사영하여 첫 축으로 삼는다. 이 대안은 \emph{항상} 사용 가능하다.
정규직교 쌍에 대해
\[
  \norm{\bar e_1}^2 + \norm{\bar e_2}^2
  = 2 - (e_1\cdot\nrm)^2 - (e_2\cdot\nrm)^2 \ge 1
\]
이므로 $\max(\norm{\bar e_1},\norm{\bar e_2}) \ge 1/\sqrt2$이고, 둘이 동시에 퇴화할 수
없다. 둘째, 등방 원반에서는 명제 \ref{prop:gauge}에 의해 면내 축의 선택이 무해하므로
퇴화 시 접평면 안의 임의의 정규직교 축을 골라도 렌더링에 영향이 없다.
\end{remark}

\section{원근 정합 래스터화 오차}
\label{sec:raster-error}

\subsection{광선-평면 교차의 적정성}

\begin{proposition}[교차의 적정성과 조건수]
\label{prop:ray-plane}
교차 깊이 \eqref{eq:ray-plane}은 분모 $(\nrm^t_i)\transp\raydir \ne 0$일 때에만 잘
정의된다. 광선과 원반 평면 사이의 각을 $\psi\in[0,\pi/2]$라 하면
$(\nrm^t_i)\transp\raydir = \pm\sin\psi$이고, 광선이 평면에 나란해질수록
$\depthq^t_i(q)$가 발산한다. 나아가 법선의 섭동 $\delta \nrm$이나 광선의 섭동
$\delta d$에 대한 깊이의 일차 민감도는
\begin{equation}
  \abs{\delta \depthq^t_i(q)}
  \lesssim \frac{C\bigl(\norm{\delta \nrm}+\norm{\delta d}\bigr)}
                {\bigl\lvert(\nrm^t_i)\transp\raydir\bigr\rvert}
  \label{eq:depth-conditioning}
\end{equation}
로, 조건수가 $1/\lvert(\nrm^t_i)\transp\raydir\rvert$에 비례한다. 따라서 컷오프
$\lvert(\nrm^t_i)\transp\raydir\rvert \ge c_0 > 0$을 두면 조건수는 $1/c_0$으로 유계가
된다.
\end{proposition}

\begin{proof}
$\nrm^t_i$와 $\raydir$가 모두 단위 벡터이므로 둘이 이루는 각을 $\theta_{nd}$라 하면
$(\nrm^t_i)\transp\raydir = \cos\theta_{nd}$이고, 광선과 평면이 이루는 각은
$\psi = \pi/2-\theta_{nd}$이므로 $\cos\theta_{nd}=\sin\psi$이다. 민감도를 보기 위해
$\depthq = A/B$로 두자. 여기서 $A = (\nrm^t_i)\transp(\anchor^t_i-\camc)$,
$B = (\nrm^t_i)\transp\raydir$이다. 몫의 미분에서
\[
  \delta\depthq = \frac{\delta A}{B} - \frac{A\,\delta B}{B^2}
  = \frac{1}{B}\bigl(\delta A - \depthq\,\delta B\bigr).
\]
$\abs{\delta A} \le \norm{\anchor^t_i-\camc}\norm{\delta \nrm}$,
$\abs{\delta B} \le \norm{\delta \nrm}+\norm{\delta d}$이고 시야 안에서
$\norm{\anchor^t_i-\camc}$와 $\abs{\depthq}$가 유계이므로, 유계 상수를 $C$로 묶으면
식 \eqref{eq:depth-conditioning}이 성립한다.
\end{proof}

\subsection{국소 좌표와 불투명도의 오차}

\begin{proposition}[국소 좌표와 불투명도의 오차 한계]
\label{prop:local-coord}
교차점 $x^t_i(q)$의 오차를 $\delta x$, 그에 따른 국소 좌표 \eqref{eq:local-coord}의
오차를 $\delta u$라 하자. 접평면 프레임이 정규직교이므로 $(\framemat^t_i)\transp$은
노름을 늘리지 않고, 따라서
\begin{equation}
  \norm{\delta u} \le \norm{\scalemat_i^{-1}}\norm{\delta x}
  = \frac{\norm{\delta x}}{\min\{s_{i,1},s_{i,2}\}}.
  \label{eq:local-coord-error}
\end{equation}
나아가 불투명도 \eqref{eq:alpha}의 오차는
\begin{equation}
  \abs{\delta \alphaq^t_i(q)} \le \opa^t_i\,\norm{u}\,\norm{\delta u}
  \label{eq:alpha-error}
\end{equation}
로 유계이다.
\end{proposition}

\begin{proof}
식 \eqref{eq:local-coord}에서 $\delta u = \scalemat_i^{-1}(\framemat^t_i)\transp\delta x$
이다. 정규직교성에서 $(\framemat^t_i)\transp$의 스펙트럼 노름은 1이므로
$\norm{(\framemat^t_i)\transp \delta x}\le\norm{\delta x}$이고,
$\scalemat_i^{-1}$의 스펙트럼 노름은 $1/\min\{s_{i,1},s_{i,2}\}$이다. 두 부등식을
결합하면 식 \eqref{eq:local-coord-error}을 얻는다. 작은 scale이 좌표 오차를 증폭함을
알 수 있다.

불투명도는 $f(w) = \opa\,e^{-w/2}$($w=\norm{u}^2\ge0$)의 꼴이며
$\abs{f'(w)}\le\tfrac12\opa$이다. 평균값 정리에 의해
$\abs{\delta\alphaq}\le\tfrac12\opa\abs{\delta w}$이고
$\delta w = 2u\transp\delta u + \norm{\delta u}^2$의 일차 항에서
$\abs{\delta w}\le2\norm{u}\norm{\delta u}$이므로 식 \eqref{eq:alpha-error}을 얻는다.
\end{proof}

교차점 오차 $\delta x$ 자체는 법선 각오차(명제 \ref{prop:normal-error})와 anchor 잔차
(보조정리 \ref{lem:quadratic})로 지배된다. 즉 오차가 기하에서 렌더링으로 전파되는
경로가 명시적으로 추적된다.

\subsection{평면 원반의 곡면 근사 오차}

\begin{proposition}[평면 원반의 곡률 기반 근사 오차]
\label{prop:planar-disk}
surfel은 참 표면 $\surft$를 anchor에서의 접평면으로 근사한다. 그 점에서 표면의
주곡률 크기의 최댓값을 $\kappa$라 하고, 원반의 반경을
$s=\max\{s_{i,1},s_{i,2}\}$ 규모라 하자. 접평면을 기준면으로 하는 높이 함수
$\eta$로 표면을 나타내면, 접점에서 옆으로 거리 $\varrho$ 떨어진 곳에서
\begin{equation}
  \eta(\varrho) = \tfrac{1}{2}\kappa\,\varrho^2 + O(\varrho^3)
  \label{eq:height-function}
\end{equation}
이 성립하여, 참 표면과 원반 사이의 높이 오차는 원반 자락에서 $O(\kappa s^2)$이고 법선
방향 편차는 $O(\kappa s)$이다.
\end{proposition}

\begin{proof}
접평면을 국소 좌표계의 기준면으로, 법선을 $z$축으로 잡으면 표면이 접평면에 접하므로
높이 함수 $\eta$는 $\eta(0)=0$, $\nabla\eta(0)=0$을 만족한다. $\eta$의 이차 Taylor
전개는 $\eta(\zeta) = \tfrac12\zeta\transp H\zeta + O(\norm{\zeta}^3)$, $\zeta\in\Real^2$
이며 $H$는 제2기본형식에 해당하는 대칭행렬로 그 고윳값이 주곡률이다. 반경 방향 거리
$\varrho=\norm{\zeta}$에 대해 Rayleigh 몫 부등식
$\abs{\zeta\transp H\zeta}\le\norm{H}\norm{\zeta}^2 = \kappa\varrho^2$을 쓰면 식
\eqref{eq:height-function}이다. 원반 자락은 $\varrho\lesssim s$이므로 높이 오차는
$O(\kappa s^2)$이다. 법선 방향 편차는 기울기
$\nabla\eta(\zeta) = H\zeta + O(\norm{\zeta}^2)$의 크기가 $O(\kappa s)$이므로 접평면
법선과 참 법선 사이의 각도 편차도 $O(\kappa s)$이다.
\end{proof}

명제 \ref{prop:planar-disk}이 식 \eqref{eq:curvature-scale}의 곡률 적응 scale의 근거이다.
원반 크기를 $s \ll 1/\kappa$로 유지하면 근사 오차가 제어된다.

\subsection{원근 정합 대 아핀 투영}

\begin{proposition}[원근 정합의 정확성]
\label{prop:perspective}
평면 원반에 대하여 원근 정합 교차 \eqref{eq:ray-plane}--\eqref{eq:local-coord}은 국소
좌표를 오차 없이 정확하게 계산한다. 반면 원반을 이미지 평면에 아핀 근사로 투영하면
자락 전역에 걸친 깊이 변화를 무시하는 데서 오차가 생기며, 그 크기는 자락에 걸친 상대
깊이 변화 $\Delta z/z$에 지배된다. 따라서 남는 오차는 곡면 근사 오차
$O(\kappa s^2)$(명제 \ref{prop:planar-disk})뿐이다.
\end{proposition}

\begin{proof}
surfel이 평면 원반이므로 그 위의 점은 정확히 평면 방정식
$(\nrm^t_i)\transp(x-\anchor^t_i)=0$을 만족한다. 화소 광선 \eqref{eq:ray}을 이 평면과
교차시키는 식 \eqref{eq:ray-plane}은 이 방정식의 정확한 해이므로, 교차점과 그로부터
얻는 국소 좌표에는 투영 모델에서 오는 오차가 없다.

아핀 근사는 화면 좌표를 anchor의 투영점 둘레에서 일차로 전개하여 자락 전역에 대해
anchor의 깊이 $z_0$를 상수로 고정한다. 그러나 자락의 실제 점은 깊이 $z = z_0+\Delta z$를
가지며, 원근 투영의 화면 좌표는 깊이에 반비례하므로 일차 전개가 무시한 항은
\[
  \frac{1}{z_0+\Delta z} - \frac{1}{z_0}\left(1-\frac{\Delta z}{z_0}\right)
  = O\!\left(\Bigl(\frac{\Delta z}{z_0}\Bigr)^2\right),
\]
곧 상대 깊이 변화의 크기에 지배되는 오차를 남긴다. $\Delta z$는 surfel이 광축에 대해
기울수록, 그리고 자락이 클수록 커지므로 아핀 오차는 가파른 각도와 큰 자락에서
증가한다. 원근 정합 교차는 각 화소마다 참 깊이를 직접 풀어 $\Delta z$를 정확히
반영하므로 이 항이 사라진다.
\end{proof}

이것이 thin-3DGS 기준선과의 비교(RQ6)가 측정하는 바이다. 그 기준선에서 깊이는 자락
전역에 걸쳐 상수이므로, 깊이 지표의 차이는 명제 \ref{prop:perspective}의 직접적
검증이 된다.

\section{외형 잔차의 오차}
\label{sec:residual-error}

\subsection{진폭 선형성}

\begin{proposition}[고정 기하에서 렌더링은 진폭에 아핀]
\label{prop:linearity}
surfel의 중심, 접평면 프레임, scale이 모두 고정되어 있으면 각 불투명도
$\alphaq^t_i(q)$는 진폭에 무관하고, alpha 합성 결과 $\Chat_t(q)$는 surfel 색 $c^t_i$에
대해 선형이다. 색이 진폭에 아핀이면 $\Chat_t(q)$는 진폭 벡터의 아핀 함수이며, 상수를
흡수하면 렌더링 사상은 행렬 $A_t$로 표현되는 선형 사상이 된다.
\end{proposition}

\begin{proof}
불투명도 \eqref{eq:alpha}은 국소 좌표 \eqref{eq:local-coord}과 불투명도 계수
$\opa^t_i$에만 의존하는데, 이들은 기하와 광선으로만 결정되고 진폭과 무관하다. 따라서
기하가 고정되면 $\alphaq^t_i(q)$와 그로부터 만들어지는 투과율 가중치
\begin{equation}
  w_i(q) := \left(\prod_{j<i}\bigl(1-\alphaq^t_j(q)\bigr)\right)\alphaq^t_i(q)
  \label{eq:weights}
\end{equation}
은 진폭에 무관한 상수이다. 식 \eqref{eq:compositing}은
$\Chat_t(q) = \sum_i w_i(q)\,c^t_i$이므로 $c^t_i$에 대해 선형이다. 색이 진폭에 아핀,
$c^t_i = \gamma_i \amp^t_i + \beta_i$이면
\[
  \Chat_t(q) = \sum_i w_i(q)\gamma_i\,\amp^t_i + \sum_i w_i(q)\beta_i
\]
이고 첫 항은 진폭에 선형, 둘째 항은 상수이다. 모든 화소에 대해 이 관계를 모으면
상수항을 흡수한 뒤 $\Chat_t = A_t \amp$의 꼴이 되며, $A_t$의 성분은
$w_i(q)\gamma_i$로 주어지는 진폭에 무관한 상수이다.
\end{proof}

명제 \ref{prop:linearity}이 식 \eqref{eq:residual-ls}을 \emph{선형화가 아니라 진짜
선형} 최소제곱 문제로 만든다. 구현은 이를 활용하여 $A_t$를 희소 행렬로 명시적으로
추출하고, 매 CG 반복을 전체 렌더링 대신 두 번의 색인 누적으로 수행한다.

\subsection{정칙화 최소제곱의 존재성과 조건수}

\begin{proposition}[존재성, 유일성, 조건수]
\label{prop:spd}
$\lambda_a + \lambda_T > 0$이면 식 \eqref{eq:residual-ls}의 목적함수는
$\Delta \amp_t$에 대한 강볼록 이차식이고, 그 정규방정식은
\begin{equation}
  \underbrace{\Bigl(A_t\transp W\transp W A_t + (\lambda_a+\lambda_T)I\Bigr)}_{=:M}
  \Delta \amp_t
  = A_t\transp W\transp W\,(y_t - A_t \amp^0) + \lambda_T \Delta \amp_{t-1}
  \label{eq:normal-equations}
\end{equation}
이다. $M$은 대칭 양의 정부호이므로 최소해가 유일하게 존재한다. 나아가 그 조건수는
\begin{equation}
  \operatorname{cond}(M)
  \le \frac{\sigma_{\max}\bigl(A_t\transp W\transp W A_t\bigr) + \lambda_a + \lambda_T}
           {\lambda_a+\lambda_T}
  \label{eq:cond-bound}
\end{equation}
로 유계이다.
\end{proposition}

\begin{proof}
목적함수를 전개하면 이차항은
$\Delta \amp_t\transp\bigl(A_t\transp W\transp W A_t + (\lambda_a+\lambda_T)I\bigr)\Delta \amp_t$
이다. $A_t\transp W\transp W A_t = (WA_t)\transp(WA_t)$은 임의의 $v$에 대해
$v\transp(WA_t)\transp(WA_t)v = \norm{WA_t v}^2 \ge 0$이므로 양의 준정부호이다. 여기에
$(\lambda_a+\lambda_T)I$를 더하면 임의의 $v\ne0$에 대해
\[
  v\transp M v = \norm{WA_t v}^2 + (\lambda_a+\lambda_T)\norm{v}^2
  \ge (\lambda_a+\lambda_T)\norm{v}^2 > 0,
\]
이므로 $M$은 대칭 양의 정부호이다. 강볼록 이차식의 최소는 기울기가 영인 점에서
유일하게 달성되며, 기울기를 영으로 두면 식 \eqref{eq:normal-equations}을 얻는다.

조건수는 대칭 양정부호 행렬에서
$\operatorname{cond}(M)=\lambda_{\max}(M)/\lambda_{\min}(M)$이다.
$A_t\transp W\transp W A_t$의 고윳값이 $[0,\sigma_{\max}]$에 있으므로 $M$의 고윳값은
$[\lambda_a+\lambda_T,\ \sigma_{\max}+\lambda_a+\lambda_T]$에 있고 식
\eqref{eq:cond-bound}이 따라온다.
\end{proof}

$M$이 SPD임은 공액경사법~\citep[제 5장]{nocedal2006_numerical}이 필요로 하는 정확한
조건이다.
$\lambda_a+\lambda_T=0$이면 반정부호에 그쳐 CG가 정체한다. 식 \eqref{eq:cond-bound}은
정칙화가 조건수에서 사는 것과 자료 적합에서 잃는 편향의 교환을 명시한다.

\subsection{저계수 절단 오차}

\begin{proposition}[저계수 절단 오차]
\label{prop:lowrank}
전체 잔차 행렬 $\Delta A = [\Delta \amp^t_i] \in \Real^{N\times T}$(행은 surfel, 열은
프레임)의 특잇값을 $\sigma_1\ge\sigma_2\ge\cdots\ge0$이라 하자. 식 \eqref{eq:lowrank}의
계수 $r$ 근사에 대하여 최적 근사의 오차는 버려진 특잇값의 꼬리로 \emph{정확히}
주어진다.
\begin{equation}
  \min_{\operatorname{rank}(\Delta A_r)\le r}\norm{\Delta A - \Delta A_r}_F^2
  = \sum_{\ell>r}\sigma_\ell^2,
  \qquad
  \min_{\operatorname{rank}(\Delta A_r)\le r}\norm{\Delta A - \Delta A_r}_2 = \sigma_{r+1}.
  \label{eq:eckart-young}
\end{equation}
\end{proposition}

\begin{proof}
Eckart--Young--Mirsky 정리~\citep{eckart1936_approximation}, 즉 절단 특잇값 분해의
최적성~\citep[제 2.4절]{golub2013_matrix}에 의해, 특잇값 분해
$\Delta A = \sum_\ell \sigma_\ell u_\ell v_\ell\transp$에 대하여 계수 $r$ 이하의 모든
행렬 가운데 Frobenius 및 스펙트럼 노름 오차를 최소화하는 것은 상위 $r$개 특잇삼중항을
남긴 $\Delta A_r = \sum_{\ell\le r}\sigma_\ell u_\ell v_\ell\transp$이며, 이때 잔차는
$\sum_{\ell>r}\sigma_\ell u_\ell v_\ell\transp$이다. 특잇벡터들이 정규직교하므로
Frobenius 노름의 제곱은 특잇값 제곱의 합이 되고, 스펙트럼 노름은 남은 최대 특잇값
$\sigma_{r+1}$이다.
\end{proof}

명제 \ref{prop:lowrank}의 실용적 의미는 압축 오차가 \emph{사전에 관측 가능}하다는
것이다. 잔차 행렬의 특잇값만 계산하면 어떤 계수에서 어떤 오차가 나는지 미리 알 수
있으므로, 계수를 조정하는 대신 목표 오차로부터 역산할 수 있다.

\caveat{%
그러나 실제 심장 자료의 잔차장이 얼마나 저계수인지, 즉 특잇값이 실제로 얼마나 빠르게
감소하는지는 \emph{경험적 질문}이다. 명제 \ref{prop:lowrank}은 오차가 꼬리와 같음을
말할 뿐 꼬리가 작음을 말하지 않는다. 구현은 단일 숫자가 아니라 전체 스펙트럼을
보고한다.}

\section{중간 시점 보간의 오차}
\label{sec:interp-error}

\begin{proposition}[보간된 level-set의 eikonal 위반]
\label{prop:interp-eikonal}
$\levt$와 $\lev_{t+1}$이 각각 정확한 부호 거리 함수여서
$\norm{\nabla\levt}=\norm{\nabla\lev_{t+1}}=1$이라 하자. 그러면 볼록 결합
\eqref{eq:interpolation}의 경사 제곱 노름은
\begin{equation}
  \norm{\nabla\lev_\tau}^2
  = (1-\beta)^2 + \beta^2 + 2\beta(1-\beta)\bigl(\nabla\levt\cdot\nabla\lev_{t+1}\bigr)
  \label{eq:interp-grad}
\end{equation}
이다. 두 경사 사이의 각을 $\omega$라 하면
\begin{equation}
  \norm{\nabla\lev_\tau}^2 = 1 - 2\beta(1-\beta)\bigl(1-\cos\omega\bigr)
  \label{eq:interp-eikonal}
\end{equation}
이고, 따라서 $0<\beta<1$에서 $\norm{\nabla\lev_\tau}=1$은 오직 $\cos\omega=1$, 즉 두
경사가 정렬될 때에만 성립한다. 그 밖의 경우 $\norm{\nabla\lev_\tau}<1$이므로
$\lev_\tau$는 부호 거리 함수가 아니다.
\end{proposition}

\begin{proof}
경사는 선형 연산자이므로
$\nabla\lev_\tau = (1-\beta)\nabla\levt + \beta\nabla\lev_{t+1}$이다. 제곱 노름을
전개하고 $\norm{\nabla\levt}=\norm{\nabla\lev_{t+1}}=1$을 대입하면 식
\eqref{eq:interp-grad}이다. $\nabla\levt\cdot\nabla\lev_{t+1}=\cos\omega$를 넣고
$(1-\beta)^2+\beta^2 = 1-2\beta(1-\beta)$를 쓰면
\[
  \norm{\nabla\lev_\tau}^2
  = 1 - 2\beta(1-\beta) + 2\beta(1-\beta)\cos\omega
  = 1 - 2\beta(1-\beta)(1-\cos\omega).
\]
$0<\beta<1$이면 $2\beta(1-\beta)>0$이고 $1-\cos\omega\ge0$이므로
$\norm{\nabla\lev_\tau}^2\le1$이며, 등호는 $\cos\omega=1$일 때에만 성립한다.
\end{proof}

\begin{proposition}[보간이 법선 정사영에 주입하는 오차]
\label{prop:interp-projection}
보간된 $\lev_\tau$ 위에서 법선 정사영을 수행하면, 명제 \ref{prop:newton-steps}의 1--2회
수렴을 뒷받침하던 eikonal 가정 $\norm{\nabla\lev_\tau}\approx1$이 명제
\ref{prop:interp-eikonal}에 의해 깨지므로 추가 오차가 주입된다. 구체적으로 반복
\eqref{eq:projection-iter}의 분모에 나타나는 $\norm{\gradh\lev_\tau}^2$이 1에서
$2\beta(1-\beta)(1-\cos\omega)$만큼 벗어나며, 정사영 한 단계에 상대 오차
\begin{equation}
  \frac{\bigl\lvert\norm{\nabla\lev_\tau}^2-1\bigr\rvert}{\norm{\nabla\lev_\tau}^2}
  = \frac{2\beta(1-\beta)(1-\cos\omega)}{1-2\beta(1-\beta)(1-\cos\omega)}
  \label{eq:interp-projection-error}
\end{equation}
를 남긴다. 두 경사가 정렬($\omega=0$)되면 이 오차는 사라지고, $\beta=1/2$과 큰 각
$\omega$에서 최대가 된다.
\end{proposition}

\begin{proof}
정사영의 한 단계는 잔차를 $\gradh\lev_\tau/\norm{\gradh\lev_\tau}^2$ 방향으로
보정한다. $\lev_\tau$가 부호 거리 함수이면 분모가 1이어서 표준 Newton 단계로
축약되지만, 명제 \ref{prop:interp-eikonal}에서
$\norm{\nabla\lev_\tau}^2 = 1-2\beta(1-\beta)(1-\cos\omega)$이므로 분모가 1에서 그만큼
벗어난다. 보조정리 \ref{lem:quadratic}의 비고가 $\varepsilon$ 정규화에서 분모가
$\norm{g}^2$에서 $\norm{g}^2+\varepsilon$으로 바뀔 때 상대 편차의 일차 잔차를
남긴다고 보인 것과 동일한 논법으로, 여기서는 분모의 참값 1 대비 편차가
$2\beta(1-\beta)(1-\cos\omega)$이므로 상대 오차는 식
\eqref{eq:interp-projection-error}이다.
\end{proof}

\caveat{%
명제 \ref{prop:interp-eikonal}과 \ref{prop:interp-projection}은 선형 보간된
$\lev_\tau$가 참 중간 표면의 부호 거리 함수가 아니며 그 위에서의 법선 정사영이
체계적 오차를 가짐을 보인다. 두 프레임 사이 심장 표면의 참 운동은 일반적으로
강체나 등척이 아니므로 두 경사가 정렬된다는 보장이 없고, 따라서 식
\eqref{eq:interp-eikonal}의 편차가 실제로 발생한다. 그러므로 중간 시점 보간은
\textbf{표시 전용} 도구이며, 보간된 표면을 참 중간 해부 구조의 임상적 복원으로
해석해서는 안 된다.}

\section{프레임 간 오차 전파}
\label{sec:propagation}

\begin{proposition}[위치 잔차는 누적되지 않고 방향 오차는 표류할 수 있다]
\label{prop:propagation}
정준 surfel의 동일성은 프레임을 넘어 유지되지만, 각 프레임에서 anchor는 그 프레임의
표면 $\surft$ 위로 새로 정사영된다. 따라서 위치 잔차 $\abs{\levt(\anchor^t_i)}$은 매
프레임 보조정리 \ref{lem:quadratic}의 이차 한계로 다시 억제되어 프레임 수 $T$에
무관하게 유계이고 \emph{누적되지 않는다}. 반면 접평면 프레임의 방향은 최소 회전
전송으로 이전 프레임에서 이어받으므로, 각 프레임의 전송 오차가 더해져
\begin{equation}
  \Theta_T \le \sum_{t=1}^{T}\theta_t
  \label{eq:drift}
\end{equation}
의 방향 오차 표류를 일으킬 수 있다. 여기서 $\theta_t$는 프레임 $t$의 전송 각오차이다.
\end{proposition}

\begin{proof}
위치에 대해서는, 각 프레임에서 반복 정사영이 잔차를 보조정리 \ref{lem:quadratic}의
한계로 억제하며, 이 억제는 이전 프레임의 결과와 무관하게 그 프레임의 표면만으로
결정된다. 곧 정사영은 매 프레임 위치 오차를 \emph{재설정}하므로 누적이 없고 상한은
$T$에 무관하다. 방향에 대해서는, 전송이 $e^t_{i,1}$을 $e^{t-1}_{i,1}$로부터 계산하므로
프레임 $t$의 오차 $\theta_t$가 이전 프레임의 방향 위에 더해진다. 합성 회전의 각은
삼각부등식으로 각의 합을 넘지 않으므로 누적 방향 오차는 식 \eqref{eq:drift}로
상한지어지며, 이는 위치와 달리 재설정되지 않고 표류할 수 있다.
\end{proof}

명제 \ref{prop:propagation}은 두 가지 실용적 함의를 가진다. 첫째, 위치 정확도는 긴
시퀀스에서도 안전하다. 둘째, 방향은 그렇지 않으므로 등방 원반(명제 \ref{prop:gauge})이
긴 시퀀스에서 구조적 이점을 가진다. 면내 방향 표류가 렌더링에 무해하기 때문이다.

\section{오차 결과와 검증 지표의 대응}
\label{sec:error-to-metric}

표 \ref{tab:error-metric}은 본 장의 각 결과를 제 \ref{ch:protocol}장의 측정 지표에
대응시킨다. 이 대응이 오차 분석을 검증 가능하게 만드는 연결고리이다.

\begin{table}[htbp]
\centering
\small
\caption{오차 분석 결과와 검증 지표의 대응. ``측정 상태'' 열은 본 논문 작성 시점의
실제 상태이다.}
\label{tab:error-metric}
\begin{tabular}{@{}p{4.4cm}p{4.6cm}p{3.4cm}@{}}
\toprule
이론 결과 & 대응 지표 & 측정 상태 \\
\midrule
명제 \ref{prop:normal-error}: $\sin\theta=O(\hmax^2)$
  & normal angular error, $\spacing$ 스윕에서의 기울기
  & \NM \\
보조정리 \ref{lem:quadratic}: 이차 잔차 수축
  & $\Esurf$ 식 \eqref{eq:e-surf}, 반복별 잔차 이력
  & \NM \\
보조정리 \ref{lem:transport-orthonormal}: 정규직교성
  & 프레임 직교성 잔차
  & \NM \\
명제 \ref{prop:gauge}: 등방 gauge 불변
  & 면내 회전 전후 렌더링 차이
  & \NM \\
명제 \ref{prop:degeneracy}: $\norm{\bar e}=\abs{\sin\psi}$
  & 전송 퇴화 빈도, $\Eflicker$
  & \NM \\
명제 \ref{prop:ray-plane}: 조건수 $1/c_0$
  & depth RMSE, 컷오프 $c_0$ 스윕
  & \NM \\
명제 \ref{prop:planar-disk}: $O(\kappa s^2)$
  & silhouette IoU, $\Esurf$
  & \NM \\
명제 \ref{prop:perspective}: 아핀 오차 $\propto \Delta z/z$
  & 2DGS 대 thin-3DGS depth RMSE (RQ6)
  & \NM \\
명제 \ref{prop:spd}: SPD와 식 \eqref{eq:cond-bound}
  & CG 반복 횟수, 조건수 상계
  & \NM \\
명제 \ref{prop:lowrank}: 특잇값 꼬리
  & 저계수 복원 오차, 스펙트럼
  & \textbf{검증됨}\footnotemark[1] \\
명제 \ref{prop:interp-eikonal}: 식 \eqref{eq:interp-eikonal}
  & 측정 eikonal 노름 대 예측값
  & \NM \\
명제 \ref{prop:propagation}: 식 \eqref{eq:drift}
  & 프레임별 $\Esurf$, 누적 회전각
  & \NM \\
\bottomrule
\end{tabular}
\footnotetext[1]{명제 \ref{prop:lowrank}의 Eckart--Young 항등식은 제
\ref{sec:kernel-verification}절에서 PyTorch 없이 독립 검증되었다. 다만 실제 심장
잔차장의 특잇값 감쇠 속도는 여전히 미측정이다.}
\end{table}

\chapter{구현}
\label{ch:impl}

제 \ref{ch:method}, \ref{ch:error}장의 이론을 PyTorch로 구현하였다. 본 장은 그 과정에서
필요했던 설계 결정을 기술한다. 이들은 원 제안서에 없던 것이며, 이론을 실제로 실행
가능한 시스템으로 바꾸는 데 필수적이었으므로 그 자체로 보고할 가치가 있다.

\section{구성}

\begin{table}[htbp]
\centering
\small
\caption{모듈 구성과 대응하는 이론.}
\label{tab:modules}
\begin{tabular}{@{}llp{6.6cm}@{}}
\toprule
모듈 & 대응 & 내용 \\
\midrule
\texttt{core/}        & --- & 격자 기하, 설정, 결정성, 단계별 계측 \\
\texttt{levelset/}    & 제 \ref{sec:spacing-aware}--\ref{sec:warmstart}절
                      & 간격 인지 연산자, Chan--Vese, 부호 거리 도구, marching tetrahedra \\
\texttt{surfel/}      & 제 \ref{sec:canonical}--\ref{sec:transport}절
                      & 정준 집합, 법선 정사영, 접선 전송, 밀도 제어 \\
\texttt{render/}      & 제 \ref{sec:raster}절
                      & 2DGS 래스터라이저, thin-3DGS 및 mesh 기준선, 광선 행진 참조 \\
\texttt{losses.py}    & 식 \eqref{eq:total-loss} & 외형/mask/normal/distortion 항 \\
\texttt{residual/}    & 제 \ref{sec:residual}절 & 정칙화 최소제곱(CG), 저계수 압축 \\
\texttt{data/}        & 제 \ref{sec:phantom}절 & 합성 팬텀, ACDC/M\&Ms-2 로더 \\
\texttt{metrics/}     & 제 \ref{sec:metrics}절 & 모든 보고 지표 \\
\texttt{pipeline/}    & 제 \ref{sec:pipeline}절 & 전처리, 정준 fitting, 재생, 직렬화 \\
\texttt{experiments/} & 제 \ref{ch:protocol}장 & 이론 검증, RQ1--RQ6, ablation \\
\bottomrule
\end{tabular}
\end{table}

전체 49개 모듈, 약 13{,}900행이다.

\section{설계 결정}
\label{sec:impl-decisions}

\subsection{기하는 최적화 불가능하다}

surfel의 anchor, 접선축, 법선을 \texttt{nn.Module}의 \emph{버퍼}로 등록하고 파라미터로
등록하지 않았다. 따라서 어떤 손실도 이들에 도달할 수 없다. 최적화되는 것은 진폭,
불투명도, 두 접선 scale뿐이다.

이는 편의가 아니라 주장의 정합성 문제이다. 제 \ref{sec:raster}절에서 언급하였듯
기하가 학습 가능하면 본 방법은 조용히 자유형 동적 2DGS fitting이 되고, ``기하가
$\surft$에서만 온다''는 주장이 무효가 된다. 학습률을 0으로 두는 것으로도 같은 효과를
낼 수 있으나, 설정 실수 한 번으로 주장이 깨진다. 버퍼로 등록하면 구조적으로 불가능해진다.
구현의 테스트는 \texttt{anchor.requires\_grad}가 참이면 실패한다.

\subsection{감독 신호가 무엇인지 정의해야 했다}

원 제안서는 감독을 ``알려진 slice plane의 mask, depth, normal, intensity''로 제한한다.
그러나 cine CMR은 복셀 볼륨이고 사진이 아니므로, 이 문장을 실행 가능한 목표로
바꾸어야 했다.

채택한 정의는 다음이다. 알려진 slice plane 카메라의 각 화소 광선에 대해 저장된
$\surft$와의 첫 교차를 광선 행진으로 찾고, 그 지점에서
\begin{itemize}[leftmargin=2em]
  \item 교차 여부 $\to$ 참조 silhouette, $\mathcal{L}_{\mathrm{mask}}$의 목표,
  \item 교차까지의 거리 $\to$ 참조 $\depthq$, depth RMSE의 기준,
  \item $\gradh\levt/\norm{\gradh\levt}$ $\to$ 참조 법선, 명제
  \ref{prop:normal-error}의 각오차 기준,
  \item 그 지점의 MRI 값 $\to$ $\mathcal{L}_{\mathrm{app}}$의 목표
\end{itemize}
를 읽는다. 이 정의의 두 성질이 중요하다. 첫째, 네 신호 모두 level-set과 MRI에서만
오므로 fitting이 저장 표현에 없는 정보를 받지 않는다. 둘째, mesh, thin-3DGS, 2DGS가
\emph{동일한} 참조에 대해 채점되므로 RQ5/RQ6가 표현을 비교하며 분할을 비교하지 않는다.

\subsection{좁은 띠는 희소 목록이 아니라 조밀 crop이다}

교과서적 좁은 띠는 활성 복셀의 희소 목록을 유지한다. GPU에서는 이것이 작은 상자에
대한 조밀 연산보다 느린 경우가 많다. 따라서 띠를
$\{\abs{\lev}<b\}$의 경계 상자에 유한차분 halo를 더한 \emph{조밀 crop}과, 그 안에서의
masked 갱신으로 실현하였다. 속도 이득은 crop에서 나온다.

한 가지 함정이 있다. 식 \eqref{eq:cin}, \eqref{eq:cout}의 영역 평균은 $\dom$ 전체에
대한 적분이다. crop 위에서 계산하면 외부 영역이 얇은 껍질로 축소되어 $\cout$이
편향된다. 따라서 영역 평균만은 \emph{전체} 볼륨에서 다시 계산한다. 이 비용은 두 번의
축약이므로 곡률 계산에 비해 무시할 만하다.

\subsection{빠른 래스터라이저의 두 근사}

성능을 위한 타일 기반 래스터라이저는 식 \eqref{eq:compositing}에 두 가지 근사를
도입한다.

\begin{enumerate}[leftmargin=2em]
  \item 타일 내 정렬은 화소별 정확한 $\depthq^t_i(q)$가 아니라 anchor 깊이로 한다.
  \item 타일당 가장 가까운 $M$개만 유지하고 나머지를 버린다.
\end{enumerate}

이 두 근사가 무해하다고 \emph{가정하지 않기 위해}, 모든 surfel을 모든 화소에 대해
평가하고 화소별로 정확한 $\depthq$로 정렬하는 브루트포스 참조 래스터라이저를 함께
구현하였다. 작은 영상에서 두 구현의 차이를 측정할 수 있으며, 구현의 테스트는 잘
분리된 surfel 배치에서 두 결과가 $10^{-3}$ 이내로 일치할 것을 요구한다. 절단된 페어
수와 절단율도 매 렌더링에서 보고한다.

\subsection{Depth distortion의 대체}

원 2DGS의 depth distortion 항은 $\sum_i\sum_j w_i w_j\abs{\depthq_i-\depthq_j}$로
$O(M^2)$이다. 구현은 $O(M)$인 가중 분산
\begin{equation}
  \sum_i w_i\bigl(\depthq_i - \bar\depthq\bigr)^2,
  \qquad \bar\depthq = \frac{\sum_i w_i \depthq_i}{\sum_i w_i}
  \label{eq:distortion-surrogate}
\end{equation}
로 대체하였다. 두 형태 모두 한 광선의 모든 기여가 한 깊이에 모일 때 정확히 소멸하므로
정규화의 목적은 보존된다. 이는 참조 구현으로부터의 의도적 이탈이므로 명시한다.

\subsection{임의 시점 탐색 문제}
\label{sec:seek-problem}

식 \eqref{eq:store}은 프레임별 anchor를 저장하지 않는다. 따라서 저장된 것만으로
프레임 $t$를 표시하려면 \emph{정준} anchor를 $\surft$로 직접 정사영해야 한다. 그러나
전처리는 anchor를 $\anchor^{t-1}\to\anchor^t$로 \emph{순차} 정사영하였다. 두 경로는
동일하지 않다. 순차 정사영은 경로를 누적하지만 직접 정사영은 훨씬 긴 한 걸음을 밟으므로
trust region에 걸리거나 잘못된 표면 분기에 착지할 가능성이 크다.

어느 한쪽을 조용히 택하는 대신 구현은 두 방식을 모두 제공하고 그 격차를 측정하는
함수를 제공한다. 직접 정사영이 $\Esurf$를 유의하게 악화시킨다면, 정직한 저장량 회계는
프레임별 anchor를 추가하거나 viewer를 순차 재생으로 제한해야 한다. 이는 표현에 대한
\emph{발견}이며 버그가 아니다.

\caveat{%
이 문제는 원 제안서에서 다루어지지 않았다. 구현 과정에서 드러난 것이며, 식
\eqref{eq:store}의 저장량 주장과 임의 탐색 기능이 양립하는지에 직접 영향을 준다.
RQ4가 이를 측정한다.}

\section{합성 팬텀}
\label{sec:phantom}

명제 \ref{prop:normal-error}, 보조정리 \ref{lem:quadratic}, 명제
\ref{prop:planar-disk}, \ref{prop:interp-eikonal}은 모두 \emph{수렴률 또는 정확한 항등식}에
대한 주장이다. 이를 확인하려면 정확한 참조가 필요하다. 즉 정확한 부호 거리 함수,
정확한 표면 법선, 그리고 $\spacing$을 임의로 바꿀 수 있는 능력이다. ACDC와 M\&Ms-2는
ED/ES label만 제공하며 어느 것도 제공하지 않는다.

따라서 해석적 참값을 갖는 합성 4D 팬텀을 구현하였다. 좌심실 내막을 수축하는 장축
회전 타원체로 모델링하고, 강체 변환이 거리를 보존하므로 기울임 아래에서도 부호 거리
함수가 \emph{정확}하게 유지된다. 비등방 표본화($h_z\gg h_x=h_y$), Rician 잡음, 매끄러운
명암 불균일, 그리고 혈액 풀 안의 심근 밝기 papillary muscle 덩어리를 포함한다.
papillary muscle은 Chan--Vese의 조각마다 상수 가정을 \emph{의도적으로} 위반하며, 정답
표면에는 포함되지 않는다(ACDC 관례와 일치). 또한 혈액과 심근 밝기가 주기에 따라
맥동하도록 하였다. 그렇지 않으면 외형 잔차가 항등적으로 0이 되어 RQ3가 공허해진다.

\caveat{%
팬텀은 \textbf{검증 도구이며 데이터셋이 아니다}. 실제 해부학, scanner 변이, 질환에
대한 강건성에 대해 아무것도 말해 주지 않는다. 실제 데이터 경로는 별도로 구현되어
있으나 실행되지 않았다.}

\section{검증 상태}
\label{sec:impl-status}

\subsection{실행되지 않았다}

본 구현은 PyTorch를 설치할 수 없는 환경에서 작성되었다(패키지 색인 접근 불가, 캐시된
wheel 없음, GPU 없음). 따라서 \textbf{PyTorch를 필요로 하는 어떤 부분도 실행된 적이
없다}. 확인된 것은 다음이다.

\begin{itemize}[leftmargin=2em]
  \item 49개 모듈과 4개 테스트 파일의 문법 (\texttt{compileall} 통과),
  \item 모든 상대 import가 정의된 이름으로 해소됨,
  \item \texttt{\_\_all\_\_} 항목과 모든 테스트 import의 해소,
  \item 텐서 형상과 수식의 행별 검토. 이 과정에서 실제 결함 4건을 발견하고 수정하였다.
  \footnote{(i) 0/1 텐서에 대한 \texttt{argmax}가 첫 최댓값을 반환한다는 보장이 없어
  광선 행진의 첫 교차 탐색이 비결정적이었던 문제, (ii) 밀도 제어로 surfel 수가 변할 때
  잔차 배열의 형상이 어긋나던 문제, (iii) 작은 격자에서 팬텀 타원체가 볼륨 경계에서
  잘려 닫힌 표면 가정이 깨지던 문제, (iv) 정사영 카메라에 원근 보정 보간을 잘못
  적용하던 문제.}
\end{itemize}

\subsection{실행된 검증: 이산 논리 커널}
\label{sec:kernel-verification}

PyTorch를 필요로 하지 않는 부분이 있다. 이산 색인 연산과 참조표이다. 이 영역은
오프바이원 하나가 예외를 발생시키지 않고 결과를 조용히 오염시키므로 위험도가 높다.
이를 표준 라이브러리만으로 재구현하여 독립 참조와 대조하였다. 41개 항목이
\emph{실제로 실행되어} 통과하였다.

\begin{table}[htbp]
\centering
\small
\caption{실행된 커널 검증. 본 논문에서 실제로 실행된 유일한 검증이다.}
\label{tab:kernel-checks}
\begin{tabular}{@{}p{6.2cm}p{7.2cm}@{}}
\toprule
검사 항목 & 독립 참조 \\
\midrule
큐브의 6면체 분해가 정확히 타일링 & 부피 합, 점 포함 횟수 \\
marching tetrahedra 16개 case의 위상 정확성 & 교차 간선 집합과의 일치 \\
(surfel, 타일) 페어 열거 & 브루트포스 이중 루프 \\
타일별 깊이 정렬과 절단 & 최근접 $M$개 보존 \\
점유 비트마스크 팩/언팩 & 8의 모든 나머지 길이에서 왕복 \\
타원체 부호 거리 이분법 브래킷 & 해석적 구, 브루트포스 최근접점 \\
삼선형 가중치의 합과 순서 & 아핀 정확성, 코너 대응 \\
공액경사법 수렴 & $n$회 내 정확해 \\
Godunov 상류 재초기화 스텐실 & $\norm{\nabla\lev}=1$이 고정점 \\
저장량 손익분기 대수 & 식 \eqref{eq:breakeven} 역산 시 $\CR=1$ \\
수축 프로파일의 연속성 & ES 및 주기 경계에서의 도약 \\
\bottomrule
\end{tabular}
\end{table}

\begin{remark}[검증이 검증 자체의 결함을 찾은 사례]
첫 실행에서 ``모든 내부 점이 정확히 하나의 사면체에 속한다''는 항목이 216개 표본 중
96개에서 실패하였다. 6개 사면체가 공유하는 내부 평면은 $x=y$, $y=z$, $x=z$ 세 개이며,
격자점 $(i,j,k)$, $i,j,k\in\{1,\dots,6\}$ 중 이 평면 위에 놓인 개수를 포함--배제로
계산하면
\[
  3\cdot36 - 3\cdot6 + 6 = 96
\]
으로 \emph{정확히 일치}하였다. 즉 분해는 옳고 검사가 사면체 경계면 위의 점을
표본화한 것이 원인이었다. 표본을 무리수 오프셋으로 옮기고, 경계면 위의 점이 정확히
두 사면체에 속하는지를 별도 항목으로 추가하여 해결하였다. 오류 개수가 예측값과
정확히 맞았다는 사실 자체가 분해의 정확성에 대한 강한 증거였다.
\end{remark}

\subsection{단계별 스모크 진단}

실행되지 않은 구현의 첫 실행에서는 여러 결함이 동시에 드러날 것이 예상된다. 첫 실패에서
멈추는 테스트 도구는 ``한 번 실행에 결함 하나'' 루프를 강요한다. 따라서 21개 독립
스테이지로 구성된 스모크 진단을 함께 구현하였다. 각 스테이지가 자신의 예외를 잡으므로
한 번 실행이 모든 실패를 보고하고, 의존성 순서로 배치되어 첫 실패가 이후 실패를 설명한다.
선행 스테이지가 실패한 스테이지는 \texttt{SKIP}으로 표시하여 무관한 실패와 구분한다.

스테이지는 실행되지 않은 구현에서 가장 발생 가능성이 높은 결함을 겨냥한다. 부호 규약
역전, 프레임 직교성 파괴, 수반 연산자 오류, 재생이 정준 집합을 변형하는 문제 등이며,
오류 메시지가 추정 원인을 명시한다. 예를 들어 구 곡률이 $-2/R$ 대신 $+2/R$로 측정되면
``내부/외부 규약이 역전되었음을 뜻한다''를 함께 출력한다.

\section{재현성}

전역 난수 상태에 의존하지 않도록 팬텀 표본 추출과 정준 초기화에 생성기를 명시적으로
전달한다. 면적 균일 표본 추출은 CPU에서 수행한 뒤 색인만 장치로 옮기므로 계산 장치와
무관하게 비트 단위로 재현된다. 모든 결과 파일은 장치, PyTorch 버전, GPU 이름을 포함한
환경 기록을 함께 저장한다. 타이밍 주장은 하드웨어 의존적이므로 이 기록 없이는 의미가
없다.

\chapter{실험 설계}
\label{ch:protocol}

본 장은 무엇을 어떻게 측정할지를 사전에 확정한다. 결과가 없는 상태에서 설계를
확정해 두는 것은 사후 선택 편향을 막기 위함이다.

\section{연구 질문}
\label{sec:rqs}

\begin{description}[leftmargin=1.6cm,style=nextline]
  \item[RQ1] 이전 level-set seed는 각 프레임의 Chan--Vese 반복 횟수와 시간을
  줄이는가? (정리 \ref{thm:warmstart})
  \item[RQ2] 법선 정사영은 프레임별 2DGS 최적화보다 적은 비용으로 표면과 렌더링
  품질을 유지하는가? (식 \eqref{eq:cost-hypothesis})
  \item[RQ3] 정준 2DGS, 프레임별 표면, 진폭 잔차 표현은 전체 2DGS 저장보다 작은가?
  (식 \eqref{eq:cr}, \eqref{eq:breakeven})
  \item[RQ4] 재생 시 정사영, orientation 갱신, 래스터화를 포함하여 30 FPS 이상을
  달성하는가? (식 \eqref{eq:frame-budget})
  \item[RQ5] 같은 Chan--Vese 표면에서 2DGS가 mesh보다 경계, depth/normal, appearance
  또는 시간적 부드러움을 개선하는가?
  \item[RQ6] 2D 원반은 같은 수의 얇은 3D Gaussian보다 표면 artifact와 시점 비일관성을
  줄이는가? (명제 \ref{prop:perspective})
\end{description}

\section{기준선}
\label{sec:baselines}

모든 기준선을 \emph{동일한 파이프라인의 설정}으로 표현하였다. 별도 구현이라면 측정된
차이가 아이디어가 아니라 엔지니어링에서 올 수 있다. Chan--Vese 솔버, 감독 뷰, 손실,
최적화기, 타일링, 컷오프, 지표가 모두 공유 코드이며 차이는 명시된 항목뿐이다.

\begin{table}[htbp]
\centering
\small
\caption{기준선과 각각이 분리하는 질문.}
\label{tab:baselines}
\begin{tabular}{@{}lp{9.0cm}@{}}
\toprule
기준선 & 분리하는 것 \\
\midrule
\texttt{mesh-only} & 동일 표면에서 부드러운 Gaussian silhouette이 텍스처 mesh보다
  나은가 (RQ5) \\
\texttt{thin-3dgs} & primitive 자체: 2D 원반 대 얇은 3D 타원체 (RQ6) \\
\texttt{independent-2dgs} & 품질 상한과 비용 하한: 매 프레임 처음부터 재최적화 (RQ2, RQ3) \\
\texttt{param-copy} & anchor가 움직여야 하는가, 외형 보정만으로 충분한가 \\
\texttt{closest-point} & 식 \eqref{eq:least-norm}의 법선 제약을 버리면 어떻게 되는가 \\
\texttt{normal-no-residual} & 외형 잔차의 기여 \\
\texttt{cold-start-cv} & 식 \eqref{eq:warmstart-init}의 웜스타트 효과 (RQ1) \\
\texttt{cv-dyn2dgs} & 전체 방법 \\
\bottomrule
\end{tabular}
\end{table}

mesh 기준선은 의도적으로 강하게 구성하였다. level-set 경사에서 얻은 정점 법선,
MRI에서 표본화한 정점 진폭, 원근 정합 보간, z-buffer에 의한 정확한 가시성을 사용한다.
mesh가 할 수 없는 것은 부드러운 silhouette뿐이며 이것이 시험 대상이다. 다만 mesh는
화소 단위 이진 덮음을 가지므로 silhouette IoU를 0.5에 고정하지 않고 여러 문턱값에
대해 스윕한다.

\section{외부 비교군: 무엇이 빠져 있는지}
\label{sec:external-baselines}

앞 절의 기준선은 모두 \emph{본 파이프라인의 설정}이다. 이는 귀속에는 최선이지만
외부 타당성에는 최악이다. 발표된 방법과의 비교가 하나도 없기 때문이다. 본 절은 그
공백을 은폐하지 않고 목록화한다. \textbf{여기 열거된 비교는 어느 것도 실시되지
않았다.}

주장은 서로 다른 세 층위에 있고, 층위를 섞으면 불공정하거나 무의미한 비교가 된다.
층위마다 고정해야 할 것이 다르다.

\begin{table}[htbp]
\centering\small
\caption{실시되지 않은 외부 비교군. ``고정''은 비교가 성립하기 위해 동일해야 하는
것을, ``축''은 그 비교군이 채점될 수 있는 지표를 뜻한다.}
\label{tab:external-baselines}
\begin{tabular}{@{}llp{3.4cm}p{2.9cm}@{}}
\toprule
층위 & 비교군 & 고정 & 축 \\
\midrule
\multirow{3}{*}{1. 표면 공급원}
 & nnU-Net mask $\to$ SDF~\citep{isensee2021_nnunet} & 표현, 렌더러, 지표 & Dice/HD95, $\Esurf$ \\
 & CSTM 4D 분할~\citep{ye2025_cstm}     & 동일 & 동일 \\
 & 정답 mask $\to$ SDF (oracle)         & 동일 & 분할/표현 오차 분리 \\
\midrule
\multirow{4}{*}{2. 표현 primitive}
 & Gaussian surfels~\citep{dai2024_gaussian_surfels} & \textbf{표면}, 감독 뷰, 손실 & PSNR/SSIM, depth, normal \\
 & 3DGS (기하 자유)~\citep{kerbl2023_3dgs}           & 감독 뷰, 지표 & 품질 대 $\Esurf$ 상충 \\
 & DG-Mesh~\citep{liu2024_dgmesh}                    & 감독 뷰 & mesh 품질, 대응 \\
 & biv-me mesh~\citep{dillon2026_bivme}              & 데이터 & 임상 지표, 저장량 \\
\midrule
\multirow{4}{*}{3. 시간 모델}
 & Dyna3DGR~\citep{fu2025_dyna3dgr}   & 데이터(ACDC) & 추적, 전처리 시간 \\
 & AT-GS~\citep{chen2025_atgs}        & 데이터, 뷰   & $\Eflicker$, 프레임당 시간 \\
 & D-2DGS~\citep{zhang2025_d2dgs}, ST-2DGS~\citep{wang2024_st2dgs} & 동일 & 표면 정확도 \\
 & Deformable 3DGS~\citep{yang2023_deformable3dgs}, 4DGS~\citep{wu2024_4dgs} & 동일 & 품질 대 비용 \\
\midrule
\multirow{3}{*}{4. 저장/재생}
 & 사전렌더링 비디오 (H.265/AV1) & 시점 집합, 화질 & 바이트, FPS, 시점 수 \\
 & Spacetime Gaussians~\citep{li2024_stg} & 데이터 & 바이트, FPS \\
 & Cinematic Anatomy 3DGS~\citep{niedermayr2024_cinematic} & --- & 바이트, FPS \\
\bottomrule
\end{tabular}
\end{table}

\caveat{%
\textbf{층위 1이 가장 시급하다.} Chan--Vese의 Dice가 nnU-Net 계열보다 낮을 것은
거의 확실하다. 이에 대한 방어는 ``우리 Dice도 좋다''가 아니라 ``본 방법은 표면
공급원에 무관하다''여야 하며, 그것은 공급원을 교체해 보아야만 알 수 있다. 결과가
어느 쪽이든 정보가 된다. 품질이 올라가면 모듈성이 입증되고, 오르지 않으면 병목이
표면이 아니라 표현이라는 뜻이며, mask에서 변환한 SDF로는 정사영이 망가진다면 그것은
\emph{진짜 부호 거리장이 필요하다}는 측정된 근거가 되어 오히려 Chan--Vese를 쓸
이유가 된다. oracle 표면은 분할 오차와 표현 오차를 분리하는 유일한 수단이므로 반드시
포함해야 한다.}

\caveat{%
\textbf{AT-GS~\citep{chen2025_atgs}는 방법 공간에서 가장 가까운 이웃이며, 원 제안서가
지목하지 못한 비교군이다.} 정준 Gaussian surfel 집합, 직전 프레임에서의 초기화,
적응적 densification, 그리고 곡률 map의 시간 일관성 항으로 흔들림을 억제한다는 점까지
본 연구와 기제가 겹친다. 결정적 차이는 단 하나로, AT-GS는 기하를 \emph{학습}하고 본
연구는 기하를 변분적 표면에 \emph{고정}한다. 프레임당 학습 시간을 보고하므로 본 연구의
전처리 비용과 직접 대응한다. 이 비교의 부재는 심사에서 지적될 수
있는 공백이다.}

\caveat{%
\textbf{층위 4는 본 논문의 저장량 주장을 위협한다.} 첫째, 시점을 몇 개로 고정한다면
사전렌더링한 비디오 코덱이 어떤 3D 표현보다도 작고 빠르다. 따라서 식
\eqref{eq:cr}의 주장은 \emph{자유 시점}이 필요할 때로 한정되며, 이
경계는 주장이 아니라 손익분기 시점 수로 \emph{측정}되어야 한다. 둘째, 압축된 4DGS
계열은 동적 장면 전체를 수 MB로 보고하고 있고, 의료 영역에서도 Cinematic Anatomy
3DGS가 수 GB 볼륨을 70\,MB 이하로 줄여 60\,FPS로 렌더링한다고
보고한다~\citep{niedermayr2024_cinematic}. 이들은 정적 볼륨 또는 일반 장면이라
과제가 다르지만, ``표면만 저장하므로 작다''는 주장을 이들과 대조하지 않고 제시하는
것은 불충분하다.}

다음 비교는 \textbf{하지 않는 것이 옳다}. (i) 본 연구의 Dice를 발표된 ACDC 딥
네트워크 수치와 같은 표에 넣는 것---과제가 다르다(본 연구는 SDF가 필요하고 그들은
mask만 낸다). (ii) 볼륨을 렌더링하는 방법과 PSNR을 비교하는 것---본 연구는 표면을
렌더링한다. (iii) Dyna3DGR의 발표 시간과 본 연구의 전처리 시간을 직접 비교하는
것---하드웨어, 데이터, 목적이 다르다.

\subsection{구현된 것과 고정된 것}
\label{sec:comparison-impl}

표 \ref{tab:external-baselines}의 비교 중 세 층위는 외부 코드 없이 본 파이프라인
안에서 수행할 수 있으며, 그 부분은 구현되어 있다.

\begin{itemize}[leftmargin=2em]
  \item \textbf{층위 1}: \texttt{levelset/surface\_source.py}. Chan--Vese, 외부 mask,
  oracle을 같은 인터페이스로 교체한다. 부분 관측 문제를 형 수준에서 강제한다---ACDC의
  oracle은 ED/ES만 덮으므로 \texttt{require\_full()}이 전 주기 지표를 거부한다.
  \texttt{sdf\_fidelity()}는 eikonal 위반과 \emph{staircase index}를 함께 보고하여,
  정사영 저하가 관측될 때 그 원인이 영등위면의 voxel 양자화인지를 지목한다.
  \item \textbf{층위 2}: \texttt{surfel/free3dgs.py}. 위치, 방향, 세 축 scale 전부를
  최적화하는 3D Gaussian이다. \texttt{SurfelSet2D}와 정확히 거울상이다---거기서는
  기하가 buffer이고 여기서는 기하가 parameter이므로, 설정으로 서로 바꿀 수 없다.
  초기화 방식이 결과를 바꾸므로 기본값을 두지 않고 이름에 포함시켰다.
  \item \textbf{층위 4}: \texttt{metrics/viewpoint.py}. 손익분기 시점 수를 계산한다.
  비트레이트 기본값을 \emph{제공하지 않는다}---실측이거나 근거를 명시한 가정이어야
  하며, 경쟁자의 수치를 추측하는 것은 자기 수치를 추측하는 것보다 나쁘다.
\end{itemize}

진짜 외부 방법은 재구현하지 않는다. 누군가의 방법을 재구현해서 졌다고 보고하는 것은
가장 신뢰할 수 없는 비교이며, 모든 차이가 재구현 탓으로 귀속될 수 있다. 대신
\texttt{external/manifest.json}이 각 대상을 commit SHA로 고정하고,
\texttt{scripts/fetch\_external.py}가 필요할 때만 clone하며,
\texttt{experiments/external.py}가 그 출력을 \emph{내부 기준선과 동일한} ray-march
기준(제 \ref{sec:metrics}절)에 대해 채점한다. 외부 방법이 제공하지 않은 채널은
0으로 채우지 않고 미측정으로 남긴다. 깊이를 내보내지 않은 방법은 깊이 RMSE에서
무한히 나쁜 점수를 받는 것이 아니라 \emph{채점되지 않는다}.

%% AUTO-GENERATED by scripts/make_tables.py -- DO NOT EDIT BY HAND.
%% Source: external/manifest.json
%% Status: STUB - no measurements available
\begin{table}[htbp]
\centering\small
\caption{외부 비교군의 고정 commit과 실행 상태. commit은
\texttt{external/manifest.json}에서 자동 생성되므로 표와 고정값이 어긋날 수 없다.
$\dagger$는 라이선스상 본 저장소에 포함할 수 없는 대상(라이선스 파일이 없거나
Inria/MPII 연구 전용)이며, 고정 commit으로 필요할 때만 clone한다.}
\label{tab:external-status}
\begin{tabular}{@{}clll@{}}
\toprule
층위 & 대상 & 고정 commit & 상태 \\
\midrule
1 & \texttt{cstm}$\dagger$ & \texttt{---} & \NM \\
1 & \texttt{nnunet} & \texttt{b4f9539f} & \NM \\
2 & \texttt{2dgs}$\dagger$ & \texttt{f3e3b9fa} & \NM \\
2 & \texttt{3dgs}$\dagger$ & \texttt{54c035f7} & \NM \\
2 & \texttt{biv-me} & \texttt{3d237a5f} & \NM \\
2 & \texttt{dg-mesh} & \texttt{754f42c6} & \NM \\
2 & \texttt{gaussian-surfels}$\dagger$ & \texttt{01d12493} & \NM \\
3 & \texttt{4dgs} & \texttt{843d5ac6} & \NM \\
3 & \texttt{at-gs} & \texttt{3fe29349} & \NM \\
3 & \texttt{dyna3dgr} & \texttt{79ad3dac} & \NM \\
3 & \texttt{dynamic-2dgs} & \texttt{ab4cafd6} & \NM \\
3 & \texttt{sc-gs}$\dagger$ & \texttt{3a9d2ad4} & \NM \\
4 & \texttt{cinematic-gaussians}$\dagger$ & \texttt{0bf4edbf} & \NM \\
4 & \texttt{spacetime-gaussians} & \texttt{427abfc5} & \NM \\
4 & \texttt{video-codec} & \texttt{---} & \NM \\
\bottomrule
\end{tabular}
\end{table}


\caveat{%
외부 코드는 본 저장소에 포함하지 않는다. 대상 중 넷은 법적으로 재배포할 수 없다.
\texttt{gaussian\_surfels}와 \texttt{cinematic-gaussians}는 LICENSE 파일이 아예 없어
기본 저작권이 적용되며, Inria/MPII의 3DGS와 2DGS는 연구 전용 라이선스로 재배포가
제한된다. 실행과 평가는 가능하지만 복제 배포는 불가하다. 부수적으로 전부 포함하면
약 1.2\,GB가 늘고 전송 patch의 순수 텍스트 성질이 깨진다. 무엇보다 commit SHA 고정이
더 재현 가능하다---고정된 SHA는 어떤 수치가 정확히 어느 commit에서 나왔는지를 기록하고,
복사본은 누군가 한때 무엇을 복사했다는 사실만 기록한다.}

%% AUTO-GENERATED by scripts/make_tables.py -- DO NOT EDIT BY HAND.
%% Source: results/comparison.json
%% Status: STUB - no measurements available
\begin{table}[htbp]
\centering\small
\caption{층위 1: 표면 공급원을 교체하고 표현 \emph{과} 지표는 고정한 비교
(제 \ref{sec:external-baselines}절). \texttt{staircase}는 mask에서 유도한 부호
거리장의 영등위면 양자화 정도이며, 정사영 저하가 관측될 경우 그 원인을 지목한다.}
\label{tab:layer1-source}
\begin{tabular}{@{}lrrrrr@{}}
\toprule
표면 공급원 & Dice & $\Esurf$ (mm) & PSNR$_{\mathrm{ROI}}$ (dB)
  & eikonal & staircase \\
\midrule
Chan--Vese (제안) & \NM & \NM
  & \NM & \NM
  & \NM \\
외부 mask $\to$ SDF & \NM & \NM
  & \NM & \NM
  & \NM \\
oracle (정확한 표면) & \NM & \NM
  & \NM & \NM
  & \NM \\
\midrule
spacing-aware 해제 & \NM
  & \NM
  & \NM
  & \NM
  & \NM \\
\bottomrule
\end{tabular}
\end{table}

%% AUTO-GENERATED by scripts/make_tables.py -- DO NOT EDIT BY HAND.
%% Source: results/comparison.json
%% Status: STUB - no measurements available
\begin{table}[htbp]
\centering\small
\caption{층위 2: 기하를 표면에 고정하는 것의 대가와 이득. 광도 지표와 기하 지표를
\emph{반드시} 함께 읽어야 한다. 자유 기하는 PSNR에서 이기고 $\Esurf$에서 크게
지도록 설계상 예상되며, 그 상충이 본 논문의 논점이다.}
\label{tab:layer2-free}
\begin{tabular}{@{}lrrrrr@{}}
\toprule
표현 & PSNR$_{\mathrm{ROI}}$ (dB) & $\Esurf$ (mm) & 법선 오차 ($^\circ$)
  & 프레임당 바이트 & 기하 고정 \\
\midrule
CV-Dyn2DGS (2D 원반) & \NM
  & \NM & \NM
  & \NM & 예 \\
Gaussian surfel (z-scale $\to0$) & \NM
  & \NM
  & \NM
  & \NM & 예 \\
thin-3DGS ($\rho=0.1$) & \NM
  & \NM & \NM
  & \NM & 예 \\
자유 3DGS (표면 초기화) & \NM
  & \NM
  & \NM
  & \NM & \textbf{아니오} \\
자유 3DGS (무작위 초기화) & \NM
  & \NM
  & \NM
  & \NM & \textbf{아니오} \\
\bottomrule
\end{tabular}
\end{table}

%% AUTO-GENERATED by scripts/make_tables.py -- DO NOT EDIT BY HAND.
%% Source: results/comparison.json
%% Status: STUB - no measurements available
\begin{table}[htbp]
\centering\small
\caption{층위 4: 시점 수에 따른 저장량. 본 방법은 시점 수에 대해 \emph{평탄}하고
사전렌더링 비디오는 선형으로 증가한다. 손익분기 시점 수가 ``자유 시점이 필요할 때
이긴다''는 진술을 주장에서 측정값으로 바꾼다. 비디오 비트레이트는 실측이거나 명시된
가정이며, 기본값은 제공되지 않는다.}
\label{tab:layer4-viewpoint}
\begin{tabular}{@{}rrrrl@{}}
\toprule
시점 수 & 본 방법 (MB) & 비디오 (MB) & 원본 볼륨 (MB) & 우세 \\
\midrule
1 & \NM & \NM & \NM & \NM \\
2 & \NM & \NM & \NM & \NM \\
4 & \NM & \NM & \NM & \NM \\
8 & \NM & \NM & \NM & \NM \\
16 & \NM & \NM & \NM & \NM \\
32 & \NM & \NM & \NM & \NM \\
\midrule
손익분기 시점 수 & \multicolumn{4}{l}{\NM} \\
\bottomrule
\end{tabular}
\end{table}


\section{데이터}
\label{sec:data}

새로운 의료 데이터셋이나 수동 label을 제작하지 않는다.

\begin{itemize}[leftmargin=2em]
  \item \textbf{합성 팬텀} (제 \ref{sec:phantom}절): 수렴률 검증과 기하 정확도.
  유일하게 정확한 부호 거리 함수와 법선을 제공한다.
  \item \textbf{ACDC}~\citep{bernard2018_acdc}: 150명의 단기축 cine CMR과 ED/ES
  분할. 웜스타트 수렴시간, ED/ES에서의 Dice/HD95, 전처리시간과 모델 크기 비교,
  전 주기 FPS와 시간적 흔들림.
  \item \textbf{M\&Ms-2}~\citep{martinisla2023_mnms2}: 360명의 multi-center,
  multi-vendor cine CMR. 다른 병원, scanner, 질환에서 표면 추적과 appearance가
  유지되는지 확인. LAX 영상은 fitting에 쓰지 않고 새로운 시점에서의 일관성 검사에만
  사용한다. 필요하면 1차 M\&Ms~\citep{campello2021_mnms}를 추가 외부 도메인으로
  쓸 수 있다(좌심실과 심근 대상이며 M\&Ms-2와 다른 데이터셋이다).
  \item \textbf{STACOM 2011 Motion Tracking}~\citep{tobongomez2013_benchmarking}
  (선택적): 운동 정답을 포함하므로, 법선 형상 이동만 모델링할 때 실제 조직 추적에서
  어느 정도 정보가 손실되는지를 보이는 보조 실험에 사용한다. 이 결과는 본 방법을
  strain 추정 방법으로 제시하기 위한 것이 \emph{아니라}, 제
  \ref{sec:tangential}절의 한계를 정량화하기 위한 것이다.
\end{itemize}

\caveat{%
정답은 ED와 ES에만 존재한다. 중간 프레임은 영상 재현과 시간적 안정성 평가에만
사용하며 분할 정답이 있다고 가정하지 않는다. 구현의 로더는 중간 프레임의 mask를
\texttt{None}으로 두어 이를 형 수준에서 강제한다.}

\subsection{데이터의 출처와 ``올바른 사본''}
\label{sec:data-provenance}

위 데이터셋은 모두 개별 등록과 데이터 사용 동의가 필요하므로 자동으로 받을 수 없다.
\texttt{external/datasets.json}이 각 데이터셋의 출처, 라이선스, 인용 의무, 어느 비교에
쓰이는지, 그리고 \texttt{cvdyn2dgs/data/real.py}가 실제로 glob하는 디스크 배치를
기록한다. \texttt{scripts/fetch\_datasets.py}는 게이트된 데이터셋을 우회하지 않고
무엇을 해야 하는지만 안내한 뒤, 수동으로 받은 사본을 \emph{검증}한다.

검증의 핵심은 파일 존재 여부가 아니라 \textbf{voxel 기하}다.

\caveat{%
\textbf{리샘플링된 사본은 기여 1을 조용히 무효화한다.} 본 논문의 첫 기여는 spacing을
반영한 이산화이며, 그 ablation(\texttt{spacing\_aware=False})은 비등방 데이터에서만
차이를 낼 수 있다. ACDC의 층간격은 통상 5--10\,mm이고 면내 간격은 1.4--1.8\,mm이다.
편의를 위해 재표본화된 공개 사본을 쓰면 ablation 결과가 \emph{평탄}하게 나오는데,
평탄한 결과는 ``효과가 없다''는 음성 결과로 읽히지 실제 원인인 ``측정 자체가
불가능하다''로 읽히지 않는다. 검증기는 이를 두 단계로 막는다. 첫째, 층간격 대 면내
간격 비가 2 미만이면 거부한다. 둘째--- 그리고 이것이 필요한 이유가 있다---층간격을
보존한 채 면내만 재표본화한 사본(예: $1\times1\times10$\,mm로 재표본화하고
$192\times192$로 중앙 crop한 공개 미러)은 비율 검사를 \emph{통과}하므로, 알려진
미러를 spacing과 면내 크기 지문으로 직접 대조한다. 또한 다기관 데이터인데 모든 피험자의
spacing이 동일하다면 그 자체를 재표본화의 증거로 보고 거부한다.}

이 검증기는 NIfTI 헤더를 표준 라이브러리만으로 해석하므로 \texttt{nibabel}이나
PyTorch를 설치하기 전에 실행된다. 즉 데이터를 받은 직후 가장 먼저 실행할 수 있으며,
잘못된 사본으로 전체 실험을 돌린 뒤에야 깨닫는 일을 방지한다. 헤더 해석과 매니페스트
정합성은 \texttt{scripts/verify\_kernels.py}의 제 15--16절 검사로 덮여 있고 실제로
실행된다(표 \ref{tab:kernel-results}).

현실적 표본 범위는 팬텀으로 구현 검증, ACDC 20--30명으로 핵심 효율 및 품질 비교,
M\&Ms-2 10--20명으로 외부 도메인 검증이다. 환자별 최적화 방식이므로 대규모 학습은
필요하지 않다.

\section{지표}
\label{sec:metrics}

\paragraph{전처리 효율.}
첫 프레임 2DGS 구축시간, 프레임별 Chan--Vese 시간, 정사영/orientation/잔차 계산시간,
전체 wall-clock time과 peak GPU memory를 \emph{구분하여} 측정한다. 식
\eqref{eq:cost-hypothesis}을 항별로 확인하기 위함이다.

\paragraph{비용은 통화가 둘이다.}
본 파이프라인은 성격이 전혀 다른 두 곳에서 시간을 쓰며, 이를 합산하거나 구분 없이 나란히
적으면 비교 불가능한 것을 비교하게 만든다. \emph{전처리}(Chan--Vese, 정사영, 전송, 밀도
제어, 잔차 해)는 환자당 오프라인으로 \textbf{한 번} 지불하므로 수 분이 걸려도 사용에
지장이 없다. \emph{재생}(표면 읽기, 정사영, 방향, rasterisation)은 표시되는 \textbf{매}
프레임마다 지불하므로 33.3\,ms 예산이 적용되는 쪽은 이것이다. 저장량은 시간이 아니라
공간으로 한 번 지불하며 위 둘과 교환된다. 세 통화 사이에 정당한 환율은 없으므로 본 논문은
단일 비용 점수를 만들지 않는다. \texttt{metrics/costquality.py}의 \texttt{Cost}는
스칼라 변환 요청에 예외를 던져 이를 구조적으로 막는다---환율을 발명하면 프레임당 성능
저하가 전처리 속도 향상 뒤에 숨을 수 있기 때문이다.

\paragraph{재생.}
표면 loading, surfel 갱신, 래스터화, UI overhead를 각각 측정하고 평균 latency, 95번째
백분위 latency, FPS를 보고한다. 성공 기준은 \emph{95번째 백분위}가 33.3\,ms 미만인
것이다. 평균만 보는 기준은 열 프레임마다 끊기는 viewer를 통과시킨다.

여기에 두 가지를 더한다. 첫째, \textbf{단계별 $p_{95}$}를 보고한다(표
\ref{tab:frame-budget-stages}). 총 $p_{95}$만으로는 끊긴다는 사실만 알 수 있고 어느 단계
때문인지는 알 수 없는데, 네 단계의 해결책은 서로 완전히 다르다. 또한 평균 비용이 가장 큰
단계가 꼬리를 만드는 단계와 일치하지 않을 수 있으므로---평소에는 저렴하지만 타일 예산을
간헐적으로 초과하는 rasteriser가 그 예다---귀속은 가장 느린 프레임들에 대해 계산한다.
둘째, \textbf{헤드룸} $33.3/p_{95}$를 보고한다. 통과/실패 이진값만으로는 ``간신히
통과''와 ``세 배 여유''를 구별할 수 없는데, 이 둘은 surfel 수나 해상도를 늘렸을 때 결과가
유지될지에 대해 전혀 다른 함의를 갖는다.

\paragraph{저장량.}
식 \eqref{eq:s-full}--\eqref{eq:cr}을 바이트 단위로, 색인과 점유 비트마스크 오버헤드를
포함하여 계산한다. 표면을 mesh, binary mask, 좁은 띠 부호 거리로 저장하는 세 경우를
\emph{모두} 보고하고, 이론적 값과 실제 파일 크기를 나란히 제시한다. 식
\eqref{eq:breakeven}의 손익분기 예산도 함께 보고한다.

\paragraph{표면 정확도.}
ED/ES label에 대해 Dice, HD95, 평균 대칭 표면 거리를 mm 단위로 계산한다. surfel 중심의
표면 잔차 $\Esurf$(식 \eqref{eq:e-surf})도 측정한다. 두 지표는 다른 질문에 답한다.
Dice/HD95는 표면이 해부학적으로 옳은지를, $\Esurf$는 anchor가 그 표면 위에 실제로
놓였는지를 묻는다. anchor가 잘못된 표면 위에 완벽하게 놓일 수도 있고, 옳은 표면
주위에 흩어질 수도 있다.

\paragraph{렌더링 품질.}
투영된 Chan--Vese mask와 rendered alpha mask 사이의 silhouette IoU(문턱값 스윕)와
boundary F-score, 동일 표면에서 얻은 참조 depth/normal에 대한 depth RMSE와 normal
angular error, hole 비율과 과도한 overlap 비율을 측정한다.

\paragraph{외형.}
원래 뷰로 투영한 영상과 관측 영상 사이의 PSNR, SSIM, NRMSE를 전체 영상과 심장 ROI에
대해 \emph{각각} 보고한다. 검은 배경이 전체 영상 PSNR을 무의미하게 부풀리기 때문이다.

\paragraph{시간적 안정성.}
동일한 고정 시점에서 인접 렌더링 프레임의 차이를 입력의 실제 프레임 차이와 비교한다.
\begin{equation}
  \Eflicker = \frac{1}{T-1}\sum_t
  \norm{\bigl(\Chat_t - \Chat_{t-1}\bigr) - \bigl(C^{\mathrm{ref}}_t - C^{\mathrm{ref}}_{t-1}\bigr)}_1.
  \label{eq:flicker}
\end{equation}
참조 차이를 빼는 것이 핵심이다. 그래야 정지한 모델이 보상받지 않는다. 시간적으로
얼어붙은 렌더러는 이 지표에서 0을 얻지 못한다.

\paragraph{품질과 비용을 함께 읽는 방법.}
\label{sec:cost-quality}
품질과 비용을 같은 행에 두는 것은 필요하지만 충분하지 않다. 열이 열다섯 개인 표는 둘 다
\emph{측정했다}는 사실을 보여줄 뿐 \emph{상충}이 무엇인지는 보여주지 않는다. 프레임당
지연을 세 배로 늘려 PSNR을 0.3\,dB 얻은 방법은 넓은 표에서는 그냥 더 좋아 보인다.

환율 없이도 말할 수 있는 것은 \textbf{지배(dominance)}다. 어떤 설정이 측정된 모든 축에서
다른 설정보다 나쁘거나 같고 최소 한 축에서 엄격히 나쁘다면, 그 설정은 지배당하며 더 논할
필요가 없다. 표 \ref{tab:cost-quality}의 Pareto 열이 이것을 보고한다. frontier에 여럿이
남으면 실제 상충이 존재하는 것이고, 어느 것을 택할지는 수치가 아니라 응용에 대한 판단이다.

\caveat{%
지배 판정은 \emph{두 설정이 모두 측정한} 축에 대해서만 계산한다. 한쪽만 측정한 축은
추측하지 않고 제외한다. 그 결과 어떤 방법이 \emph{불리한 축을 보고하지 않음으로써} 이기는
일이 구조적으로 불가능하다---보고하지 않으면 그 축은 판정에서 빠지고 이익도 생기지 않는다.
\texttt{shared\_axis\_count()}가 판정의 근거가 된 축 수를 함께 보고하므로, 축 하나에만
의존한 판정을 독자가 알아볼 수 있다.

자유 기하 3DGS가 좋은 예다. PSNR에서 이기고 $\Esurf$에서 크게 지도록 설계상 예상되며,
따라서 두 설정 모두 frontier에 남는다. 그 상충이 평균으로 뭉개지지 않고 그대로 보이는
것이 본 논문의 논점이다.}

%% AUTO-GENERATED by scripts/make_tables.py -- DO NOT EDIT BY HAND.
%% Source: results/cost_quality.json
%% Status: STUB - no measurements available
\begin{table}[htbp]
\centering\small
\caption{품질과 비용을 함께. 비용 통화가 \emph{두 개}임에 유의한다. 전처리는 오프라인에서
한 번만 지불하고, p95 프레임 지연은 표시되는 \emph{매} 프레임마다 지불하므로 식
\eqref{eq:frame-budget}의 33.3\,ms 예산이 적용되는 쪽은 후자다. 헤드룸은 $33.3/p_{95}$이며
1을 넘으면 통과이고 그 값이 여유를 말한다. Pareto 열이 이 표를 진술로 만든다---dominated는
다른 설정보다 모든 축에서 나쁘거나 같으므로 더 논할 필요가 없고, frontier가 여럿이면 실제
상충이 존재한다는 뜻이다. 세 통화 사이에 정당한 환율이 없으므로 단일 점수는 계산하지 않는다.}
\label{tab:cost-quality}
\begin{tabular}{@{}lrrrrrrrl@{}}
\toprule
 & \multicolumn{3}{c}{품질} & \multicolumn{2}{c}{비용} & & & \\
\cmidrule(lr){2-4}\cmidrule(lr){5-6}
설정 & PSNR$_{\mathrm{ROI}}$ & $\Esurf$ & $\Eflicker$
  & 전처리 & $p_{95}$ & 헤드룸 & 저장량 & Pareto \\
 & (dB) & (mm) & & (s) & (ms) & ($\times$) & (MB) & \\
\midrule
\texttt{cv-dyn2dgs} & \NM & \NM & \NM & \NM & \NM & \NM & \NM & \NM \\
\texttt{independent-2dgs} & \NM & \NM & \NM & \NM & \NM & \NM & \NM & \NM \\
\texttt{mesh-only} & \NM & \NM & \NM & \NM & \NM & \NM & \NM & \NM \\
\texttt{thin-3dgs} & \NM & \NM & \NM & \NM & \NM & \NM & \NM & \NM \\
\texttt{free-3dgs-surface} & \NM & \NM & \NM & \NM & \NM & \NM & \NM & \NM \\
\texttt{v1-minimal} & \NM & \NM & \NM & \NM & \NM & \NM & \NM & \NM \\
\bottomrule
\end{tabular}
\end{table}

%% AUTO-GENERATED by scripts/make_tables.py -- DO NOT EDIT BY HAND.
%% Source: results/research_questions.json
%% Status: STUB - no measurements available
\begin{table}[htbp]
\centering\small
\caption{프레임 예산의 단계별 분해를 평균이 아니라 $p_{95}$로 보고한다(식
\eqref{eq:frame-budget}). 총 $p_{95}$만으로는 재생이 끊긴다는 사실만 알 수 있고 어느
단계 때문인지는 알 수 없는데, 네 단계의 해결책은 서로 완전히 다르다. 평균 비용이 가장 큰
단계가 꼬리를 만드는 단계와 일치하지 않을 수 있으므로---예컨대 평소에는 저렴하지만 타일
예산을 간헐적으로 초과하는 rasteriser---귀속은 가장 느린 프레임들에 대해 계산한다.}
\label{tab:frame-budget-stages}
\begin{tabular}{@{}lrrrrrr@{}}
\toprule
설정 & 표면 & 정사영 & 방향 & rasteriser & 총 $p_{95}$ & 헤드룸 \\
 & (ms) & (ms) & (ms) & (ms) & (ms) & ($\times$) \\
\midrule
\NM & \NM & \NM & \NM & \NM & \NM & \NM \\
\bottomrule
\end{tabular}
\end{table}


\paragraph{임상적으로 해석 가능한 지표.}
Chan--Vese 표면에서 계산한 좌심실 volume curve와
\begin{equation}
  \mathrm{EF} = \frac{V_{\mathrm{ED}}-V_{\mathrm{ES}}}{V_{\mathrm{ED}}}\times100\%
  \label{eq:ef}
\end{equation}
를 label 기반 값과 비교한다. 이는 viewer의 시각적 품질과 분할의 계량적 유효성을
분리해서 보여준다. EF 오차는 퍼센트 포인트로 보고한다.

\section{Ablation}
\label{sec:ablation-design}

전체 설정에 대해 한 번에 하나의 기제만 토글한다.

\begin{itemize}[leftmargin=2em]
  \item 이전 level-set seed 사용 여부,
  \item \textbf{물리적 복셀 간격 반영 여부} --- 첫 번째 기여의 유일한 직접 시험. 등방
  격자와 $h_z/h_x\approx6.4$인 격자에서 비교하며, 격차가 비등방성과 함께 커지지 않으면
  그 기여는 이론이 주장하는 일을 하지 않는 것이다.
  \item 정사영 반복 횟수 $K_p\in\{1,2,3,5\}$ (명제 \ref{prop:newton-steps}),
  \item closest-point 대 normal-only 정사영,
  \item 등방 대 비등방 원반 (명제 \ref{prop:gauge}, \ref{prop:aniso}),
  \item 2DGS, thin-3DGS, mesh primitive 비교,
  \item mask/normal/depth-distortion 손실의 개별 효과,
  \item 잔차 및 시간 정칙화 사용 여부,
  \item 전체 대 저계수 잔차 (명제 \ref{prop:lowrank}),
  \item repulsion/densification 사용 여부,
  \item surfel 수와 viewer 해상도에 따른 FPS 변화.
\end{itemize}

\section{개발 단계의 평가}

방법을 세 단계로 나누어 각 단계가 실제로 무엇을 얻는지 측정한다.

\begin{table}[htbp]
\centering
\small
\caption{개발 단계. 각 단계가 추가하는 기제.}
\label{tab:stages}
\begin{tabular}{@{}llp{7.2cm}@{}}
\toprule
단계 & 범위 & 추가 기제 \\
\midrule
v1 & 최소 성공 & 등방 원반, $K_p=1$, 기하 손실 없음, 밀도 제어 없음 \\
v2 & --- & 식 \eqref{eq:total-loss}의 기하 손실, 비등방 최소회전 전송 \\
v3 & 확장 & 곡률 적응 scale(식 \eqref{eq:curvature-scale}), repulsion, 제한된
  densification/pruning, 시간 정칙화, 저계수 잔차 \\
\bottomrule
\end{tabular}
\end{table}

품질, 비용, 저장량을 함께 보고하므로 추가 계산만으로 산 품질 향상은 그것으로 드러난다.

\section{가설과 통계}
\label{sec:statistics}

가설을 하나의 모호한 ``품질'' 점수로 합치지 않고 세 부분으로 분리한다.

\begin{enumerate}[leftmargin=2em]
  \item CV-Dyn2DGS는 independent 2DGS보다 전체 전처리시간과 저장량이 작으면서
  silhouette IoU, normal error, 심장 ROI SSIM에서 사전에 정한 비열등성 마진 안에
  들어온다.
  \item 동일 표면과 비슷한 primitive 수에서 2DGS는 mesh 및 thin-3DGS보다 silhouette,
  normal/depth, appearance 중 적어도 하나를 개선한다.
  \item viewer의 95번째 백분위 latency는 33.3\,ms 미만이다.
\end{enumerate}

시간과 품질의 한쪽만 보고 결론내리지 않기 위해 각 방법의 전처리시간--품질,
저장량--품질, FPS--품질 Pareto curve를 제시한다. 동일 환자에 여러 방법을 적용한
paired comparison을 사용하고, 환자별 평균을 독립 표본으로 두어 bootstrap 95\% 신뢰구간을
보고하며 필요한 경우 Wilcoxon signed-rank test를 사용한다. 여러 프레임을 서로 독립
환자처럼 계산하지 않는다.

\caveat{%
비열등성 마진은 \textbf{본 논문에서 정하지 않았다}. 5명 규모의 pilot 결과와 측정
noise를 보고 고정한 뒤 본 실험에서 바꾸지 않아야 하며, 결과를 본 뒤에 정하면 검정이
무의미해진다. 마진을 결과와 함께 제시하는 것은 순환논증이다.}

\caveat{%
또한 현재 구현은 환자 수준 paired 통계를 \emph{구현하지 않았다}. 지표 산출은
환자별로 이루어지지만 bootstrap과 Wilcoxon 집계는 실제 데이터와 함께 추가되어야 한다.
단일 팬텀 결과는 환자에 대한 증거로 읽을 수 없다.}

\chapter{결과}
\label{ch:results}

\ifdraftmode
\begin{center}
\fcolorbox{red}{red!5}{%
\begin{minipage}{0.93\linewidth}
\small
\textcolor{red}{\textbf{본 장의 상태}}\par\smallskip
제 \ref{sec:kernel-results}절을 제외한 본 장의 모든 표는 \textbf{비어 있다}.
구현은 완성되었으나 실행되지 않았으므로(제 \ref{sec:impl-status}절) RQ1--RQ6의 어떤
수치도 존재하지 않는다. 미측정 항목은 \NM 로 표시하였다.\par\smallskip
표들은 손으로 채우도록 설계되지 않았다. \texttt{cvdyn2dgs experiments}가 산출하는
JSON을 \texttt{scripts/make\_tables.py}가 읽어 \texttt{paper/tables/}의 \texttt{.tex}
파일을 생성하며, 본 장은 그것을 \texttt{\textbackslash input}한다. 숫자를 원고에
직접 타이핑하는 경로는 존재하지 않는다.
\end{minipage}}
\end{center}
\medskip
\fi

\section{실행된 검증: 이산 논리 커널}
\label{sec:kernel-results}

본 절은 본 논문에서 유일하게 \emph{측정된} 결과이다. 제
\ref{sec:kernel-verification}절에서 기술한 커널 검증을 실행한 결과 41개 항목 전부가
통과하였다. PyTorch를 필요로 하지 않으므로 어떤 환경에서도 재현할 수 있다.

%% AUTO-GENERATED by scripts/make_tables.py -- DO NOT EDIT BY HAND.
%% Source: scripts/verify_kernels.py
%% Status: filled from measurements
\begin{table}[htbp]
\centering\small
\caption{실행된 커널 검증 결과. \texttt{python scripts/verify\_kernels.py}로 재현 가능.}
\label{tab:kernel-results}
\begin{tabular}{@{}lrr@{}}
\toprule
검증 묶음 & 항목 수 & 통과 \\
\midrule
큐브 $\to$ 6면체 분해            & 5  & 5 \\
marching tetrahedra 참조표      & 2  & 2 \\
(surfel, 타일) 페어 열거        & 3  & 3 \\
타일별 깊이 정렬과 절단          & 4  & 4 \\
점유 비트마스크 팩/언팩          & 3  & 3 \\
타원체 부호 거리 이분법          & 5  & 5 \\
삼선형 보간 가중치               & 3  & 3 \\
공액경사법                      & 2  & 2 \\
Godunov 상류 재초기화            & 2  & 2 \\
저장량 손익분기 대수             & 6  & 6 \\
수축 프로파일                   & 6  & 6 \\
\midrule
\textbf{합계}                   & \textbf{41} & \textbf{41} \\
\bottomrule
\end{tabular}
\end{table}


이 결과의 범위는 제한적이다. 통과는 이산 색인 논리가 건전함을 뜻하며, 파이프라인이
동작함을 뜻하지 않는다. 특히 래스터화, 자동미분, Chan--Vese 솔버, 전체 파이프라인은
전혀 포함되지 않는다.

\section{이론 수렴률}
\label{sec:results-rates}

제 \ref{ch:error}장의 수렴률 주장은 $\spacing$을 변화시키며 $\log$--$\log$ 기울기를
적합하여 확인한다.

%% AUTO-GENERATED by scripts/make_tables.py -- DO NOT EDIT BY HAND.
%% Source: results/theory_checks.json
%% Status: STUB - no measurements available
\begin{table}[htbp]
\centering\small
\caption{이론이 예측하는 수렴률과 측정된 경험적 지수.}
\label{tab:theory-rates}
\begin{tabular}{@{}lccl@{}}
\toprule
결과 & 예측 & 측정 & 비고 \\
\midrule
명제 \ref{prop:normal-error} $\sin\theta=O(\hmax^2)$
  & 2.0 & \NM & $\spacing$ 스윕 기울기 \\
보조정리 \ref{lem:quadratic} 이차 수축
  & 2.0 & \NM & 잔차 이력 적합 \\
명제 \ref{prop:interp-eikonal} 식 \eqref{eq:interp-eikonal}
  & 0.0 & \NM & 측정 대 예측 편차 \\
명제 \ref{prop:lowrank} Eckart--Young
  & 0.0 & \NM & 상대 편차 \\
\bottomrule
\end{tabular}
\end{table}


\pending{%
명제 \ref{prop:normal-error}이 예측하는 $\sin\theta = O(\hmax^2)$의 경험적 지수,
보조정리 \ref{lem:quadratic}이 예측하는 이차 수축의 경험적 차수, 명제
\ref{prop:planar-disk}의 $O(\kappa s^2)$ 비례상수, 명제 \ref{prop:interp-eikonal}의
측정 대 예측 eikonal 노름. 실행 명령: \texttt{cvdyn2dgs theory}.}

\section{RQ1: 웜스타트}
\label{sec:results-rq1}

%% AUTO-GENERATED by scripts/make_tables.py -- DO NOT EDIT BY HAND.
%% Source: results/research_questions.json
%% Status: STUB - no measurements available
\begin{table}[htbp]
\centering\small
\caption{RQ1: 웜스타트(식 \eqref{eq:warmstart-init})의 효과.}
\label{tab:rq1}
\begin{tabular}{@{}lcc@{}}
\toprule
 & 웜스타트 & 콜드스타트 \\
\midrule
총 반복 횟수            & \NM & \NM \\
프레임당 평균 반복      & \NM & \NM \\
총 시간 (ms)            & \NM & \NM \\
$K(e_0)$ 평균           & \NM & \NM \\
초기 오차 $e_0$         & \NM & \NM \\
평균 crop 비율          & \NM & \NM \\
\bottomrule
\end{tabular}
\end{table}


\pending{%
웜스타트와 콜드스타트의 총 반복 횟수 및 벽시계 시간, 참조 해에 대한 $K(e_0)$(식
\eqref{eq:iteration-bound}), 초기 오차 $e_0$. 정리 \ref{thm:warmstart}은 \emph{순서}만
예측하므로, 확인할 것은 $K(e^{\mathrm{warm}}_0)\le K(e^{\mathrm{cold}}_0)$이다.
벽시계 시간의 격차는 좁은 띠 crop 효과와 섞이므로 반복 횟수와 분리하여 해석해야 한다.}

\section{RQ2: 정사영 대 독립 재최적화}
\label{sec:results-rq2}

%% AUTO-GENERATED by scripts/make_tables.py -- DO NOT EDIT BY HAND.
%% Source: results/research_questions.json
%% Status: STUB - no measurements available
\begin{table}[htbp]
\centering\scriptsize
\setlength{\tabcolsep}{3.5pt}
\caption{RQ2: 법선 정사영 대 프레임별 독립 재최적화.}
\label{tab:rq2}
\begin{tabular}{@{}lrrrrrrrrrrr@{}}
\toprule
방법 & Dice & HD95 & $\Esurf$ & IoU & depth & normal & PSNR & SSIM & CR & 전처리 & p95 \\
     &      & (mm) & (mm)     &     & (mm)  & (deg)  & (dB) &      &    & (ms)   & (ms) \\
\midrule
CV-Dyn2DGS & \NM & \NM & \NM & \NM & \NM & \NM & \NM & \NM & \NM & \NM & \NM \\
Independent 2DGS & \NM & \NM & \NM & \NM & \NM & \NM & \NM & \NM & \NM & \NM & \NM \\
\bottomrule
\end{tabular}
\end{table}

\noindent 전처리 속도 배율: \NM$\times$.
프레임당 갱신 시간: \NM\,ms (제안) 대
\NM\,ms (재최적화).


\pending{%
식 \eqref{eq:cost-hypothesis}의 좌변 네 항($C_{\mathrm{CV}}$, $C_{\mathrm{project}}$,
$C_{\mathrm{orient}}$, $C_{\mathrm{res}}$)을 각각, 그리고 우변 $C_{\mathrm{full}}$을
측정하여 부등식이 성립하는지 확인한다. 품질 측면에서는 independent 2DGS를 상한으로
두고 비열등성 마진 안에 들어오는지 본다. 목표는 이기는 것이 아니라 마진 안에 있으면서
비용이 작은 것이다.}

\section{RQ3: 저장량}
\label{sec:results-rq3}

%% AUTO-GENERATED by scripts/make_tables.py -- DO NOT EDIT BY HAND.
%% Source: results/research_questions.json
%% Status: STUB - no measurements available
\begin{table}[htbp]
\centering\small
\caption{RQ3: 저장량과 압축비(식 \eqref{eq:s-full}--\eqref{eq:cr}). 세 가지 표면
저장 방식 모두를 보고한다.}
\label{tab:rq3}
\begin{tabular}{@{}lrrrr@{}}
\toprule
표면 저장 방식 & $\Sfull$ (MB) & $\Sours$ (MB) & $\CR$ & 손익분기 (B/프레임) \\
\midrule
좁은 띠 부호 거리 & \NM
  & \NM
  & \NM
  & \NM \\
binary mask & \NM
  & \NM
  & \NM
  & \NM \\
triangle mesh & \NM
  & \NM
  & \NM
  & \NM \\
\midrule
실제 파일 크기 (MB) & \multicolumn{4}{l}{\NM} \\
\bottomrule
\end{tabular}
\end{table}


\pending{%
$\Sfull$, $\Sours$, $\CR$을 세 가지 표면 저장 방식에 대해, 그리고 이론적 바이트 수와
실제 파일 크기를 나란히. 식 \eqref{eq:breakeven}의 손익분기 예산과 실제 프레임당 표면
크기의 비교가 핵심이다. $\CR<1$이면 제안 표현이 이 구성에서 이기지 못한다는 뜻이며,
그것도 유효한 결과이다.}

\section{RQ4: 재생 속도}
\label{sec:results-rq4}

%% AUTO-GENERATED by scripts/make_tables.py -- DO NOT EDIT BY HAND.
%% Source: results/research_questions.json
%% Status: STUB - no measurements available
\begin{table}[htbp]
\centering\small
\caption{RQ4: 재생 latency 분해(식 \eqref{eq:frame-budget}). 성공 기준은 p95 $<33.3$\,ms.}
\label{tab:rq4}
\begin{tabular}{@{}lrrrrrrrc@{}}
\toprule
해상도 & $N$ & surface & project & orient & raster & 평균 & p95 & 30 FPS \\
       &     & (ms)    & (ms)    & (ms)   & (ms)   & (ms) & (ms) &        \\
\midrule
256 & 10000 & \NM & \NM & \NM & \NM & \NM & \NM & \NM \\
256 & 20000 & \NM & \NM & \NM & \NM & \NM & \NM & \NM \\
512 & 10000 & \NM & \NM & \NM & \NM & \NM & \NM & \NM \\
512 & 20000 & \NM & \NM & \NM & \NM & \NM & \NM & \NM \\
\bottomrule
\end{tabular}
\end{table}


\pending{%
식 \eqref{eq:frame-budget}의 네 항을 분해하여, surfel 수와 viewer 해상도를 스윕하며
95번째 백분위 latency가 33.3\,ms 미만인지 확인한다. 또한 제 \ref{sec:seek-problem}절의
순차 대 정준 정사영 격차를 측정한다. 이 격차가 크면 임의 탐색을 식 \eqref{eq:store}의
저장량만으로 제공할 수 없다는 뜻이다.}

\section{RQ5: mesh와의 비교}
\label{sec:results-rq5}

%% AUTO-GENERATED by scripts/make_tables.py -- DO NOT EDIT BY HAND.
%% Source: results/research_questions.json
%% Status: STUB - no measurements available
\begin{table}[htbp]
\centering\small
\caption{RQ5: 동일 Chan--Vese 표면에서 2DGS와 텍스처 mesh의 지표별 비교. 단일 점수로
합치지 않고 지표별 승자를 보고한다.}
\label{tab:rq5}
\begin{tabular}{@{}lrrll@{}}
\toprule
지표 & 2DGS & mesh & 방향 & 승자 \\
\midrule
silhouette IoU & \NM & \NM & 높을수록 & \NM \\
boundary F-score & \NM & \NM & 높을수록 & \NM \\
depth RMSE (mm) & \NM & \NM & 낮을수록 & \NM \\
normal error (deg) & \NM & \NM & 낮을수록 & \NM \\
ROI PSNR (dB) & \NM & \NM & 높을수록 & \NM \\
ROI SSIM & \NM & \NM & 높을수록 & \NM \\
$\Eflicker$ & \NM & \NM & 낮을수록 & \NM \\
\bottomrule
\end{tabular}
\end{table}


\pending{%
동일 표면, 비슷한 primitive 예산에서 silhouette(문턱값 스윕), boundary F-score, depth
RMSE, normal angular error, ROI PSNR/SSIM, $\Eflicker$를 지표별로 비교한다. 단일 점수로
합치지 않고 지표별 승패를 보고한다. 2DGS가 어느 지표에서도 개선하지 못하면 ``이
목적에는 mesh가 더 적절하다''가 올바른 결론이다.}

\section{RQ6: 2D 원반 대 얇은 3D Gaussian}
\label{sec:results-rq6}

%% AUTO-GENERATED by scripts/make_tables.py -- DO NOT EDIT BY HAND.
%% Source: results/research_questions.json
%% Status: STUB - no measurements available
\begin{table}[htbp]
\centering\small
\caption{RQ6: 2D 원반과 얇은 3D Gaussian. $\rho$는 법선방향 두께 비.}
\label{tab:rq6}
\begin{tabular}{@{}lrrrr@{}}
\toprule
표현 & IoU & depth RMSE (mm) & normal (deg) & ROI PSNR (dB) \\
\midrule
2DGS (2D 원반) & \NM & \NM & \NM & \NM \\
\midrule
thin-3DGS rho=0.05 & \NM & \NM & \NM & \NM \\
thin-3DGS rho=0.1 & \NM & \NM & \NM & \NM \\
thin-3DGS rho=0.3 & \NM & \NM & \NM & \NM \\
\bottomrule
\end{tabular}
\end{table}


\pending{%
법선방향 두께 비 $\rho$를 스윕하며 depth RMSE와 normal angular error를 비교한다.
thin-3DGS의 깊이는 자락 전역에서 상수이므로 2DGS의 깊이 우위는 명제
\ref{prop:perspective}의 직접 검증이며 독립적 발견이 아니다. 더 얇은 타원체는
기하학적으로 원반에 가깝지만 조건수가 나빠지므로 그 절충이 스윕에서 드러난다.}

\section{기준선 종합}
\label{sec:results-baselines}

%% AUTO-GENERATED by scripts/make_tables.py -- DO NOT EDIT BY HAND.
%% Source: results/baselines.json
%% Status: STUB - no measurements available
\begin{table}[htbp]
\centering\scriptsize
\setlength{\tabcolsep}{3.5pt}
\caption{기준선 종합 비교 (표 \ref{tab:baselines}).}
\label{tab:baselines-results}
\begin{tabular}{@{}lrrrrrrrrrrr@{}}
\toprule
방법 & Dice & HD95 & $\Esurf$ & IoU & depth & normal & PSNR & SSIM & CR & 전처리 & p95 \\
     &      & (mm) & (mm)     &     & (mm)  & (deg)  & (dB) &      &    & (ms)   & (ms) \\
\midrule
cv-dyn2dgs & \NM & \NM & \NM & \NM & \NM & \NM & \NM & \NM & \NM & \NM & \NM \\
independent-2dgs & \NM & \NM & \NM & \NM & \NM & \NM & \NM & \NM & \NM & \NM & \NM \\
mesh-only & \NM & \NM & \NM & \NM & \NM & \NM & \NM & \NM & \NM & \NM & \NM \\
thin-3dgs & \NM & \NM & \NM & \NM & \NM & \NM & \NM & \NM & \NM & \NM & \NM \\
param-copy & \NM & \NM & \NM & \NM & \NM & \NM & \NM & \NM & \NM & \NM & \NM \\
closest-point & \NM & \NM & \NM & \NM & \NM & \NM & \NM & \NM & \NM & \NM & \NM \\
normal-no-residual & \NM & \NM & \NM & \NM & \NM & \NM & \NM & \NM & \NM & \NM & \NM \\
\bottomrule
\end{tabular}
\end{table}


\pending{%
표 \ref{tab:baselines}의 모든 기준선에 대한 주요 지표. 실행 명령:
\texttt{cvdyn2dgs experiments --size small}.}

\section{Ablation}
\label{sec:results-ablation}

%% AUTO-GENERATED by scripts/make_tables.py -- DO NOT EDIT BY HAND.
%% Source: results/ablations.json
%% Status: STUB - no measurements available
\begin{table}[htbp]
\centering\scriptsize
\setlength{\tabcolsep}{3.5pt}
\caption{Ablation (제 \ref{sec:ablation-design}절). 전체 설정에서 한 번에 하나의 기제만
토글한다.}
\label{tab:ablations}
\begin{tabular}{@{}lrrrrrr@{}}
\toprule
설정 & $\Esurf$ (mm) & IoU & depth (mm) & normal (deg) & PSNR (dB) & 전처리 (ms) \\
\midrule
full & \NM & \NM & \NM & \NM & \NM & \NM \\
no-warm-start & \NM & \NM & \NM & \NM & \NM & \NM \\
kp=1 & \NM & \NM & \NM & \NM & \NM & \NM \\
kp=2 & \NM & \NM & \NM & \NM & \NM & \NM \\
kp=3 & \NM & \NM & \NM & \NM & \NM & \NM \\
kp=5 & \NM & \NM & \NM & \NM & \NM & \NM \\
closest-point & \NM & \NM & \NM & \NM & \NM & \NM \\
isotropic & \NM & \NM & \NM & \NM & \NM & \NM \\
no-mask-loss & \NM & \NM & \NM & \NM & \NM & \NM \\
no-normal-loss & \NM & \NM & \NM & \NM & \NM & \NM \\
no-dist-loss & \NM & \NM & \NM & \NM & \NM & \NM \\
no-geometry-losses & \NM & \NM & \NM & \NM & \NM & \NM \\
no-residual & \NM & \NM & \NM & \NM & \NM & \NM \\
no-temporal-reg & \NM & \NM & \NM & \NM & \NM & \NM \\
full-rank-residual & \NM & \NM & \NM & \NM & \NM & \NM \\
lowrank-r4 & \NM & \NM & \NM & \NM & \NM & \NM \\
no-repulsion & \NM & \NM & \NM & \NM & \NM & \NM \\
no-densify & \NM & \NM & \NM & \NM & \NM & \NM \\
no-density-control & \NM & \NM & \NM & \NM & \NM & \NM \\
no-curvature-scale & \NM & \NM & \NM & \NM & \NM & \NM \\
\bottomrule
\end{tabular}
\end{table}


\pending{%
제 \ref{sec:ablation-design}절의 각 항목. 특히 간격 인지 여부의 ablation은 첫 번째
기여의 유일한 직접 시험이므로, 등방 격자와 비등방 격자에서의 격차를 함께 제시해야
한다.}

\section{개발 단계}
\label{sec:results-stages}

%% AUTO-GENERATED by scripts/make_tables.py -- DO NOT EDIT BY HAND.
%% Source: results/progressive.json
%% Status: STUB - no measurements available
\begin{table}[htbp]
\centering\scriptsize
\setlength{\tabcolsep}{3.5pt}
\caption{개발 단계별 결과 (표 \ref{tab:stages}).}
\label{tab:stages-results}
\begin{tabular}{@{}lrrrrrrrrrrr@{}}
\toprule
방법 & Dice & HD95 & $\Esurf$ & IoU & depth & normal & PSNR & SSIM & CR & 전처리 & p95 \\
     &      & (mm) & (mm)     &     & (mm)  & (deg)  & (dB) &      &    & (ms)   & (ms) \\
\midrule
v1-minimal & \NM & \NM & \NM & \NM & \NM & \NM & \NM & \NM & \NM & \NM & \NM \\
v2-geometry & \NM & \NM & \NM & \NM & \NM & \NM & \NM & \NM & \NM & \NM & \NM \\
v3-adaptive & \NM & \NM & \NM & \NM & \NM & \NM & \NM & \NM & \NM & \NM & \NM \\
\bottomrule
\end{tabular}
\end{table}


\pending{%
v1 $\to$ v2 $\to$ v3의 품질, 비용, 저장량 변화. 각 단계가 무엇을 얻고 무엇을
지불하는지를 함께 보아야 하며, 추가 계산만으로 산 향상은 그것으로 드러난다.}

\chapter{논의}
\label{ch:discussion}

\ifdraftmode
\begin{center}
\fcolorbox{red}{red!5}{%
\begin{minipage}{0.93\linewidth}
\small
\textcolor{red}{\textbf{본 장의 상태}}\par\smallskip
결과가 존재하지 않으므로 본 장은 결과 \emph{해석}을 담을 수 없다. 대신 결과를
보기 전에 확정해 두는 \textbf{해석 규칙}을 기술한다. 어떤 결과가 나오면 어떤 결론을
내릴 것인지를 사전에 고정하는 것은, 결과를 본 뒤에 유리한 서사를 고르는 것을
막는 유일한 방법이다.
\end{minipage}}
\end{center}
\medskip
\fi

\section{사전에 확정한 해석 규칙}
\label{sec:interpretation-rules}

본 연구의 결과는 네 가지 중 하나로 나올 수 있으며, 각 경우의 결론을 미리 정한다.

\paragraph{경우 1: 빠르고 품질 유지.}
식 \eqref{eq:cost-hypothesis}이 성립하고, 저장량이 $\CR>1$이며, 품질이 independent
2DGS 대비 사전 마진 안에 들어오는 경우. 이는 계산 재사용과 compact representation
가설을 지지한다. 다만 이것만으로 mesh보다 낫다는 결론은 나오지 않는다. RQ5가 별도로
답해야 한다.

\paragraph{경우 2: 빠르지만 품질 저하.}
비용은 줄지만 품질이 마진을 벗어나는 경우. 결론은 ``방법이 실패했다''가 아니라
``절충점이 어디인가''이다. 잔차 계수, surfel 수, 정사영 보정 횟수에 따른 Pareto
curve를 제시하여 어느 계산 예산에서 어떤 표현이 유리한지 보인다.

\paragraph{경우 3: 2DGS가 mesh보다 이득 없음.}
RQ5의 어떤 지표에서도 2DGS가 개선하지 못하는 경우. 이 경우의 결론은 명확하다.
\emph{단색 또는 단순 표면 시각화에는 mesh가 더 적절하다}. 이는 유효한 연구 결과이며
본 논문은 이를 부정적 결과로 취급하지 않는다. 그 경우에도 제 \ref{ch:method},
\ref{ch:error}장의 이론적 기여는 유지된다. 법선 정사영과 오차 분석은 primitive 선택과
독립적이며, mesh 정점에도 같은 방식으로 적용된다.

\paragraph{경우 4: 정사영이 병목.}
$C_{\mathrm{project}}$가 지배적인 경우. surfel 수, 정사영 반복 수, 표면 자료구조에
따른 손익분기점을 분석한다. 특히 명제 \ref{prop:newton-steps}이 예측하는 $K_p\in\{1,2\}$가
실제로 충분한지가 이 경우의 핵심이다.

어떤 결과가 나오더라도 ``표면 변화만으로 어느 수준의 동적 Gaussian 표현을 저비용으로
재사용할 수 있는지''는 밝힐 수 있다.

\section{결과를 보기 전에 정해야 하는 것}

\caveat{%
제 \ref{sec:statistics}절에서 밝힌 대로 비열등성 마진은 본 논문에서 정하지 않았다.
5명 규모 pilot의 측정 noise를 보고 고정해야 하며, 결과를 본 뒤에 정하면 검정이
무의미해진다. 따라서 경우 1과 경우 2를 구분하는 기준선 자체가 아직 존재하지 않는다.
이는 원고의 미완성이 아니라 실험 설계의 필수 순서이다.}

\section{예상되는 실패 모드}

이론이 미리 알려 주는 취약점이 있으며, 실험은 이들을 \emph{확인하러} 설계되었다.

\begin{itemize}[leftmargin=2em]
  \item \textbf{큰 프레임 간 이동.} 법선 정사영은 인접 프레임 사이의 이동이 작고 동일한
  표면 분기가 유지될 때 가장 안정적이다(제 \ref{sec:tangential}절). 수축 프로파일이
  가파른 구간에서 trust region 발동 빈도가 오를 것으로 예상된다.
  \item \textbf{높은 곡률.} 명제 \ref{prop:planar-disk}의 $O(\kappa s^2)$ 오차는 심첨부와
  papillary muscle 근방에서 가장 크다. 곡률 적응 scale(식 \eqref{eq:curvature-scale})이
  이를 억제하는지가 v3 단계의 핵심 검증이다.
  \item \textbf{밝기 비균질.} papillary muscle과 blood pool이 섞이면 Chan--Vese의
  조각마다 상수 가정이 깨진다. 팬텀은 이를 의도적으로 포함하므로 분할 정확도의 저하가
  관측될 것으로 예상된다.
  \item \textbf{방향 표류.} 명제 \ref{prop:propagation}에 의해 위치는 재설정되지만 방향은
  누적된다. 긴 시퀀스에서 $\Eflicker$가 상승하는지, 그리고 등방 원반이 실제로 이를
  무해하게 만드는지(명제 \ref{prop:gauge}) 확인해야 한다.
  \item \textbf{임의 탐색.} 제 \ref{sec:seek-problem}절의 순차 대 정준 정사영 격차.
  이것이 크면 식 \eqref{eq:store}의 저장량 주장과 임의 탐색 기능이 양립하지 않는다.
\end{itemize}

\section{이론과 구현의 관계에 대하여}

본 연구에서 구현은 이론의 부속물이 아니었다. 구현 과정에서 이론 문서가 다루지 않은
세 가지 질문이 강제로 드러났다.

첫째, \emph{감독 신호가 정확히 무엇인가}. 제안서는 ``알려진 slice plane의 mask, depth,
normal, intensity''로 제한했으나 그것을 계산 가능한 목표로 바꾸는 정의가 없었다.
제 \ref{sec:impl-decisions}절의 광선 행진 정의가 그 공백을 채운다.

둘째, \emph{기하가 최적화되지 않음을 어떻게 보장하는가}. 이론적으로는 자명하지만
구현에서는 설정 하나로 깨질 수 있다. 버퍼 등록이 이를 구조적으로 만든다.

셋째, \emph{저장 표현이 임의 탐색을 지원하는가}. 식 \eqref{eq:store}만 보면 자명해
보이지만, 전처리의 순차 정사영과 재생의 직접 정사영이 다르다는 사실은 구현하지
않으면 드러나지 않는다.

이 세 가지는 이론의 오류가 아니라 이론이 명시하지 않은 자유도였다. 구현이 그것을
특정 방식으로 고정했고, 본 논문은 그 선택을 명시적으로 기록하였다.

\chapter{한계}
\label{ch:limitations}

본 장의 항목들은 조용히 고칠 결함이 아니라 방법 또는 현재 연구 상태의 성질이다.

\section{방법의 한계}

\paragraph{조각마다 상수 가정.}
기본 Chan--Vese는 내부와 외부의 밝기가 각각 일정하다는 단순한 가정을 사용한다.
Papillary muscle, blood pool, myocardium이 섞인 영역에서는 국소 영역 항이나 multiphase
model이 필요할 수 있다. 본 연구의 팬텀은 이 위반을 의도적으로 포함하지만, 그것을
해결하지는 않는다.

\paragraph{접선 운동의 미결정성.}
법선 정사영은 식 \eqref{eq:least-norm}의 최소 노름 해를 택하므로 접선 방향 성분을
결정하지 않는다. 따라서 결과는 기하학적 표면 운동이며 실제 조직 대응이나 strain이
아니다. Vertex correspondence가 중요한 기능적 계측에는 mesh 운동 추적이 더 적합하다.

\paragraph{위상 변화와 근접 분기.}
정사영은 인접 프레임 사이의 이동이 작고 동일한 표면 분기가 유지될 때 안정적이다.
큰 이동, 높은 곡률, 가까운 두 경계, 위상 변화에서는 잘못된 표면으로 이동할 수 있다.
Trust region과 되돌림이 이를 줄이지만 제거하지는 못한다.

\paragraph{표면 렌더러의 범위.}
2DGS는 명시적 표면 렌더러이다. 심장 내부 전체의 정확한 volumetric MRI나 임의 단면을
복원하려면 별도 volumetric model이 필요하며, 이는 본 연구의 주장이 아니다.

\paragraph{잔차 저장량.}
진폭 잔차를 surfel마다 모든 프레임에 저장하면 surfel 수가 매우 클 때 저장량이
증가한다. 식 \eqref{eq:lowrank}의 저계수 압축이 이를 완화하지만, 명제
\ref{prop:lowrank}은 오차가 특잇값 꼬리와 같음을 말할 뿐 꼬리가 작음을 보장하지 않는다.

\paragraph{입력 모달리티의 차이.}
원 2DGS는 calibrated multi-view RGB 영상을 가정하지만 cine CMR은 물리 좌표가 알려진
여러 MRI slice이다. 본 연구는 2DGS를 새로운 장면복원기로 쓰지 않고 알려진 geometry로
투영 가능한 표면 렌더러로 사용한다. 감독도 알려진 slice plane의 네 신호로 제한된다.

\paragraph{원근 정합의 계산 비용.}
화소별 광선--surfel 교차는 단순 삼각형 래스터화보다 primitive당 계산이 더 클 수 있다.
따라서 실시간성과 mesh 대비 품질 이득은 \emph{별개의} 가설이며 둘 다 실제 하드웨어에서
측정해야 한다.

\paragraph{중간 시점 보간.}
명제 \ref{prop:interp-eikonal}, \ref{prop:interp-projection}에 의해 보간된 level-set은
참 중간 표면의 부호 거리 함수가 아니다. 표시 전용이며 임상적 복원이 아니다.

\section{범위의 한계}

\paragraph{단일 구조.}
좌심실 cavity만 구현되었다. 전체 좌/우심실 및 심근은 multiphase level-set을 요구하는
확장 범위이며, 실제 데이터 로더는 다른 구조를 요청하면 예외를 발생시킨다.

\paragraph{Digital twin이 아니다.}
조직 물성, 전기전도, 혈류, 역학 모델을 포함하지 않으며 치료 반응을 예측하지 않는다.
결과물은 시간가변 해부학 시각화 모델이다.

\paragraph{정답의 부재.}
ACDC와 M\&Ms-2는 ED/ES에만 label을 제공한다. 중간 프레임에는 분할 정답이 없으므로
그 프레임들은 외형과 시간적 안정성 지표로만 평가할 수 있다.

\section{현재 연구 상태의 한계}

\caveat{%
\textbf{구현이 실행되지 않았다.} 본 논문에서 가장 중대한 한계이다. PyTorch를 설치할 수
없는 환경에서 작성되었으므로 RQ1--RQ6의 어떤 주장도 확인되지도 반박되지도 않았다.
문법, import 해소, 이산 논리 커널 41개 항목은 검증되었으나(제
\ref{sec:kernel-verification}절) 그것은 파이프라인이 동작함을 뜻하지 않는다. 첫 실행에서
결함이 드러날 것으로 예상하는 것이 합리적이며, 그 목적으로 21스테이지 스모크 진단을
함께 제공하였다.}

\paragraph{통계가 구현되지 않았다.}
제 \ref{sec:statistics}절이 요구하는 환자 수준 paired 통계(bootstrap 신뢰구간,
Wilcoxon signed-rank)는 구현되지 않았다. 지표는 환자별로 산출되지만 집계는 실제
데이터와 함께 추가되어야 한다. 단일 팬텀 결과는 환자에 대한 증거로 읽을 수 없다.

\paragraph{비열등성 마진이 정해지지 않았다.}
결과를 본 뒤에 정하면 검정이 무의미해지므로 의도적으로 비워 두었다. pilot이 선행되어야
한다.

\paragraph{Dyna3DGR이 재현되지 않았다.}
가장 가까운 선행연구이지만 재현하지 않았고 어떤 수치도 귀속시키지 않았다. 구조적
비교만 제시하였다(표 \ref{tab:dyna3dgr}).

\paragraph{실제 데이터가 실행되지 않았다.}
ACDC/M\&Ms-2 로더는 구현되었으나 데이터에 접근할 수 없어 실행되지 않았다. 따라서
실제 해부학, scanner 변이, 질환에 대한 강건성은 전혀 평가되지 않았다.

\paragraph{PDF가 컴파일되지 않았다.}
본 원고를 작성한 환경에는 \LaTeX{} 배포판이 없었다. 따라서 이 문서 자체도 컴파일되어
확인된 적이 없으며, 조판 오류가 남아 있을 수 있다.

\section{다음에 해야 할 일}

우선순위 순서이다.

\begin{enumerate}[leftmargin=2em]
  \item PyTorch가 있는 환경에서 \texttt{python scripts/verify\_kernels.py}로 41/41을
  재확인한 뒤 \texttt{cvdyn2dgs smoke}로 파이프라인을 올린다.
  \item \texttt{cvdyn2dgs theory}로 제 \ref{ch:error}장의 수렴률을 측정한다. 이것이
  실패하면 이후 품질 수치는 무의미하다.
  \item 팬텀에서 RQ1--RQ6을 실행하고 \texttt{scripts/make\_tables.py}로 제
  \ref{ch:results}장을 채운다.
  \item ACDC 5명으로 pilot을 수행하고 비열등성 마진을 고정한다.
  \item ACDC 20--30명, M\&Ms-2 10--20명으로 본 실험을 수행하고 환자 수준 paired 통계를
  추가한다.
  \item 제 \ref{ch:discussion}장의 사전 해석 규칙에 따라 결론을 작성한다.
\end{enumerate}

\chapter{결론}
\label{ch:conclusion}

본 논문은 cine 심장 MRI에서 얻은 Chan--Vese 분할 표면 위에 법선방향 두께가 없는 2차원
Gaussian 서펠을 배치하여 심장 표면을 실시간으로 시각화하는 방법 CV-Dyn2DGS의 수학적
이론을 전개하고, 이를 PyTorch로 구현하였다.

\section{확립된 것}

\paragraph{이론.}
네 가지 결과를 증명하였다.

\begin{enumerate}[leftmargin=2em]
  \item \textbf{정리 \ref{thm:gradient-flow}}: Mumford--Shah 범함수의 이상 축약과 매끄러운
  Heaviside/Dirac 근사로부터 Chan--Vese 에너지 \eqref{eq:cv-energy}를 얻고, 변분법으로
  Euler--Lagrange 방정식 \eqref{eq:euler-lagrange}과 경사 흐름 \eqref{eq:gradient-flow}을
  부분적분 단계를 생략하지 않고 유도하였다. 여기에 복셀 간격을 명시적으로 반영한 간격
  인지 이산화 \eqref{eq:central-grad}, \eqref{eq:backward-div}, \eqref{eq:euler-update}를
  첫 번째 기여로 제시하였다.
  \item \textbf{정리 \ref{thm:warmstart}}: 경사 흐름의 에너지 소산(보조정리
  \ref{lem:dissipation})과 국소 PL 조건 아래의 선형 수렴(보조정리
  \ref{lem:linear-conv})이라는 기지 결과를 적용하여, 좁은 띠 웜스타트 초기값에 대한
  초기 오차와 반복 횟수의 관계 \eqref{eq:iteration-bound}을 유도하고 웜스타트가
  콜드스타트보다 적은 반복으로 수렴함 \eqref{eq:warmstart-gain}을 보였다. 전역
  비볼록성은 인정하고 논의를 국소 근방으로 제한하였다.
  \item \textbf{정리 \ref{thm:projection}}: 일차 Taylor 제약 \eqref{eq:linear-constraint}에
  대한 최소 노름 논법으로 최소 법선 보정 \eqref{eq:minimal-step}과 Newton 유사 반복
  \eqref{eq:projection-iter}을 직접 유도하고, 부호 거리 근방에서의 1--2회 수렴(명제
  \ref{prop:newton-steps})과 이차 잔차 한계(보조정리 \ref{lem:quadratic})를 Taylor
  나머지항 상쇄로 증명하였다.
  \item \textbf{네 축의 오차 분석}: 접평면 프레임(명제 \ref{prop:normal-error},
  보조정리 \ref{lem:transport-orthonormal}, 명제 \ref{prop:gauge}--\ref{prop:degeneracy}),
  원근 정합 래스터화(명제 \ref{prop:ray-plane}--\ref{prop:perspective}), 외형 잔차(명제
  \ref{prop:linearity}--\ref{prop:lowrank}), 중간 시점 보간(명제
  \ref{prop:interp-eikonal}--\ref{prop:interp-projection})에 대해 오차 한계를 유도하고,
  각 한계를 검증 지표에 대응시켰다(표 \ref{tab:error-metric}). 프레임 간 전파에서 위치
  잔차는 매 프레임 재설정되어 누적되지 않으나 방향 오차는 표류할 수 있음을 보였다(명제
  \ref{prop:propagation}).
\end{enumerate}

\paragraph{구현.}
49개 모듈로 위 이론을 구현하였다. 이 과정에서 이론이 명시하지 않은 자유도 세 가지를
특정 방식으로 고정하였고 그 선택을 기록하였다. 감독 신호를 광선 행진으로 정의한 것,
기하를 버퍼로 등록하여 최적화 불가능성을 구조적으로 보장한 것, 그리고 저장 표현이
임의 시점 탐색을 지원하는지의 문제를 명시한 것이다(제 \ref{sec:impl-decisions}절).

\paragraph{검증.}
PyTorch를 필요로 하지 않는 이산 논리 커널 41개 항목을 독립 참조와 대조하여 검증하였다
(제 \ref{sec:kernel-results}절). 이 과정에서 검증 자체의 결함을 찾은 사례가 있었고,
오류 개수가 포함--배제 예측값과 정확히 일치한다는 사실이 대상 알고리즘의 정확성에 대한
증거가 되었다.

\section{확립되지 않은 것}

\caveat{%
본 연구의 \textbf{실험적 주장은 하나도 확인되지 않았다}. 구현은 완성되었으나 실행되지
않았다. 따라서 다음은 모두 열려 있다.
\begin{itemize}[leftmargin=1.5em,topsep=2pt]
  \item 식 \eqref{eq:cost-hypothesis}의 비용 부등식이 성립하는가 (RQ1, RQ2),
  \item $\CR>1$인가, 즉 제안 표현이 애초에 더 작은가 (RQ3),
  \item 95번째 백분위 프레임 latency가 33.3\,ms 미만인가 (RQ4),
  \item 2DGS가 동일 표면의 mesh보다 어느 지표에서든 나은가 (RQ5),
  \item 2D 원반이 얇은 3D Gaussian보다 나은가 (RQ6).
\end{itemize}
제 \ref{ch:discussion}장은 이 중 어느 결과가 나오든 어떤 결론을 내릴지를 사전에
확정해 두었다. 특히 RQ5가 부정적이면 ``이 목적에는 mesh가 더 적절하다''가 올바른
결론이며, 그 경우에도 위 이론적 기여는 유지된다.}

\section{맺음말}

본 연구의 수학적 중심은 식 \eqref{eq:cv-energy}의 변분 분할, 식 \eqref{eq:minimal-step}의
최소법선 이동 유도, 식 \eqref{eq:surfel-point}의 2차원 표면 매개변수화, 식
\eqref{eq:transport-project}--\eqref{eq:transport-normalize}의 접평면 갱신, 그리고 식
\eqref{eq:residual-ls}의 정칙화된 잔차 추정이다. Viewer는 이러한 모델의 효율과 활용
가능성을 보여주는 응용 결과이지 그 자체가 기여는 아니다.

본 논문이 기여하는 바를 한 문장으로 요약하면 다음이다. \emph{분할이 이미 계산한 미분
정보를 렌더링 primitive의 시간 갱신에 직접 사용할 수 있으며, 그 갱신의 오차를 네 축으로
분해하여 각각 한계지을 수 있다.} 이 갱신이 실제로 프레임별 재최적화보다 싸고 충분히
정확한지는 실험이 답할 문제이며, 그 실험을 수행할 수 있는 시스템과 그것을 정직하게
보고할 수 있는 구조를 함께 제시하였다.


\appendix
\chapter{기호 정리}
\label{app:symbols}

\section{분할과 level-set}

\begin{longtable}{@{}lp{10.2cm}@{}}
\toprule
기호 & 의미 \\
\midrule
\endhead
$\dom \subset \Real^3$ & 영상이 정의된 유계 열린 영역. $\partial\dom$은 조각마다 매끄럽다. \\
$\imgt : \dom \to \Real$ & 프레임 $t$의 3차원 영상. 유계이며 제곱적분가능. \\
$\levt : \dom \to \Real$ & 프레임 $t$의 level-set 함수. Lipschitz 이상으로 매끄럽다. \\
$\surft = \{\levt = 0\}$ & 프레임 $t$의 표면. \textbf{$\levt>0$이 내부}(식 \eqref{eq:sign}). \\
$\spacing = (h_x,h_y,h_z)$ & 복셀 간격 (mm). MRI에서 보통 $h_z \ge h_x \approx h_y$. \\
$\hmax$ & $\max\{h_x,h_y,h_z\}$. 명제 \ref{prop:normal-error}의 수렴률 변수. \\
$\gradh, \divh$ & 간격 인지 경사와 발산 (식 \eqref{eq:central-grad}, \eqref{eq:backward-div}). \\
$D^\pm_a$ & 축 $a$의 전진/후진 차분 (식 \eqref{eq:forward-diff}, \eqref{eq:backward-diff}). \\
$\cin, \cout$ & 프레임 $t$의 내부/외부 영역 평균 (식 \eqref{eq:cin}, \eqref{eq:cout}). \\
$\mu > 0$ & 길이/정규화 항의 가중치. \\
$\lin, \lout > 0$ & 내부/외부 자료 충실도 항의 가중치. \\
$\varepsilon > 0$ & Heaviside/Dirac 근사 폭 및 분모 정규화 상수. \\
$\Heps, \deps$ & 매끄러운 Heaviside와 Dirac (식 \eqref{eq:heaviside}, \eqref{eq:dirac}). \\
$b > 0$ & 웜스타트 좁은 띠의 폭 (식 \eqref{eq:warmstart-init}). \\
$s \ge 0$ & \textbf{인위적} 최적화 시간. 물리적 심장 시간이 아니다. \\
$t$ & cine 영상의 프레임 지표, 즉 실제 시점. \\
$\kappa$ & 평균 곡률 또는 주곡률 크기의 최댓값 (문맥에 따라). \\
$T$ & 프레임 수. \\
$\eta$ & 인접 프레임 해의 근접성 상수 (가정 \ref{ass:small-change}). \\
$m, L$ & 국소 PL 상수와 Lipschitz 상수 (가정 \ref{ass:pl}). \\
$\rho = 1-\alpha m$ & 선형 수렴 비율 (보조정리 \ref{lem:linear-conv}). \\
\bottomrule
\end{longtable}

\section{서펠 기하}

서펠은 총 $N$개이며 지표 $i=1,\dots,N$로 센다.

\begin{longtable}{@{}lp{10.2cm}@{}}
\toprule
기호 & 의미 \\
\midrule
\endhead
$\anchor^t_i \in \dom$ & 서펠 $i$의 프레임 $t$에서의 anchor(원반 중심). \\
$e^t_{i,1}, e^t_{i,2} \in \Real^3$ & 두 접선 축. 서로 직교하는 단위 벡터. \\
$s_{i,1}, s_{i,2} > 0$ & 두 접선 축 방향의 scale. \textbf{법선방향 scale은 없다.} \\
$\amp^t_i \in \Real$ & 외형 진폭. $\amp^0_i + \Delta\amp^t_i$ (식 \eqref{eq:amplitude-split}). \\
$\opa^t_i \in [0,1]$ & 불투명도. \\
$\nrm^t_i \in \Real^3$ & 단위 법선 (식 \eqref{eq:unit-normal}). \\
$\framemat^t_i \in \Real^{3\times2}$ & 접평면 프레임 행렬 (식 \eqref{eq:frame}). \\
$\scalemat_i \in \Real^{2\times2}$ & $\operatorname{diag}(s_{i,1},s_{i,2})$. \\
$u \in \Real^2$ & 서펠 국소 좌표 (식 \eqref{eq:surfel-point}, \eqref{eq:local-coord}). \\
$\Gcan$ & 정준 서펠 집합 (식 \eqref{eq:canonical-set}). \\
$\store$ & 저장 표현 (식 \eqref{eq:store}). \\
$K_p$ & 정사영 반복 횟수 (식 \eqref{eq:projection-iter}). \\
$\psi$ & 이전 접선 축과 새 법선 사이의 각 (명제 \ref{prop:degeneracy}). \\
$\delta$ & 접평면 프레임의 면내 정렬 오차각 (명제 \ref{prop:aniso}). \\
\bottomrule
\end{longtable}

\section{렌더링과 지표}

\begin{longtable}{@{}lp{10.2cm}@{}}
\toprule
기호 & 의미 \\
\midrule
\endhead
$\camc \in \Real^3$ & 카메라 중심. \\
$\raydir \in \Real^3$ & 화소 $q$ 방향의 \textbf{단위} 광선 벡터. \\
$\depthq^t_i(q)$ & 광선--서펠 교차 깊이 (mm) (식 \eqref{eq:ray-plane}). \\
$\alphaq^t_i(q)$ & 화소 불투명도 (식 \eqref{eq:alpha}). \\
$w_i(q)$ & 투과율 가중치 (식 \eqref{eq:weights}). \\
$\Chat_t(q)$ & 합성된 화소 값 (식 \eqref{eq:compositing}). \\
$A_t$ & 렌더링 행렬. 진폭에 대해 선형 (명제 \ref{prop:linearity}). \\
$W$ & 화소 가중 행렬 (보통 심장 ROI). \\
$\lambda_a, \lambda_T$ & 잔차의 Tikhonov 및 시간 정칙화 계수. \\
$r$ & 저계수 잔차의 계수 (식 \eqref{eq:lowrank}). \\
$c_0$ & 스침각 컷오프 (명제 \ref{prop:ray-plane}). \\
$\Esurf$ & 표면 정렬 오차 (식 \eqref{eq:e-surf}). \\
$\Eflicker$ & 시간적 흔들림 (식 \eqref{eq:flicker}). \\
$\Sfull, \Sours, \CR$ & 저장량과 압축비 (식 \eqref{eq:s-full}--\eqref{eq:cr}). \\
$P_{2D}, P_r$ & 서펠당, 잔차당 parameter 수. \\
$\beta \in [0,1]$ & 중간 시점 보간 계수 (식 \eqref{eq:interpolation}). \\
$\omega$ & 인접 프레임 경사 사이의 각 (명제 \ref{prop:interp-eikonal}). \\
\bottomrule
\end{longtable}

\chapter{수식과 구현의 대응}
\label{app:code-map}

본 부록은 원고의 각 수식 및 정리를 구현의 어느 위치가 실현하는지 명시한다. 이 대응이
없으면 원고와 구현이 시간이 지나며 조용히 어긋난다.

\section{Level-set과 Chan--Vese}

\begin{longtable}{@{}p{3.6cm}p{4.6cm}p{5.2cm}@{}}
\toprule
수식/결과 & 모듈 & 함수 \\
\midrule
\endhead
식 \eqref{eq:heaviside}, \eqref{eq:dirac} & \texttt{levelset/operators.py}
  & \texttt{heaviside\_eps}, \texttt{dirac\_eps} \\
식 \eqref{eq:forward-diff}--\eqref{eq:backward-div} & \texttt{levelset/operators.py}
  & \texttt{forward\_diff}, \texttt{backward\_diff}, \texttt{central\_diff},
    \texttt{gradient\_central}, \texttt{divergence\_backward}, \texttt{curvature} \\
식 \eqref{eq:cv-energy} & \texttt{levelset/operators.py} & \texttt{chanvese\_energy} \\
식 \eqref{eq:cin}, \eqref{eq:cout} & \texttt{levelset/operators.py} & \texttt{region\_means} \\
식 \eqref{eq:gradient-flow} & \texttt{levelset/operators.py} & \texttt{chanvese\_speed} \\
식 \eqref{eq:euler-update}, \eqref{eq:cfl} & \texttt{levelset/chanvese.py} & \texttt{solve\_frame} \\
식 \eqref{eq:warmstart-init} & \texttt{levelset/chanvese.py}
  & \texttt{solve\_sequence} (\texttt{warm\_start}), \texttt{\_bbox\_from\_mask} \\
정리 \ref{thm:warmstart} 의 $K(e_0)$ & \texttt{levelset/chanvese.py}
  & \texttt{solve\_frame(reference\_phi=..., tol\_to\_reference=...)} \\
명제 \ref{prop:newton-steps}의 재초기화 & \texttt{levelset/sdf.py} & \texttt{reinitialize} \\
식 \eqref{eq:interpolation} & \texttt{levelset/sdf.py} & \texttt{interpolate\_levelsets} \\
표면 추출 & \texttt{levelset/mesh\_extract.py} & \texttt{marching\_tetrahedra} \\
\bottomrule
\end{longtable}

\section{서펠 기하}

\begin{longtable}{@{}p{3.6cm}p{4.6cm}p{5.2cm}@{}}
\toprule
수식/결과 & 모듈 & 함수 \\
\midrule
\endhead
식 \eqref{eq:canonical-set} & \texttt{surfel/model.py} & \texttt{SurfelSet2D} \\
식 \eqref{eq:store} & \texttt{pipeline/precompute.py} & \texttt{PrecomputedModel} \\
식 \eqref{eq:scale-init}, \eqref{eq:curvature-scale} & \texttt{surfel/canonical.py}
  & \texttt{initialize\_canonical\_surfels} \\
식 \eqref{eq:unit-normal} & \texttt{surfel/canonical.py} & \texttt{surface\_normals\_at} \\
식 \eqref{eq:minimal-step}, \eqref{eq:projection-iter} & \texttt{surfel/projection.py}
  & \texttt{project\_to\_surface} \\
식 \eqref{eq:e-surf} & \texttt{surfel/projection.py}, \texttt{metrics/segmentation.py}
  & \texttt{e\_surf} \\
보조정리 \ref{lem:quadratic}의 잔차 이력 & \texttt{surfel/projection.py}
  & \texttt{ProjectionResult.residual\_history} \\
식 \eqref{eq:transport-project}, \eqref{eq:transport-normalize} & \texttt{surfel/transport.py}
  & \texttt{transport\_tangent\_frame} \\
명제 \ref{prop:degeneracy}의 퇴화 처리 & \texttt{surfel/transport.py}
  & \texttt{degeneracy\_thresh}, \texttt{TransportDiagnostics} \\
밀도 제어 (제 \ref{sec:transport}절) & \texttt{surfel/density.py}
  & \texttt{tangential\_repulsion}, \texttt{densify}, \texttt{prune} \\
\bottomrule
\end{longtable}

\section{렌더링}

\begin{longtable}{@{}p{3.6cm}p{4.6cm}p{5.2cm}@{}}
\toprule
수식/결과 & 모듈 & 함수 \\
\midrule
\endhead
식 \eqref{eq:ray}--\eqref{eq:compositing} & \texttt{render/raster2dgs.py} & \texttt{render\_2dgs} \\
제 \ref{sec:impl-decisions}절의 참조 구현 & \texttt{render/raster2dgs.py}
  & \texttt{render\_2dgs\_reference} \\
식 \eqref{eq:weights}, 명제 \ref{prop:linearity} & \texttt{render/raster2dgs.py}
  & \texttt{render\_weights}, \texttt{WeightMatrix} \\
명제 \ref{prop:ray-plane}의 컷오프 & \texttt{render/raster2dgs.py}
  & \texttt{RenderConfig.cull\_cos} \\
식 \eqref{eq:distortion-surrogate} & \texttt{render/raster2dgs.py} & \texttt{\_distortion} \\
명제 \ref{prop:perspective}의 아핀 비교 & \texttt{render/raster3dgs.py} & \texttt{render\_thin\_3dgs} \\
mesh 기준선 (제 \ref{sec:baselines}절) & \texttt{render/mesh\_render.py} & \texttt{render\_mesh} \\
감독 신호 정의 (제 \ref{sec:impl-decisions}절) & \texttt{render/raymarch.py}
  & \texttt{raymarch\_levelset} \\
카메라 모델 & \texttt{render/camera.py}
  & \texttt{Camera.orbit}, \texttt{Camera.from\_slice\_plane} \\
식 \eqref{eq:total-loss} & \texttt{losses.py} & \texttt{total\_loss} \\
\bottomrule
\end{longtable}

\section{잔차와 저장}

\begin{longtable}{@{}p{3.6cm}p{4.6cm}p{5.2cm}@{}}
\toprule
수식/결과 & 모듈 & 함수 \\
\midrule
\endhead
식 \eqref{eq:residual-ls}, \eqref{eq:normal-equations} & \texttt{residual/solver.py}
  & \texttt{solve\_residual}, \texttt{conjugate\_gradient} \\
식 \eqref{eq:cond-bound} & \texttt{residual/solver.py}
  & \texttt{\_estimate\_sigma\_max}, \texttt{ResidualResult.condition\_bound} \\
식 \eqref{eq:lowrank}, 명제 \ref{prop:lowrank} & \texttt{residual/lowrank.py}
  & \texttt{compress\_residual}, \texttt{rank\_for\_error} \\
식 \eqref{eq:s-full}--\eqref{eq:cr} & \texttt{metrics/storage.py} & \texttt{storage\_report} \\
식 \eqref{eq:breakeven} & \texttt{metrics/storage.py}
  & \texttt{StorageReport.break\_even\_surface\_bytes\_per\_frame} \\
직렬화 & \texttt{pipeline/storage\_io.py}
  & \texttt{save\_model}, \texttt{pack\_narrow\_band} \\
\bottomrule
\end{longtable}

\section{파이프라인, 지표, 실험}

\begin{longtable}{@{}p{3.6cm}p{4.6cm}p{5.2cm}@{}}
\toprule
수식/결과 & 모듈 & 함수 \\
\midrule
\endhead
제 \ref{sec:pipeline}절의 7단계 & \texttt{pipeline/precompute.py} & \texttt{precompute} \\
식 \eqref{eq:cost-ours}의 네 항 & \texttt{pipeline/precompute.py}
  & \texttt{FrameRecord.per\_frame\_update\_ms} \\
식 \eqref{eq:frame-budget} & \texttt{pipeline/playback.py}
  & \texttt{PlaybackEngine.step\_to}, \texttt{FrameTiming} \\
제 \ref{sec:seek-problem}절의 탐색 문제 & \texttt{pipeline/playback.py}
  & \texttt{compare\_projection\_modes} \\
식 \eqref{eq:flicker} & \texttt{metrics/photometric.py} & \texttt{flicker}, \texttt{temporal\_report} \\
식 \eqref{eq:ef} & \texttt{metrics/clinical.py} & \texttt{volume\_curve\_report} \\
팬텀 (제 \ref{sec:phantom}절) & \texttt{data/phantom.py}
  & \texttt{make\_phantom}, \texttt{ellipsoid\_sdf} \\
RQ1--RQ6 & \texttt{experiments/research\_questions.py} & \texttt{rq1\_...}--\texttt{rq6\_...} \\
제 \ref{sec:ablation-design}절 & \texttt{experiments/ablations.py}
  & \texttt{ABLATIONS}, \texttt{spacing\_ablation} \\
표 \ref{tab:error-metric}의 수치 확인 & \texttt{experiments/theory\_checks.py}
  & \texttt{run\_all\_checks} \\
제 \ref{sec:kernel-verification}절 & \texttt{scripts/verify\_kernels.py} & --- \\
제 \ref{sec:impl-status}절의 스모크 & \texttt{cvdyn2dgs/smoke.py} & \texttt{run\_smoke} \\
제 \ref{ch:results}장의 표 생성 & \texttt{scripts/make\_tables.py} & --- \\
\bottomrule
\end{longtable}


\bibliographystyle{plainnat}
\bibliography{refs}

\end{document}
